% comprehension_exp.tex
% Stand-alone plan for the comprehension experiment.
% Can be embedded in main.tex via % comprehension_exp.tex
% Stand-alone plan for the comprehension experiment.
% Can be embedded in main.tex via % comprehension_exp.tex
% Stand-alone plan for the comprehension experiment.
% Can be embedded in main.tex via % comprehension_exp.tex
% Stand-alone plan for the comprehension experiment.
% Can be embedded in main.tex via \input{comprehension_exp}.
%
% Compile stand-alone:
%   pdflatex comprehension_exp.tex
%
% Embed in main.tex:
%   \input{comprehension_exp}   % (inside a \section{Experiment} environment)

\documentclass[12pt]{article}
\usepackage{CSPArticleStyle}
\usepackage{csquotes}
\usepackage[margin=1in]{geometry}
\usepackage{booktabs}
\usepackage{parskip}
\usepackage{amsmath}
\usepackage{times}
\usepackage[T1]{fontenc}
\usepackage[utf8]{inputenc}
\usepackage{xcolor}
\usepackage{tabularray}
\usepackage{gb4e}
\noautomath

\title{Comprehension Experiment: Inferring Speaker Goal and Action Adoption
       from Consensus-Marker Use}
\author{}
\date{}

\begin{document}

%% ============================================================
%% When embedding in main.tex, delete from \documentclass to
%% \begin{document} and delete \end{document} at the bottom.
%% ============================================================

\section{Experiment: Consensus-Marker Comprehension}

\subsection{Overview}

The comprehension experiment is the listener-side counterpart of the
production experiment and tests whether listeners can invert the
speaker's marker choice to recover the latent context that motivated it.
Where the production experiment gave participants a context (fixing
\textsc{pc\textsubscript{prag}} and $g$) and asked which marker they
would use, the comprehension experiment gives participants a context
(fixing \textsc{pc\textsubscript{prag}}) together with a marker already
embedded in the sentence and asks what goal strength the speaker likely
had and how likely the listener is to adopt $q$.

The specific marker presented in each vignette is not chosen arbitrarily
from the three candidate markers. Our preregistered plan was to assign,
for each item$\times$condition cell, the \emph{Bayesian best-choice marker}:
the marker with the highest posterior probability under the cost-augmented
noisy-threshold RSA model fitted to the production data, given the context's
\textsc{pc\textsubscript{prag}} level and a specified $g$ level. This was
intended to mirror the production design closely and to ensure that the
comprehension stimuli reflect model-based marker use rather than arbitrary
marker$\times$item pairings.

Applying this strict modal-assignment rule to the fitted production model,
however, yielded a highly imbalanced stimulus set: most critical cells were
assigned the strongest marker, with very little variation across marker
categories. In addition, many of these modal decisions were supported only by
very small probability margins over the next-best alternative. Because such a
stimulus set would make the comprehension experiment substantially less
informative about marker-based listener inferences, we introduced a documented,
exploratory balancing step for stimulus construction. Starting from the same
posterior predictive probabilities, we selected a near-balanced assignment of
the three markers across the 32 critical cells while maximizing total posterior
support under the fitted production model. The strict modal assignment is
retained as the preregistered baseline, but the balanced assignment is used for
the actual comprehension materials.

\textsc{pc\textsubscript{prag}} is retained as an explicit context
manipulation (low vs.\ high), so that listeners can be expected to
have a well-defined prior on pragmatic controversy when interpreting
the marker. \textsc{pc\textsubscript{prop}} is not manipulated within
the experiment; it enters as a continuous topic-level covariate from
the norming study.

After reading each vignette, participants provide two slider ratings
operationalising the key listener-side quantities:
\begin{enumerate}
  \item \textbf{Inferred speaker goal strength} ($\widehat{g}$):
        how strongly does the speaker seem to want the listener to
        accept and act on $q$?
  \item \textbf{Adoption likelihood} ($\widehat{A}$):
        how likely is the listener to follow the recommended action $q$?
\end{enumerate}

\noindent Perceived controversy is \emph{not} elicited as a rating:
because \textsc{pc\textsubscript{prag}} is already manipulated as a
context factor, asking participants to rate controversy would be
redundant with respect to the manipulation check and circular with
respect to the research question. The theoretically central inference
to test is whether listeners can decode the speaker's \emph{goal} from
marker use, and whether that inference translates into adoption
behaviour.

The two slider measures, together with the \textsc{pc\textsubscript{prag}}
context and the \textsc{pc\textsubscript{prop}} covariate, provide the
data needed to fit the model's listener parameters
$(\eta_0, \eta_g, \eta_{pc}, \eta_{\mathrm{int}})$.

\subsection{Participants}

Participants will be recruited via Prolific and will be required to be
native speakers of German (language-at-home filter).
Data collection is planned after completion of the production experiment
(model fitting required to determine the best-choice markers).
A target of $N = 80$ participants is planned ($n = 20$ per list across
4 counterbalancing lists), each completing 16 trials (8 critical + 8
filler), matching the structure of the production experiment.

\subsection{Materials}

\subsubsection{Markers}

Each critical vignette is presented with one specific consensus marker
already embedded in the Mittelfeld of the sentence frame (see
example~(\ref{ex:frame_comp})):

\begin{exe}
  \ex \label{ex:frame_comp}
  Digitale Kompetenzen sind \textbf{[MARKER]}
  in der heutigen Arbeitswelt unverzichtbar.\\
  `Digital skills are \textbf{[MARKER]} indispensable in today's
  working world.'
\end{exe}

Unlike the production experiment, participants do \emph{not} choose the
marker; they read the sentence with the marker already in place and respond
to the two rating probes.

Critically, the marker assigned to each condition is not drawn arbitrarily
from the three candidates. Under the preregistered plan, for each of the four
production conditions (pc\textsubscript{prag} $\times$ $g$), the marker used in
comprehension was to be the \textbf{Bayesian best-choice marker}: the marker
$u^*$ with the highest posterior probability under the fitted
cost-augmented noisy-threshold RSA production model,
\[
u^* = \arg\max_{u} \; P(u \mid pc_{\mathrm{prag}}, g,
                           pc_{\mathrm{prop}\_c}).
\]
Because $pc_{\mathrm{prop}}$ varies continuously across items,
$u^*$ may differ per item even within the same
(pc\textsubscript{prag}, $g$) cell; the final stimulus list is therefore
constructed item-by-item once the production model has been fitted.
In practice, the fitted production model yielded a nearly degenerate modal map,
with most cells assigned to \textit{bekanntlich}. We therefore retained the
model-based logic of the procedure but replaced strict cellwise modal selection
with an uncertainty-aware balancing step: among the marker assignments with
high posterior support, we selected a near-balanced configuration across the 32
critical cells. This preserves the production--comprehension link while also
ensuring that the comprehension experiment contains sufficient marker variation
to test listener-side inferences.

All three markers remain candidates; the mapping from conditions to
markers is an empirical output of the production analysis, not a design
pre-commitment. In the final comprehension materials, we use the
\textbf{balanced model-based assignment}, rather than the strict modal
assignment, in order to maintain informative variation across marker types.

\subsubsection{Topic Items}

The same eight topic domains used in the norming study and the production
experiment are used here (Table~\ref{tab:items_comp}).
The 4-condition $\times$ 8-item Latin square (Section~\ref{subsub:counterbalancing_comp})
requires no additional items.

\begin{table}[ht]
\centering
\caption{Topic items: sentence frame for $p$ (with marker position [M])
         and consequent $q$.}
\label{tab:items_comp}
\small
\begin{tabular}{clll}
\toprule
\textbf{\#} & \textbf{Topic} & \textbf{Sentence frame} & \textbf{Consequent $q$} \\
\midrule
1 & Klimawandel &
  \textit{Der Klimawandel wird [M]} &
  auf Flugreisen verzichten \\
  & & \textit{durch menschliche Aktivit\"aten verursacht.} & \\[3pt]
2 & Ern\"ahrung &
  \textit{Eine \"uberwiegend pflanzliche Ern\"ahrung} &
  mehr pflanzliche Lebens- \\
  & & \textit{f\"ordert [M] die Gesundheit.} &
  mittel einbauen \\[3pt]
3 & Stadtverkehr &
  \textit{Ein ausgebautes Nahverkehrsnetz} &
  h\"aufiger auf \\
  & & \textit{entlastet [M] den Stadtverkehr} &
  \"offentliche Verkehrs- \\
  & & \textit{deutlich.} &
  mittel umsteigen \\[3pt]
4 & Digitalisierung &
  \textit{Digitale Kompetenzen sind [M] in} &
  regelm\"a\ss{}ig Weiter- \\
  & & \textit{der heutigen Arbeitswelt unverzichtbar.} &
  bildungsangebote nutzen \\[3pt]
5 & Schlaf &
  \textit{Ausreichend Schlaf ist [M] entscheidend} &
  auf feste Schlafzeiten \\
  & & \textit{f\"ur die kognitive Leistungsf\"ahigkeit.} &
  achten \\[3pt]
6 & Plastik &
  \textit{Einwegplastik schadet [M] den} &
  konsequent auf Einweg- \\
  & & \textit{Meeres\"okosystemen erheblich.} &
  plastik verzichten \\[3pt]
7 & Sport &
  \textit{Regelm\"a\ss{}ige k\"orperliche Bewegung} &
  regelm\"a\ss{}ig Sport in \\
  & & \textit{senkt [M] das Risiko f\"ur} &
  den Alltag integrieren \\
  & & \textit{Herz-Kreislauf-Erkrankungen deutlich.} & \\[3pt]
8 & Lokalkauf &
  \textit{Das Kaufen bei lokalen H\"andlern} &
  \"ofter bei lokalen \\
  & & \textit{st\"arkt [M] die regionale Wirtschaft} &
  Gesch\"aften einkaufen \\
  & & \textit{nachhaltig.} & \\
\bottomrule
\end{tabular}
\end{table}

\noindent [M] denotes the position of the pre-filled best-choice marker.

\subsubsection{Vignettes}

\paragraph{Structure.}
Each trial presents a vignette consisting of three parts:
(a) a context paragraph establishing the communicative situation and
fixing the \textsc{pc\textsubscript{prag}} level;
(b) the target sentence with the best-choice marker already embedded
in the Mittelfeld;
(c) the consequent action clause.
The two rating probes follow immediately below the vignette.

\paragraph{Propositional perceived controversy (\textsc{pc\textsubscript{prop}}).}
As in the production experiment, \textsc{pc\textsubscript{prop}} is not
manipulated within the comprehension experiment.
Topic-level propositional acceptance ratings from the norming study
(0--100 scale; e.g., \textit{Digitalisierung}: 85/100,
\textit{Lokalkauf}: 58/100) serve as a continuous, mean-centred
covariate \textsc{pc\textsubscript{prop}}$_c$ in analysis.
Context paragraphs present $p$ without editorial evaluation; variation in
listeners' individual assessments of propositional controversy is
captured by the norming-derived covariate.

\paragraph{Pragmatic perceived controversy (\textsc{pc\textsubscript{prag}}).}
\textsc{pc\textsubscript{prag}} is manipulated between item instances
using the same two context versions as in the production experiment:
the speaker's social circle either \emph{largely agrees} (low) or
\emph{holds divergent views} (high) about $p$.
For example, for the \textit{Digitalisierung} item:

\begin{itemize}
  \item $\textsc{pc}_\text{prag} = \text{low}$: \textit{Deine Kollegin spricht
        mit dir \"uber Berufsausbildung und Weiterbildung.
        In eurem Team teilen die meisten die Ansicht, dass digitale Kompetenzen
        in der heutigen Arbeitswelt unverzichtbar sind.}
  \item $\textsc{pc}_\text{prag} = \text{high}$: \textit{Deine Kollegin spricht
        mit dir \"uber Berufsausbildung und Weiterbildung.
        In eurem Team gehen die Meinungen dazu stark auseinander.}
\end{itemize}

The context is framed from the listener's perspective as a direct addressee,
mirroring the production study where participants adopt the speaker's
perspective.

\paragraph{Marker manipulation.}
The dependent measures of the production experiment (which marker is chosen)
become the independent variable here.
One of the three consensus markers is inserted into the Mittelfeld of the
target sentence.
A complete sample item for the \textit{Digitalisierung} topic
($\textsc{pc}_\text{prag} = \text{high}$, marker = \textit{bekanntlich};
norming $\textsc{pc}_\text{prop}$ rating: 85/100) is given in
(\ref{ex:comp_sample}):

\begin{exe}
  \ex \label{ex:comp_sample}
  \textit{Deine Kollegin spricht mit dir \"uber Berufsausbildung und
  Weiterbildung.
  In eurem Team gehen die Meinungen dazu stark auseinander.}

  \textit{Sie sagt zu dir:} Digitale Kompetenzen sind \textbf{bekanntlich}
  in der heutigen Arbeitswelt unverzichtbar.
  Deshalb solltest du regelm\"a\ss{}ig digitale Weiterbildungsangebote nutzen.

  \noindent\textit{[Rating probes follow.]}
\end{exe}

\subsection{Design}

\subsubsection{Factorial structure}

The comprehension experiment mirrors the production experiment's
$\mathbf{2}$~\textbf{(pc\textsubscript{prag}: low vs.\ high)}
$\mathbf{\times}$~$\mathbf{2}$~\textbf{($g_{\mathrm{implied}}$: low
vs.\ high) design}, yielding \textbf{4 conditions}.
The $g_{\mathrm{implied}}$ dimension is not manipulated via an explicit
goal instruction (as in production) but is instead conveyed entirely
through the choice of marker: for each condition, the marker presented
is the best-choice marker for the corresponding
(\textsc{pc\textsubscript{prag}},~$g$) cell from the production model.
\textsc{pc\textsubscript{prop}} enters as a continuous norming covariate.
With 8 topic items and a Latin square rotation across 4 lists, each
participant sees exactly \textbf{8 critical trials}---one item per
condition, each condition appearing twice---plus 8 filler trials,
for \textbf{16 trials} in total (Table~\ref{tab:design_comp}).

\begin{table}[ht]
\centering
\caption{%
  The $2\times2$ design (4 conditions), mirroring the production experiment.
  The marker presented in each condition is the Bayesian best-choice
  marker from the fitted production model; specific markers are
  determined after production data collection.}
\label{tab:design_comp}
\small
\begin{tabular}{ccll}
\toprule
\textbf{pc\textsubscript{prag}} &
\textbf{$g_{\mathrm{implied}}$} &
\textbf{Marker (TBD from production)} &
\textbf{Condition label} \\
\midrule
low  & low  & best-choice for $(\text{low}, g_{\text{low}})$  &
  (0) community agrees, weak signal \\
low  & high & best-choice for $(\text{low}, g_{\text{high}})$ &
  (1) community agrees, strong signal \\
high & low  & best-choice for $(\text{high}, g_{\text{low}})$ &
  (2) community divided, weak signal \\
high & high & best-choice for $(\text{high}, g_{\text{high}})$ &
  (3) community divided, strong signal \\
\bottomrule
\end{tabular}
\end{table}

\subsubsection{Counterbalancing}
\label{subsub:counterbalancing_comp}

The experiment uses the \textbf{same $4 \times 8$ Latin square} as the
production experiment: each of the 4 lists assigns item $i$
($i = 0,\ldots,7$) to condition $(i + l)\bmod 4$, where
$l \in \{0,1,2,3\}$ is the list number.
This guarantees that:
\begin{itemize}
  \item each participant sees every item exactly once;
  \item each participant sees each of the 4 conditions exactly twice
        (two items per condition), ensuring balanced coverage of the
        $2\times2$ design;
  \item across all 4 lists, each item appears equally often in each
        condition.
\end{itemize}

List assignment follows the same Prolific-and-URL strategy as the production
experiment: each Prolific study embeds \texttt{?list=N} (0--3) in the
experiment URL, capping recruitment at $n = 20$ participants per list
for a total of $N = 80$.

In addition to the 8 critical trials, each participant completed 8
\textbf{filler trials} using the same filler items as the production
experiment (train travel, reading light, cooking, coffee, vacation,
hydration, fresh air, noise), each presented with a randomly assigned
marker.
Fillers are excluded from analysis; they serve to dilute the critical
items and reduce demand characteristics.

\subsubsection{Dependent variables}
\label{subsub:dvs_comp}

Two slider ratings are collected after each vignette (0--100 continuous
scale with no visible numbers, labelled at the endpoints as shown):

\begin{enumerate}
  \item \textbf{Inferred speaker goal strength} (operationalises $\widehat{g}$):
  \begin{quote}
    \textit{Wie stark hast du den Eindruck, dass die Sprecherin m\"ochte,
    dass du ihre Empfehlung befolgst?}\\
    (`How strongly do you get the impression that the speaker wants you to
    follow their recommendation?')\\
    Endpoints: 0 = \textit{\"uberhaupt nicht}
    (`not at all') ---
    100 = \textit{sehr stark} (`very strongly').
  \end{quote}

  \item \textbf{Adoption likelihood of $q$} (operationalises $\widehat{A}$):
  \begin{quote}
    \textit{Wie wahrscheinlich ist es, dass du die empfohlene Handlung
    dann auch tats\"achlich umsetzt?}\\
    (`How likely is it that you will actually carry out the recommended
    action?')\\
    Endpoints: 0 = \textit{sehr unwahrscheinlich}
    (`very unlikely') ---
    100 = \textit{sehr wahrscheinlich} (`very likely').
  \end{quote}
\end{enumerate}

\noindent Perceived controversy is deliberately \emph{not} included as a
rating probe: \textsc{pc\textsubscript{prag}} is experimentally
manipulated via the context paragraph, so a controversy rating would
be redundant and risks obscuring the marker's contribution to inferred
goal strength.

Recorded per trial: \texttt{trial\_id}, \texttt{topic},
\texttt{is\_filler}, \texttt{condition\_index}, \texttt{marker},
\texttt{pc\_prag}, \texttt{g\_implied}, \texttt{rating\_g},
\texttt{rating\_adopt}, \texttt{rt} (ms); and once per participant:
\texttt{list\_num}.
The continuous \texttt{pc\_prop\_rating} covariate (topic-level norming
mean) is joined to the trial data during analysis.

\subsection{Procedure}

Participants complete the experiment online via the magpie platform
\citep{kochetkova2024magpie} in a browser.
After reading instructions in German, participants are assigned to one of
4 counterbalancing lists (via URL \texttt{?list=N} parameter assigned by
Prolific; no random fallback in deployment).
They then proceed through 16 trials---8 critical and 8 fillers---in a
fully randomised order, followed by a post-test questionnaire eliciting
native language, age, gender, education level, and optional comments.

Each trial displays:
(a) a context paragraph establishing the communicative situation and the
\textsc{pc\textsubscript{prag}} level;
(b) the target sentence with the best-choice marker already in place in
the Mittelfeld;
(c) the consequent action clause;
(d) the two rating sliders in fixed order ($\widehat{g}$, $\widehat{A}$),
each labelled in German.
The ``Weiter'' (Next) button is disabled until both sliders have been
moved at least once, ensuring a complete response on every trial.

\subsection{Planned Analysis}

\subsubsection{Empirical plots}

For each dependent variable ($\widehat{g}$, $\widehat{A}$):

\begin{itemize}
  \item \textbf{Condition means}: a dot-and-whisker plot shows mean
        ratings per condition (the $2\times2$ of pc\textsubscript{prag}
        $\times$ $g_{\mathrm{implied}}$) with 95\% bootstrapped CIs.
        The expected pattern: $\widehat{g}$ and $\widehat{A}$ should
        be higher in the high-$g_{\mathrm{implied}}$ conditions
        (where the model assigns a stronger marker), and lower in the
        high-pc\textsubscript{prag} conditions (credibility discounting).
  \item \textbf{Marker identity}: because the best-choice marker may
        differ across items within the same condition, a secondary plot
        shows ratings grouped by actual marker identity, allowing
        inspection of whether marker type (rather than just condition)
        drives the effects.
  \item \textbf{pc\textsubscript{prop} covariate}: scatter plots display
        each DV as a function of the continuous norming rating per
        condition, testing whether higher propositional acceptance
        amplifies the $\widehat{g}$ and $\widehat{A}$ signals for
        strong markers.
\end{itemize}

\subsubsection{Bayesian mixed-effects regression}

Two Bayesian mixed-effects regression models are fitted via
\texttt{brms} \citep{burkner2017}, one per DV ($\widehat{g}$ and
$\widehat{A}$):

\begin{align}
y &\sim \mathrm{Beta}(\mu, \phi)
\quad \text{(ratings rescaled to }(0,1)\text{)} \\
\mathrm{logit}(\mu) &= \alpha
+ \beta_{g} \cdot \mathrm{g\_implied\_c}
+ \beta_{\mathrm{prag}} \cdot \mathrm{pc\_prag\_c}
+ \beta_{\mathrm{prop}} \cdot \mathrm{pc\_prop\_c}
+ \beta_{\mathrm{int}} \cdot \mathrm{g\_implied\_c}
  \times \mathrm{pc\_prag\_c}
+ (1 \mid \mathrm{submission\_id})
+ (1 \mid \mathrm{topic}),
\end{align}

where \texttt{g\_implied\_c} is the $g$ level from the corresponding
production condition ($-0.5$ = low, $+0.5$ = high), effectively
coding the marker's position on the production utility scale.
A secondary specification replaces \texttt{g\_implied\_c} with
actualy marker identity as an unordered factor, to examine
non-monotonic patterns and verify that marker rather than
condition label drives effects.

Model comparison via approximate leave-one-out cross-validation (LOO)
selects among additive, interaction, and marker-as-categorical
formulations.

\subsubsection{RSA listener parameter estimation}
\label{subsub:rsa_fit}

The two empirical rating means (per condition cell: 4 conditions) constrain
the listener parameters of the noisy-threshold RSA model.
Specifically, the model predicts:

\begin{align}
\mathbb{E}[\widehat{g} \mid u, \text{context}]
  &= \int g \cdot P(g \mid u, pc_{\mathrm{prag}})\, dg,
\label{eq:pred_g} \\
P(\widehat{A} \mid u, \text{context})
  &= \iiint \sigma(\eta_0 + \eta_g g - \eta_{pc}\, pc
     - \eta_{\mathrm{int}}\, g \cdot pc)\;
     P(pc_{\mathrm{prop}}, pc_{\mathrm{prag}}, g \mid u)\,
     d{pc_{\mathrm{prop}}}\,d{pc_{\mathrm{prag}}}\,dg,
\label{eq:adoption}
\end{align}

\noindent where the posterior $P(pc_{\mathrm{prop}}, pc_{\mathrm{prag}}, g
\mid u)$ is conditioned on the context's \textsc{pc\textsubscript{prag}}
level as a known prior (since it is explicitly given to the listener).

The threshold parameters $(\mu_1,\mu_2, \sigma_1,\sigma_2,
\lambda)$ together with the marker-specific costs $(c_u)$ are estimated
jointly from the production experiment data via the Stan model.
Given those threshold estimates, the listener parameters
$(\eta_0, \eta_g, \eta_{pc}, \eta_{\mathrm{int}})$ are then estimated by
minimising the weighted squared deviation between the model-predicted
means (Eqs.~\ref{eq:pred_g}--\ref{eq:adoption}) and the empirical rating
means from the comprehension experiment, or via a full joint Bayesian
fit over both experiments simultaneously.

\subsection{Model Predictions}

The noisy-threshold RSA model makes the following predictions for
$\widehat{g}$ and $\widehat{A}$, which serve as pre-registered hypotheses.
In all cases, ``strong-marker condition'' refers to conditions 1 and 3
(high $g_{\mathrm{implied}}$, i.e.\ the marker selected by the production
model under a high-$g$ context), and ``weak-marker condition'' to conditions
0 and 2.

\begin{enumerate}
  \item \textbf{Marker encodes inferred goal $\widehat{g}$}: $\widehat{g}$
        ratings should be higher in strong-marker conditions than in
        weak-marker conditions, across both pc\textsubscript{prag} levels.
        This is the core inversion prediction: if the production model
        correctly captures speaker behaviour, listeners should be able to
        decode $g$ from the marker.

  \item \textbf{Marker encodes adoption $\widehat{A}$}: $\widehat{A}$
        ratings should likewise be higher in strong-marker conditions,
        reflecting that higher inferred $g$ translates into greater
        adoption pressure downstream.

  \item \textbf{pc\textsubscript{prag} modulates $\widehat{A}$ but not
        $\widehat{g}$}: because $g$ is orthogonal to pc\textsubscript{prag}
        in the model's utility function, the effect of the marker on
        $\widehat{g}$ should be largely the same in low- and high-
        pc\textsubscript{prag} contexts.
        By contrast, $\widehat{A}$ should be reduced in high-
        pc\textsubscript{prag} contexts via the interaction term
        $\eta_{\mathrm{int}}\,g \cdot pc$: even a strong-marker signal
        for high $g$ is partly cancelled when controversy is high.

  \item \textbf{Credibility discounting}: for strong-marker conditions
        presented in a high-pc\textsubscript{prag} context (condition~3),
        $\widehat{A}$ should be \emph{lower} than for the same marker in
        a low-pc\textsubscript{prag} context (condition~1).
        This is the model's signature prediction distinguishing Bayesian
        inversion from a simpler marker-strength heuristic: a sceptical
        listener who knows controversy is high will discount the adoption
        implication of even the strongest marker.
        The $g_{\mathrm{implied}} \times \textsc{pc\textsubscript{prag}}$
        interaction on $\widehat{A}$ provides the empirical test.
\end{enumerate}

\end{document}
.
%
% Compile stand-alone:
%   pdflatex comprehension_exp.tex
%
% Embed in main.tex:
%   % comprehension_exp.tex
% Stand-alone plan for the comprehension experiment.
% Can be embedded in main.tex via \input{comprehension_exp}.
%
% Compile stand-alone:
%   pdflatex comprehension_exp.tex
%
% Embed in main.tex:
%   \input{comprehension_exp}   % (inside a \section{Experiment} environment)

\documentclass[12pt]{article}
\usepackage{CSPArticleStyle}
\usepackage{csquotes}
\usepackage[margin=1in]{geometry}
\usepackage{booktabs}
\usepackage{parskip}
\usepackage{amsmath}
\usepackage{times}
\usepackage[T1]{fontenc}
\usepackage[utf8]{inputenc}
\usepackage{xcolor}
\usepackage{tabularray}
\usepackage{gb4e}
\noautomath

\title{Comprehension Experiment: Inferring Speaker Goal and Action Adoption
       from Consensus-Marker Use}
\author{}
\date{}

\begin{document}

%% ============================================================
%% When embedding in main.tex, delete from \documentclass to
%% \begin{document} and delete \end{document} at the bottom.
%% ============================================================

\section{Experiment: Consensus-Marker Comprehension}

\subsection{Overview}

The comprehension experiment is the listener-side counterpart of the
production experiment and tests whether listeners can invert the
speaker's marker choice to recover the latent context that motivated it.
Where the production experiment gave participants a context (fixing
\textsc{pc\textsubscript{prag}} and $g$) and asked which marker they
would use, the comprehension experiment gives participants a context
(fixing \textsc{pc\textsubscript{prag}}) together with a marker already
embedded in the sentence and asks what goal strength the speaker likely
had and how likely the listener is to adopt $q$.

The specific marker presented in each vignette is not chosen arbitrarily
from the three candidate markers. Our preregistered plan was to assign,
for each item$\times$condition cell, the \emph{Bayesian best-choice marker}:
the marker with the highest posterior probability under the cost-augmented
noisy-threshold RSA model fitted to the production data, given the context's
\textsc{pc\textsubscript{prag}} level and a specified $g$ level. This was
intended to mirror the production design closely and to ensure that the
comprehension stimuli reflect model-based marker use rather than arbitrary
marker$\times$item pairings.

Applying this strict modal-assignment rule to the fitted production model,
however, yielded a highly imbalanced stimulus set: most critical cells were
assigned the strongest marker, with very little variation across marker
categories. In addition, many of these modal decisions were supported only by
very small probability margins over the next-best alternative. Because such a
stimulus set would make the comprehension experiment substantially less
informative about marker-based listener inferences, we introduced a documented,
exploratory balancing step for stimulus construction. Starting from the same
posterior predictive probabilities, we selected a near-balanced assignment of
the three markers across the 32 critical cells while maximizing total posterior
support under the fitted production model. The strict modal assignment is
retained as the preregistered baseline, but the balanced assignment is used for
the actual comprehension materials.

\textsc{pc\textsubscript{prag}} is retained as an explicit context
manipulation (low vs.\ high), so that listeners can be expected to
have a well-defined prior on pragmatic controversy when interpreting
the marker. \textsc{pc\textsubscript{prop}} is not manipulated within
the experiment; it enters as a continuous topic-level covariate from
the norming study.

After reading each vignette, participants provide two slider ratings
operationalising the key listener-side quantities:
\begin{enumerate}
  \item \textbf{Inferred speaker goal strength} ($\widehat{g}$):
        how strongly does the speaker seem to want the listener to
        accept and act on $q$?
  \item \textbf{Adoption likelihood} ($\widehat{A}$):
        how likely is the listener to follow the recommended action $q$?
\end{enumerate}

\noindent Perceived controversy is \emph{not} elicited as a rating:
because \textsc{pc\textsubscript{prag}} is already manipulated as a
context factor, asking participants to rate controversy would be
redundant with respect to the manipulation check and circular with
respect to the research question. The theoretically central inference
to test is whether listeners can decode the speaker's \emph{goal} from
marker use, and whether that inference translates into adoption
behaviour.

The two slider measures, together with the \textsc{pc\textsubscript{prag}}
context and the \textsc{pc\textsubscript{prop}} covariate, provide the
data needed to fit the model's listener parameters
$(\eta_0, \eta_g, \eta_{pc}, \eta_{\mathrm{int}})$.

\subsection{Participants}

Participants will be recruited via Prolific and will be required to be
native speakers of German (language-at-home filter).
Data collection is planned after completion of the production experiment
(model fitting required to determine the best-choice markers).
A target of $N = 80$ participants is planned ($n = 20$ per list across
4 counterbalancing lists), each completing 16 trials (8 critical + 8
filler), matching the structure of the production experiment.

\subsection{Materials}

\subsubsection{Markers}

Each critical vignette is presented with one specific consensus marker
already embedded in the Mittelfeld of the sentence frame (see
example~(\ref{ex:frame_comp})):

\begin{exe}
  \ex \label{ex:frame_comp}
  Digitale Kompetenzen sind \textbf{[MARKER]}
  in der heutigen Arbeitswelt unverzichtbar.\\
  `Digital skills are \textbf{[MARKER]} indispensable in today's
  working world.'
\end{exe}

Unlike the production experiment, participants do \emph{not} choose the
marker; they read the sentence with the marker already in place and respond
to the two rating probes.

Critically, the marker assigned to each condition is not drawn arbitrarily
from the three candidates. Under the preregistered plan, for each of the four
production conditions (pc\textsubscript{prag} $\times$ $g$), the marker used in
comprehension was to be the \textbf{Bayesian best-choice marker}: the marker
$u^*$ with the highest posterior probability under the fitted
cost-augmented noisy-threshold RSA production model,
\[
u^* = \arg\max_{u} \; P(u \mid pc_{\mathrm{prag}}, g,
                           pc_{\mathrm{prop}\_c}).
\]
Because $pc_{\mathrm{prop}}$ varies continuously across items,
$u^*$ may differ per item even within the same
(pc\textsubscript{prag}, $g$) cell; the final stimulus list is therefore
constructed item-by-item once the production model has been fitted.
In practice, the fitted production model yielded a nearly degenerate modal map,
with most cells assigned to \textit{bekanntlich}. We therefore retained the
model-based logic of the procedure but replaced strict cellwise modal selection
with an uncertainty-aware balancing step: among the marker assignments with
high posterior support, we selected a near-balanced configuration across the 32
critical cells. This preserves the production--comprehension link while also
ensuring that the comprehension experiment contains sufficient marker variation
to test listener-side inferences.

All three markers remain candidates; the mapping from conditions to
markers is an empirical output of the production analysis, not a design
pre-commitment. In the final comprehension materials, we use the
\textbf{balanced model-based assignment}, rather than the strict modal
assignment, in order to maintain informative variation across marker types.

\subsubsection{Topic Items}

The same eight topic domains used in the norming study and the production
experiment are used here (Table~\ref{tab:items_comp}).
The 4-condition $\times$ 8-item Latin square (Section~\ref{subsub:counterbalancing_comp})
requires no additional items.

\begin{table}[ht]
\centering
\caption{Topic items: sentence frame for $p$ (with marker position [M])
         and consequent $q$.}
\label{tab:items_comp}
\small
\begin{tabular}{clll}
\toprule
\textbf{\#} & \textbf{Topic} & \textbf{Sentence frame} & \textbf{Consequent $q$} \\
\midrule
1 & Klimawandel &
  \textit{Der Klimawandel wird [M]} &
  auf Flugreisen verzichten \\
  & & \textit{durch menschliche Aktivit\"aten verursacht.} & \\[3pt]
2 & Ern\"ahrung &
  \textit{Eine \"uberwiegend pflanzliche Ern\"ahrung} &
  mehr pflanzliche Lebens- \\
  & & \textit{f\"ordert [M] die Gesundheit.} &
  mittel einbauen \\[3pt]
3 & Stadtverkehr &
  \textit{Ein ausgebautes Nahverkehrsnetz} &
  h\"aufiger auf \\
  & & \textit{entlastet [M] den Stadtverkehr} &
  \"offentliche Verkehrs- \\
  & & \textit{deutlich.} &
  mittel umsteigen \\[3pt]
4 & Digitalisierung &
  \textit{Digitale Kompetenzen sind [M] in} &
  regelm\"a\ss{}ig Weiter- \\
  & & \textit{der heutigen Arbeitswelt unverzichtbar.} &
  bildungsangebote nutzen \\[3pt]
5 & Schlaf &
  \textit{Ausreichend Schlaf ist [M] entscheidend} &
  auf feste Schlafzeiten \\
  & & \textit{f\"ur die kognitive Leistungsf\"ahigkeit.} &
  achten \\[3pt]
6 & Plastik &
  \textit{Einwegplastik schadet [M] den} &
  konsequent auf Einweg- \\
  & & \textit{Meeres\"okosystemen erheblich.} &
  plastik verzichten \\[3pt]
7 & Sport &
  \textit{Regelm\"a\ss{}ige k\"orperliche Bewegung} &
  regelm\"a\ss{}ig Sport in \\
  & & \textit{senkt [M] das Risiko f\"ur} &
  den Alltag integrieren \\
  & & \textit{Herz-Kreislauf-Erkrankungen deutlich.} & \\[3pt]
8 & Lokalkauf &
  \textit{Das Kaufen bei lokalen H\"andlern} &
  \"ofter bei lokalen \\
  & & \textit{st\"arkt [M] die regionale Wirtschaft} &
  Gesch\"aften einkaufen \\
  & & \textit{nachhaltig.} & \\
\bottomrule
\end{tabular}
\end{table}

\noindent [M] denotes the position of the pre-filled best-choice marker.

\subsubsection{Vignettes}

\paragraph{Structure.}
Each trial presents a vignette consisting of three parts:
(a) a context paragraph establishing the communicative situation and
fixing the \textsc{pc\textsubscript{prag}} level;
(b) the target sentence with the best-choice marker already embedded
in the Mittelfeld;
(c) the consequent action clause.
The two rating probes follow immediately below the vignette.

\paragraph{Propositional perceived controversy (\textsc{pc\textsubscript{prop}}).}
As in the production experiment, \textsc{pc\textsubscript{prop}} is not
manipulated within the comprehension experiment.
Topic-level propositional acceptance ratings from the norming study
(0--100 scale; e.g., \textit{Digitalisierung}: 85/100,
\textit{Lokalkauf}: 58/100) serve as a continuous, mean-centred
covariate \textsc{pc\textsubscript{prop}}$_c$ in analysis.
Context paragraphs present $p$ without editorial evaluation; variation in
listeners' individual assessments of propositional controversy is
captured by the norming-derived covariate.

\paragraph{Pragmatic perceived controversy (\textsc{pc\textsubscript{prag}}).}
\textsc{pc\textsubscript{prag}} is manipulated between item instances
using the same two context versions as in the production experiment:
the speaker's social circle either \emph{largely agrees} (low) or
\emph{holds divergent views} (high) about $p$.
For example, for the \textit{Digitalisierung} item:

\begin{itemize}
  \item $\textsc{pc}_\text{prag} = \text{low}$: \textit{Deine Kollegin spricht
        mit dir \"uber Berufsausbildung und Weiterbildung.
        In eurem Team teilen die meisten die Ansicht, dass digitale Kompetenzen
        in der heutigen Arbeitswelt unverzichtbar sind.}
  \item $\textsc{pc}_\text{prag} = \text{high}$: \textit{Deine Kollegin spricht
        mit dir \"uber Berufsausbildung und Weiterbildung.
        In eurem Team gehen die Meinungen dazu stark auseinander.}
\end{itemize}

The context is framed from the listener's perspective as a direct addressee,
mirroring the production study where participants adopt the speaker's
perspective.

\paragraph{Marker manipulation.}
The dependent measures of the production experiment (which marker is chosen)
become the independent variable here.
One of the three consensus markers is inserted into the Mittelfeld of the
target sentence.
A complete sample item for the \textit{Digitalisierung} topic
($\textsc{pc}_\text{prag} = \text{high}$, marker = \textit{bekanntlich};
norming $\textsc{pc}_\text{prop}$ rating: 85/100) is given in
(\ref{ex:comp_sample}):

\begin{exe}
  \ex \label{ex:comp_sample}
  \textit{Deine Kollegin spricht mit dir \"uber Berufsausbildung und
  Weiterbildung.
  In eurem Team gehen die Meinungen dazu stark auseinander.}

  \textit{Sie sagt zu dir:} Digitale Kompetenzen sind \textbf{bekanntlich}
  in der heutigen Arbeitswelt unverzichtbar.
  Deshalb solltest du regelm\"a\ss{}ig digitale Weiterbildungsangebote nutzen.

  \noindent\textit{[Rating probes follow.]}
\end{exe}

\subsection{Design}

\subsubsection{Factorial structure}

The comprehension experiment mirrors the production experiment's
$\mathbf{2}$~\textbf{(pc\textsubscript{prag}: low vs.\ high)}
$\mathbf{\times}$~$\mathbf{2}$~\textbf{($g_{\mathrm{implied}}$: low
vs.\ high) design}, yielding \textbf{4 conditions}.
The $g_{\mathrm{implied}}$ dimension is not manipulated via an explicit
goal instruction (as in production) but is instead conveyed entirely
through the choice of marker: for each condition, the marker presented
is the best-choice marker for the corresponding
(\textsc{pc\textsubscript{prag}},~$g$) cell from the production model.
\textsc{pc\textsubscript{prop}} enters as a continuous norming covariate.
With 8 topic items and a Latin square rotation across 4 lists, each
participant sees exactly \textbf{8 critical trials}---one item per
condition, each condition appearing twice---plus 8 filler trials,
for \textbf{16 trials} in total (Table~\ref{tab:design_comp}).

\begin{table}[ht]
\centering
\caption{%
  The $2\times2$ design (4 conditions), mirroring the production experiment.
  The marker presented in each condition is the Bayesian best-choice
  marker from the fitted production model; specific markers are
  determined after production data collection.}
\label{tab:design_comp}
\small
\begin{tabular}{ccll}
\toprule
\textbf{pc\textsubscript{prag}} &
\textbf{$g_{\mathrm{implied}}$} &
\textbf{Marker (TBD from production)} &
\textbf{Condition label} \\
\midrule
low  & low  & best-choice for $(\text{low}, g_{\text{low}})$  &
  (0) community agrees, weak signal \\
low  & high & best-choice for $(\text{low}, g_{\text{high}})$ &
  (1) community agrees, strong signal \\
high & low  & best-choice for $(\text{high}, g_{\text{low}})$ &
  (2) community divided, weak signal \\
high & high & best-choice for $(\text{high}, g_{\text{high}})$ &
  (3) community divided, strong signal \\
\bottomrule
\end{tabular}
\end{table}

\subsubsection{Counterbalancing}
\label{subsub:counterbalancing_comp}

The experiment uses the \textbf{same $4 \times 8$ Latin square} as the
production experiment: each of the 4 lists assigns item $i$
($i = 0,\ldots,7$) to condition $(i + l)\bmod 4$, where
$l \in \{0,1,2,3\}$ is the list number.
This guarantees that:
\begin{itemize}
  \item each participant sees every item exactly once;
  \item each participant sees each of the 4 conditions exactly twice
        (two items per condition), ensuring balanced coverage of the
        $2\times2$ design;
  \item across all 4 lists, each item appears equally often in each
        condition.
\end{itemize}

List assignment follows the same Prolific-and-URL strategy as the production
experiment: each Prolific study embeds \texttt{?list=N} (0--3) in the
experiment URL, capping recruitment at $n = 20$ participants per list
for a total of $N = 80$.

In addition to the 8 critical trials, each participant completed 8
\textbf{filler trials} using the same filler items as the production
experiment (train travel, reading light, cooking, coffee, vacation,
hydration, fresh air, noise), each presented with a randomly assigned
marker.
Fillers are excluded from analysis; they serve to dilute the critical
items and reduce demand characteristics.

\subsubsection{Dependent variables}
\label{subsub:dvs_comp}

Two slider ratings are collected after each vignette (0--100 continuous
scale with no visible numbers, labelled at the endpoints as shown):

\begin{enumerate}
  \item \textbf{Inferred speaker goal strength} (operationalises $\widehat{g}$):
  \begin{quote}
    \textit{Wie stark hast du den Eindruck, dass die Sprecherin m\"ochte,
    dass du ihre Empfehlung befolgst?}\\
    (`How strongly do you get the impression that the speaker wants you to
    follow their recommendation?')\\
    Endpoints: 0 = \textit{\"uberhaupt nicht}
    (`not at all') ---
    100 = \textit{sehr stark} (`very strongly').
  \end{quote}

  \item \textbf{Adoption likelihood of $q$} (operationalises $\widehat{A}$):
  \begin{quote}
    \textit{Wie wahrscheinlich ist es, dass du die empfohlene Handlung
    dann auch tats\"achlich umsetzt?}\\
    (`How likely is it that you will actually carry out the recommended
    action?')\\
    Endpoints: 0 = \textit{sehr unwahrscheinlich}
    (`very unlikely') ---
    100 = \textit{sehr wahrscheinlich} (`very likely').
  \end{quote}
\end{enumerate}

\noindent Perceived controversy is deliberately \emph{not} included as a
rating probe: \textsc{pc\textsubscript{prag}} is experimentally
manipulated via the context paragraph, so a controversy rating would
be redundant and risks obscuring the marker's contribution to inferred
goal strength.

Recorded per trial: \texttt{trial\_id}, \texttt{topic},
\texttt{is\_filler}, \texttt{condition\_index}, \texttt{marker},
\texttt{pc\_prag}, \texttt{g\_implied}, \texttt{rating\_g},
\texttt{rating\_adopt}, \texttt{rt} (ms); and once per participant:
\texttt{list\_num}.
The continuous \texttt{pc\_prop\_rating} covariate (topic-level norming
mean) is joined to the trial data during analysis.

\subsection{Procedure}

Participants complete the experiment online via the magpie platform
\citep{kochetkova2024magpie} in a browser.
After reading instructions in German, participants are assigned to one of
4 counterbalancing lists (via URL \texttt{?list=N} parameter assigned by
Prolific; no random fallback in deployment).
They then proceed through 16 trials---8 critical and 8 fillers---in a
fully randomised order, followed by a post-test questionnaire eliciting
native language, age, gender, education level, and optional comments.

Each trial displays:
(a) a context paragraph establishing the communicative situation and the
\textsc{pc\textsubscript{prag}} level;
(b) the target sentence with the best-choice marker already in place in
the Mittelfeld;
(c) the consequent action clause;
(d) the two rating sliders in fixed order ($\widehat{g}$, $\widehat{A}$),
each labelled in German.
The ``Weiter'' (Next) button is disabled until both sliders have been
moved at least once, ensuring a complete response on every trial.

\subsection{Planned Analysis}

\subsubsection{Empirical plots}

For each dependent variable ($\widehat{g}$, $\widehat{A}$):

\begin{itemize}
  \item \textbf{Condition means}: a dot-and-whisker plot shows mean
        ratings per condition (the $2\times2$ of pc\textsubscript{prag}
        $\times$ $g_{\mathrm{implied}}$) with 95\% bootstrapped CIs.
        The expected pattern: $\widehat{g}$ and $\widehat{A}$ should
        be higher in the high-$g_{\mathrm{implied}}$ conditions
        (where the model assigns a stronger marker), and lower in the
        high-pc\textsubscript{prag} conditions (credibility discounting).
  \item \textbf{Marker identity}: because the best-choice marker may
        differ across items within the same condition, a secondary plot
        shows ratings grouped by actual marker identity, allowing
        inspection of whether marker type (rather than just condition)
        drives the effects.
  \item \textbf{pc\textsubscript{prop} covariate}: scatter plots display
        each DV as a function of the continuous norming rating per
        condition, testing whether higher propositional acceptance
        amplifies the $\widehat{g}$ and $\widehat{A}$ signals for
        strong markers.
\end{itemize}

\subsubsection{Bayesian mixed-effects regression}

Two Bayesian mixed-effects regression models are fitted via
\texttt{brms} \citep{burkner2017}, one per DV ($\widehat{g}$ and
$\widehat{A}$):

\begin{align}
y &\sim \mathrm{Beta}(\mu, \phi)
\quad \text{(ratings rescaled to }(0,1)\text{)} \\
\mathrm{logit}(\mu) &= \alpha
+ \beta_{g} \cdot \mathrm{g\_implied\_c}
+ \beta_{\mathrm{prag}} \cdot \mathrm{pc\_prag\_c}
+ \beta_{\mathrm{prop}} \cdot \mathrm{pc\_prop\_c}
+ \beta_{\mathrm{int}} \cdot \mathrm{g\_implied\_c}
  \times \mathrm{pc\_prag\_c}
+ (1 \mid \mathrm{submission\_id})
+ (1 \mid \mathrm{topic}),
\end{align}

where \texttt{g\_implied\_c} is the $g$ level from the corresponding
production condition ($-0.5$ = low, $+0.5$ = high), effectively
coding the marker's position on the production utility scale.
A secondary specification replaces \texttt{g\_implied\_c} with
actualy marker identity as an unordered factor, to examine
non-monotonic patterns and verify that marker rather than
condition label drives effects.

Model comparison via approximate leave-one-out cross-validation (LOO)
selects among additive, interaction, and marker-as-categorical
formulations.

\subsubsection{RSA listener parameter estimation}
\label{subsub:rsa_fit}

The two empirical rating means (per condition cell: 4 conditions) constrain
the listener parameters of the noisy-threshold RSA model.
Specifically, the model predicts:

\begin{align}
\mathbb{E}[\widehat{g} \mid u, \text{context}]
  &= \int g \cdot P(g \mid u, pc_{\mathrm{prag}})\, dg,
\label{eq:pred_g} \\
P(\widehat{A} \mid u, \text{context})
  &= \iiint \sigma(\eta_0 + \eta_g g - \eta_{pc}\, pc
     - \eta_{\mathrm{int}}\, g \cdot pc)\;
     P(pc_{\mathrm{prop}}, pc_{\mathrm{prag}}, g \mid u)\,
     d{pc_{\mathrm{prop}}}\,d{pc_{\mathrm{prag}}}\,dg,
\label{eq:adoption}
\end{align}

\noindent where the posterior $P(pc_{\mathrm{prop}}, pc_{\mathrm{prag}}, g
\mid u)$ is conditioned on the context's \textsc{pc\textsubscript{prag}}
level as a known prior (since it is explicitly given to the listener).

The threshold parameters $(\mu_1,\mu_2, \sigma_1,\sigma_2,
\lambda)$ together with the marker-specific costs $(c_u)$ are estimated
jointly from the production experiment data via the Stan model.
Given those threshold estimates, the listener parameters
$(\eta_0, \eta_g, \eta_{pc}, \eta_{\mathrm{int}})$ are then estimated by
minimising the weighted squared deviation between the model-predicted
means (Eqs.~\ref{eq:pred_g}--\ref{eq:adoption}) and the empirical rating
means from the comprehension experiment, or via a full joint Bayesian
fit over both experiments simultaneously.

\subsection{Model Predictions}

The noisy-threshold RSA model makes the following predictions for
$\widehat{g}$ and $\widehat{A}$, which serve as pre-registered hypotheses.
In all cases, ``strong-marker condition'' refers to conditions 1 and 3
(high $g_{\mathrm{implied}}$, i.e.\ the marker selected by the production
model under a high-$g$ context), and ``weak-marker condition'' to conditions
0 and 2.

\begin{enumerate}
  \item \textbf{Marker encodes inferred goal $\widehat{g}$}: $\widehat{g}$
        ratings should be higher in strong-marker conditions than in
        weak-marker conditions, across both pc\textsubscript{prag} levels.
        This is the core inversion prediction: if the production model
        correctly captures speaker behaviour, listeners should be able to
        decode $g$ from the marker.

  \item \textbf{Marker encodes adoption $\widehat{A}$}: $\widehat{A}$
        ratings should likewise be higher in strong-marker conditions,
        reflecting that higher inferred $g$ translates into greater
        adoption pressure downstream.

  \item \textbf{pc\textsubscript{prag} modulates $\widehat{A}$ but not
        $\widehat{g}$}: because $g$ is orthogonal to pc\textsubscript{prag}
        in the model's utility function, the effect of the marker on
        $\widehat{g}$ should be largely the same in low- and high-
        pc\textsubscript{prag} contexts.
        By contrast, $\widehat{A}$ should be reduced in high-
        pc\textsubscript{prag} contexts via the interaction term
        $\eta_{\mathrm{int}}\,g \cdot pc$: even a strong-marker signal
        for high $g$ is partly cancelled when controversy is high.

  \item \textbf{Credibility discounting}: for strong-marker conditions
        presented in a high-pc\textsubscript{prag} context (condition~3),
        $\widehat{A}$ should be \emph{lower} than for the same marker in
        a low-pc\textsubscript{prag} context (condition~1).
        This is the model's signature prediction distinguishing Bayesian
        inversion from a simpler marker-strength heuristic: a sceptical
        listener who knows controversy is high will discount the adoption
        implication of even the strongest marker.
        The $g_{\mathrm{implied}} \times \textsc{pc\textsubscript{prag}}$
        interaction on $\widehat{A}$ provides the empirical test.
\end{enumerate}

\end{document}
   % (inside a \section{Experiment} environment)

\documentclass[12pt]{article}
\usepackage{CSPArticleStyle}
\usepackage{csquotes}
\usepackage[margin=1in]{geometry}
\usepackage{booktabs}
\usepackage{parskip}
\usepackage{amsmath}
\usepackage{times}
\usepackage[T1]{fontenc}
\usepackage[utf8]{inputenc}
\usepackage{xcolor}
\usepackage{tabularray}
\usepackage{gb4e}
\noautomath

\title{Comprehension Experiment: Inferring Speaker Goal and Action Adoption
       from Consensus-Marker Use}
\author{}
\date{}

\begin{document}

%% ============================================================
%% When embedding in main.tex, delete from \documentclass to
%% \begin{document} and delete \end{document} at the bottom.
%% ============================================================

\section{Experiment: Consensus-Marker Comprehension}

\subsection{Overview}

The comprehension experiment is the listener-side counterpart of the
production experiment and tests whether listeners can invert the
speaker's marker choice to recover the latent context that motivated it.
Where the production experiment gave participants a context (fixing
\textsc{pc\textsubscript{prag}} and $g$) and asked which marker they
would use, the comprehension experiment gives participants a context
(fixing \textsc{pc\textsubscript{prag}}) together with a marker already
embedded in the sentence and asks what goal strength the speaker likely
had and how likely the listener is to adopt $q$.

The specific marker presented in each vignette is not chosen arbitrarily
from the three candidate markers. Our preregistered plan was to assign,
for each item$\times$condition cell, the \emph{Bayesian best-choice marker}:
the marker with the highest posterior probability under the cost-augmented
noisy-threshold RSA model fitted to the production data, given the context's
\textsc{pc\textsubscript{prag}} level and a specified $g$ level. This was
intended to mirror the production design closely and to ensure that the
comprehension stimuli reflect model-based marker use rather than arbitrary
marker$\times$item pairings.

Applying this strict modal-assignment rule to the fitted production model,
however, yielded a highly imbalanced stimulus set: most critical cells were
assigned the strongest marker, with very little variation across marker
categories. In addition, many of these modal decisions were supported only by
very small probability margins over the next-best alternative. Because such a
stimulus set would make the comprehension experiment substantially less
informative about marker-based listener inferences, we introduced a documented,
exploratory balancing step for stimulus construction. Starting from the same
posterior predictive probabilities, we selected a near-balanced assignment of
the three markers across the 32 critical cells while maximizing total posterior
support under the fitted production model. The strict modal assignment is
retained as the preregistered baseline, but the balanced assignment is used for
the actual comprehension materials.

\textsc{pc\textsubscript{prag}} is retained as an explicit context
manipulation (low vs.\ high), so that listeners can be expected to
have a well-defined prior on pragmatic controversy when interpreting
the marker. \textsc{pc\textsubscript{prop}} is not manipulated within
the experiment; it enters as a continuous topic-level covariate from
the norming study.

After reading each vignette, participants provide two slider ratings
operationalising the key listener-side quantities:
\begin{enumerate}
  \item \textbf{Inferred speaker goal strength} ($\widehat{g}$):
        how strongly does the speaker seem to want the listener to
        accept and act on $q$?
  \item \textbf{Adoption likelihood} ($\widehat{A}$):
        how likely is the listener to follow the recommended action $q$?
\end{enumerate}

\noindent Perceived controversy is \emph{not} elicited as a rating:
because \textsc{pc\textsubscript{prag}} is already manipulated as a
context factor, asking participants to rate controversy would be
redundant with respect to the manipulation check and circular with
respect to the research question. The theoretically central inference
to test is whether listeners can decode the speaker's \emph{goal} from
marker use, and whether that inference translates into adoption
behaviour.

The two slider measures, together with the \textsc{pc\textsubscript{prag}}
context and the \textsc{pc\textsubscript{prop}} covariate, provide the
data needed to fit the model's listener parameters
$(\eta_0, \eta_g, \eta_{pc}, \eta_{\mathrm{int}})$.

\subsection{Participants}

Participants will be recruited via Prolific and will be required to be
native speakers of German (language-at-home filter).
Data collection is planned after completion of the production experiment
(model fitting required to determine the best-choice markers).
A target of $N = 80$ participants is planned ($n = 20$ per list across
4 counterbalancing lists), each completing 16 trials (8 critical + 8
filler), matching the structure of the production experiment.

\subsection{Materials}

\subsubsection{Markers}

Each critical vignette is presented with one specific consensus marker
already embedded in the Mittelfeld of the sentence frame (see
example~(\ref{ex:frame_comp})):

\begin{exe}
  \ex \label{ex:frame_comp}
  Digitale Kompetenzen sind \textbf{[MARKER]}
  in der heutigen Arbeitswelt unverzichtbar.\\
  `Digital skills are \textbf{[MARKER]} indispensable in today's
  working world.'
\end{exe}

Unlike the production experiment, participants do \emph{not} choose the
marker; they read the sentence with the marker already in place and respond
to the two rating probes.

Critically, the marker assigned to each condition is not drawn arbitrarily
from the three candidates. Under the preregistered plan, for each of the four
production conditions (pc\textsubscript{prag} $\times$ $g$), the marker used in
comprehension was to be the \textbf{Bayesian best-choice marker}: the marker
$u^*$ with the highest posterior probability under the fitted
cost-augmented noisy-threshold RSA production model,
\[
u^* = \arg\max_{u} \; P(u \mid pc_{\mathrm{prag}}, g,
                           pc_{\mathrm{prop}\_c}).
\]
Because $pc_{\mathrm{prop}}$ varies continuously across items,
$u^*$ may differ per item even within the same
(pc\textsubscript{prag}, $g$) cell; the final stimulus list is therefore
constructed item-by-item once the production model has been fitted.
In practice, the fitted production model yielded a nearly degenerate modal map,
with most cells assigned to \textit{bekanntlich}. We therefore retained the
model-based logic of the procedure but replaced strict cellwise modal selection
with an uncertainty-aware balancing step: among the marker assignments with
high posterior support, we selected a near-balanced configuration across the 32
critical cells. This preserves the production--comprehension link while also
ensuring that the comprehension experiment contains sufficient marker variation
to test listener-side inferences.

All three markers remain candidates; the mapping from conditions to
markers is an empirical output of the production analysis, not a design
pre-commitment. In the final comprehension materials, we use the
\textbf{balanced model-based assignment}, rather than the strict modal
assignment, in order to maintain informative variation across marker types.

\subsubsection{Topic Items}

The same eight topic domains used in the norming study and the production
experiment are used here (Table~\ref{tab:items_comp}).
The 4-condition $\times$ 8-item Latin square (Section~\ref{subsub:counterbalancing_comp})
requires no additional items.

\begin{table}[ht]
\centering
\caption{Topic items: sentence frame for $p$ (with marker position [M])
         and consequent $q$.}
\label{tab:items_comp}
\small
\begin{tabular}{clll}
\toprule
\textbf{\#} & \textbf{Topic} & \textbf{Sentence frame} & \textbf{Consequent $q$} \\
\midrule
1 & Klimawandel &
  \textit{Der Klimawandel wird [M]} &
  auf Flugreisen verzichten \\
  & & \textit{durch menschliche Aktivit\"aten verursacht.} & \\[3pt]
2 & Ern\"ahrung &
  \textit{Eine \"uberwiegend pflanzliche Ern\"ahrung} &
  mehr pflanzliche Lebens- \\
  & & \textit{f\"ordert [M] die Gesundheit.} &
  mittel einbauen \\[3pt]
3 & Stadtverkehr &
  \textit{Ein ausgebautes Nahverkehrsnetz} &
  h\"aufiger auf \\
  & & \textit{entlastet [M] den Stadtverkehr} &
  \"offentliche Verkehrs- \\
  & & \textit{deutlich.} &
  mittel umsteigen \\[3pt]
4 & Digitalisierung &
  \textit{Digitale Kompetenzen sind [M] in} &
  regelm\"a\ss{}ig Weiter- \\
  & & \textit{der heutigen Arbeitswelt unverzichtbar.} &
  bildungsangebote nutzen \\[3pt]
5 & Schlaf &
  \textit{Ausreichend Schlaf ist [M] entscheidend} &
  auf feste Schlafzeiten \\
  & & \textit{f\"ur die kognitive Leistungsf\"ahigkeit.} &
  achten \\[3pt]
6 & Plastik &
  \textit{Einwegplastik schadet [M] den} &
  konsequent auf Einweg- \\
  & & \textit{Meeres\"okosystemen erheblich.} &
  plastik verzichten \\[3pt]
7 & Sport &
  \textit{Regelm\"a\ss{}ige k\"orperliche Bewegung} &
  regelm\"a\ss{}ig Sport in \\
  & & \textit{senkt [M] das Risiko f\"ur} &
  den Alltag integrieren \\
  & & \textit{Herz-Kreislauf-Erkrankungen deutlich.} & \\[3pt]
8 & Lokalkauf &
  \textit{Das Kaufen bei lokalen H\"andlern} &
  \"ofter bei lokalen \\
  & & \textit{st\"arkt [M] die regionale Wirtschaft} &
  Gesch\"aften einkaufen \\
  & & \textit{nachhaltig.} & \\
\bottomrule
\end{tabular}
\end{table}

\noindent [M] denotes the position of the pre-filled best-choice marker.

\subsubsection{Vignettes}

\paragraph{Structure.}
Each trial presents a vignette consisting of three parts:
(a) a context paragraph establishing the communicative situation and
fixing the \textsc{pc\textsubscript{prag}} level;
(b) the target sentence with the best-choice marker already embedded
in the Mittelfeld;
(c) the consequent action clause.
The two rating probes follow immediately below the vignette.

\paragraph{Propositional perceived controversy (\textsc{pc\textsubscript{prop}}).}
As in the production experiment, \textsc{pc\textsubscript{prop}} is not
manipulated within the comprehension experiment.
Topic-level propositional acceptance ratings from the norming study
(0--100 scale; e.g., \textit{Digitalisierung}: 85/100,
\textit{Lokalkauf}: 58/100) serve as a continuous, mean-centred
covariate \textsc{pc\textsubscript{prop}}$_c$ in analysis.
Context paragraphs present $p$ without editorial evaluation; variation in
listeners' individual assessments of propositional controversy is
captured by the norming-derived covariate.

\paragraph{Pragmatic perceived controversy (\textsc{pc\textsubscript{prag}}).}
\textsc{pc\textsubscript{prag}} is manipulated between item instances
using the same two context versions as in the production experiment:
the speaker's social circle either \emph{largely agrees} (low) or
\emph{holds divergent views} (high) about $p$.
For example, for the \textit{Digitalisierung} item:

\begin{itemize}
  \item $\textsc{pc}_\text{prag} = \text{low}$: \textit{Deine Kollegin spricht
        mit dir \"uber Berufsausbildung und Weiterbildung.
        In eurem Team teilen die meisten die Ansicht, dass digitale Kompetenzen
        in der heutigen Arbeitswelt unverzichtbar sind.}
  \item $\textsc{pc}_\text{prag} = \text{high}$: \textit{Deine Kollegin spricht
        mit dir \"uber Berufsausbildung und Weiterbildung.
        In eurem Team gehen die Meinungen dazu stark auseinander.}
\end{itemize}

The context is framed from the listener's perspective as a direct addressee,
mirroring the production study where participants adopt the speaker's
perspective.

\paragraph{Marker manipulation.}
The dependent measures of the production experiment (which marker is chosen)
become the independent variable here.
One of the three consensus markers is inserted into the Mittelfeld of the
target sentence.
A complete sample item for the \textit{Digitalisierung} topic
($\textsc{pc}_\text{prag} = \text{high}$, marker = \textit{bekanntlich};
norming $\textsc{pc}_\text{prop}$ rating: 85/100) is given in
(\ref{ex:comp_sample}):

\begin{exe}
  \ex \label{ex:comp_sample}
  \textit{Deine Kollegin spricht mit dir \"uber Berufsausbildung und
  Weiterbildung.
  In eurem Team gehen die Meinungen dazu stark auseinander.}

  \textit{Sie sagt zu dir:} Digitale Kompetenzen sind \textbf{bekanntlich}
  in der heutigen Arbeitswelt unverzichtbar.
  Deshalb solltest du regelm\"a\ss{}ig digitale Weiterbildungsangebote nutzen.

  \noindent\textit{[Rating probes follow.]}
\end{exe}

\subsection{Design}

\subsubsection{Factorial structure}

The comprehension experiment mirrors the production experiment's
$\mathbf{2}$~\textbf{(pc\textsubscript{prag}: low vs.\ high)}
$\mathbf{\times}$~$\mathbf{2}$~\textbf{($g_{\mathrm{implied}}$: low
vs.\ high) design}, yielding \textbf{4 conditions}.
The $g_{\mathrm{implied}}$ dimension is not manipulated via an explicit
goal instruction (as in production) but is instead conveyed entirely
through the choice of marker: for each condition, the marker presented
is the best-choice marker for the corresponding
(\textsc{pc\textsubscript{prag}},~$g$) cell from the production model.
\textsc{pc\textsubscript{prop}} enters as a continuous norming covariate.
With 8 topic items and a Latin square rotation across 4 lists, each
participant sees exactly \textbf{8 critical trials}---one item per
condition, each condition appearing twice---plus 8 filler trials,
for \textbf{16 trials} in total (Table~\ref{tab:design_comp}).

\begin{table}[ht]
\centering
\caption{%
  The $2\times2$ design (4 conditions), mirroring the production experiment.
  The marker presented in each condition is the Bayesian best-choice
  marker from the fitted production model; specific markers are
  determined after production data collection.}
\label{tab:design_comp}
\small
\begin{tabular}{ccll}
\toprule
\textbf{pc\textsubscript{prag}} &
\textbf{$g_{\mathrm{implied}}$} &
\textbf{Marker (TBD from production)} &
\textbf{Condition label} \\
\midrule
low  & low  & best-choice for $(\text{low}, g_{\text{low}})$  &
  (0) community agrees, weak signal \\
low  & high & best-choice for $(\text{low}, g_{\text{high}})$ &
  (1) community agrees, strong signal \\
high & low  & best-choice for $(\text{high}, g_{\text{low}})$ &
  (2) community divided, weak signal \\
high & high & best-choice for $(\text{high}, g_{\text{high}})$ &
  (3) community divided, strong signal \\
\bottomrule
\end{tabular}
\end{table}

\subsubsection{Counterbalancing}
\label{subsub:counterbalancing_comp}

The experiment uses the \textbf{same $4 \times 8$ Latin square} as the
production experiment: each of the 4 lists assigns item $i$
($i = 0,\ldots,7$) to condition $(i + l)\bmod 4$, where
$l \in \{0,1,2,3\}$ is the list number.
This guarantees that:
\begin{itemize}
  \item each participant sees every item exactly once;
  \item each participant sees each of the 4 conditions exactly twice
        (two items per condition), ensuring balanced coverage of the
        $2\times2$ design;
  \item across all 4 lists, each item appears equally often in each
        condition.
\end{itemize}

List assignment follows the same Prolific-and-URL strategy as the production
experiment: each Prolific study embeds \texttt{?list=N} (0--3) in the
experiment URL, capping recruitment at $n = 20$ participants per list
for a total of $N = 80$.

In addition to the 8 critical trials, each participant completed 8
\textbf{filler trials} using the same filler items as the production
experiment (train travel, reading light, cooking, coffee, vacation,
hydration, fresh air, noise), each presented with a randomly assigned
marker.
Fillers are excluded from analysis; they serve to dilute the critical
items and reduce demand characteristics.

\subsubsection{Dependent variables}
\label{subsub:dvs_comp}

Two slider ratings are collected after each vignette (0--100 continuous
scale with no visible numbers, labelled at the endpoints as shown):

\begin{enumerate}
  \item \textbf{Inferred speaker goal strength} (operationalises $\widehat{g}$):
  \begin{quote}
    \textit{Wie stark hast du den Eindruck, dass die Sprecherin m\"ochte,
    dass du ihre Empfehlung befolgst?}\\
    (`How strongly do you get the impression that the speaker wants you to
    follow their recommendation?')\\
    Endpoints: 0 = \textit{\"uberhaupt nicht}
    (`not at all') ---
    100 = \textit{sehr stark} (`very strongly').
  \end{quote}

  \item \textbf{Adoption likelihood of $q$} (operationalises $\widehat{A}$):
  \begin{quote}
    \textit{Wie wahrscheinlich ist es, dass du die empfohlene Handlung
    dann auch tats\"achlich umsetzt?}\\
    (`How likely is it that you will actually carry out the recommended
    action?')\\
    Endpoints: 0 = \textit{sehr unwahrscheinlich}
    (`very unlikely') ---
    100 = \textit{sehr wahrscheinlich} (`very likely').
  \end{quote}
\end{enumerate}

\noindent Perceived controversy is deliberately \emph{not} included as a
rating probe: \textsc{pc\textsubscript{prag}} is experimentally
manipulated via the context paragraph, so a controversy rating would
be redundant and risks obscuring the marker's contribution to inferred
goal strength.

Recorded per trial: \texttt{trial\_id}, \texttt{topic},
\texttt{is\_filler}, \texttt{condition\_index}, \texttt{marker},
\texttt{pc\_prag}, \texttt{g\_implied}, \texttt{rating\_g},
\texttt{rating\_adopt}, \texttt{rt} (ms); and once per participant:
\texttt{list\_num}.
The continuous \texttt{pc\_prop\_rating} covariate (topic-level norming
mean) is joined to the trial data during analysis.

\subsection{Procedure}

Participants complete the experiment online via the magpie platform
\citep{kochetkova2024magpie} in a browser.
After reading instructions in German, participants are assigned to one of
4 counterbalancing lists (via URL \texttt{?list=N} parameter assigned by
Prolific; no random fallback in deployment).
They then proceed through 16 trials---8 critical and 8 fillers---in a
fully randomised order, followed by a post-test questionnaire eliciting
native language, age, gender, education level, and optional comments.

Each trial displays:
(a) a context paragraph establishing the communicative situation and the
\textsc{pc\textsubscript{prag}} level;
(b) the target sentence with the best-choice marker already in place in
the Mittelfeld;
(c) the consequent action clause;
(d) the two rating sliders in fixed order ($\widehat{g}$, $\widehat{A}$),
each labelled in German.
The ``Weiter'' (Next) button is disabled until both sliders have been
moved at least once, ensuring a complete response on every trial.

\subsection{Planned Analysis}

\subsubsection{Empirical plots}

For each dependent variable ($\widehat{g}$, $\widehat{A}$):

\begin{itemize}
  \item \textbf{Condition means}: a dot-and-whisker plot shows mean
        ratings per condition (the $2\times2$ of pc\textsubscript{prag}
        $\times$ $g_{\mathrm{implied}}$) with 95\% bootstrapped CIs.
        The expected pattern: $\widehat{g}$ and $\widehat{A}$ should
        be higher in the high-$g_{\mathrm{implied}}$ conditions
        (where the model assigns a stronger marker), and lower in the
        high-pc\textsubscript{prag} conditions (credibility discounting).
  \item \textbf{Marker identity}: because the best-choice marker may
        differ across items within the same condition, a secondary plot
        shows ratings grouped by actual marker identity, allowing
        inspection of whether marker type (rather than just condition)
        drives the effects.
  \item \textbf{pc\textsubscript{prop} covariate}: scatter plots display
        each DV as a function of the continuous norming rating per
        condition, testing whether higher propositional acceptance
        amplifies the $\widehat{g}$ and $\widehat{A}$ signals for
        strong markers.
\end{itemize}

\subsubsection{Bayesian mixed-effects regression}

Two Bayesian mixed-effects regression models are fitted via
\texttt{brms} \citep{burkner2017}, one per DV ($\widehat{g}$ and
$\widehat{A}$):

\begin{align}
y &\sim \mathrm{Beta}(\mu, \phi)
\quad \text{(ratings rescaled to }(0,1)\text{)} \\
\mathrm{logit}(\mu) &= \alpha
+ \beta_{g} \cdot \mathrm{g\_implied\_c}
+ \beta_{\mathrm{prag}} \cdot \mathrm{pc\_prag\_c}
+ \beta_{\mathrm{prop}} \cdot \mathrm{pc\_prop\_c}
+ \beta_{\mathrm{int}} \cdot \mathrm{g\_implied\_c}
  \times \mathrm{pc\_prag\_c}
+ (1 \mid \mathrm{submission\_id})
+ (1 \mid \mathrm{topic}),
\end{align}

where \texttt{g\_implied\_c} is the $g$ level from the corresponding
production condition ($-0.5$ = low, $+0.5$ = high), effectively
coding the marker's position on the production utility scale.
A secondary specification replaces \texttt{g\_implied\_c} with
actualy marker identity as an unordered factor, to examine
non-monotonic patterns and verify that marker rather than
condition label drives effects.

Model comparison via approximate leave-one-out cross-validation (LOO)
selects among additive, interaction, and marker-as-categorical
formulations.

\subsubsection{RSA listener parameter estimation}
\label{subsub:rsa_fit}

The two empirical rating means (per condition cell: 4 conditions) constrain
the listener parameters of the noisy-threshold RSA model.
Specifically, the model predicts:

\begin{align}
\mathbb{E}[\widehat{g} \mid u, \text{context}]
  &= \int g \cdot P(g \mid u, pc_{\mathrm{prag}})\, dg,
\label{eq:pred_g} \\
P(\widehat{A} \mid u, \text{context})
  &= \iiint \sigma(\eta_0 + \eta_g g - \eta_{pc}\, pc
     - \eta_{\mathrm{int}}\, g \cdot pc)\;
     P(pc_{\mathrm{prop}}, pc_{\mathrm{prag}}, g \mid u)\,
     d{pc_{\mathrm{prop}}}\,d{pc_{\mathrm{prag}}}\,dg,
\label{eq:adoption}
\end{align}

\noindent where the posterior $P(pc_{\mathrm{prop}}, pc_{\mathrm{prag}}, g
\mid u)$ is conditioned on the context's \textsc{pc\textsubscript{prag}}
level as a known prior (since it is explicitly given to the listener).

The threshold parameters $(\mu_1,\mu_2, \sigma_1,\sigma_2,
\lambda)$ together with the marker-specific costs $(c_u)$ are estimated
jointly from the production experiment data via the Stan model.
Given those threshold estimates, the listener parameters
$(\eta_0, \eta_g, \eta_{pc}, \eta_{\mathrm{int}})$ are then estimated by
minimising the weighted squared deviation between the model-predicted
means (Eqs.~\ref{eq:pred_g}--\ref{eq:adoption}) and the empirical rating
means from the comprehension experiment, or via a full joint Bayesian
fit over both experiments simultaneously.

\subsection{Model Predictions}

The noisy-threshold RSA model makes the following predictions for
$\widehat{g}$ and $\widehat{A}$, which serve as pre-registered hypotheses.
In all cases, ``strong-marker condition'' refers to conditions 1 and 3
(high $g_{\mathrm{implied}}$, i.e.\ the marker selected by the production
model under a high-$g$ context), and ``weak-marker condition'' to conditions
0 and 2.

\begin{enumerate}
  \item \textbf{Marker encodes inferred goal $\widehat{g}$}: $\widehat{g}$
        ratings should be higher in strong-marker conditions than in
        weak-marker conditions, across both pc\textsubscript{prag} levels.
        This is the core inversion prediction: if the production model
        correctly captures speaker behaviour, listeners should be able to
        decode $g$ from the marker.

  \item \textbf{Marker encodes adoption $\widehat{A}$}: $\widehat{A}$
        ratings should likewise be higher in strong-marker conditions,
        reflecting that higher inferred $g$ translates into greater
        adoption pressure downstream.

  \item \textbf{pc\textsubscript{prag} modulates $\widehat{A}$ but not
        $\widehat{g}$}: because $g$ is orthogonal to pc\textsubscript{prag}
        in the model's utility function, the effect of the marker on
        $\widehat{g}$ should be largely the same in low- and high-
        pc\textsubscript{prag} contexts.
        By contrast, $\widehat{A}$ should be reduced in high-
        pc\textsubscript{prag} contexts via the interaction term
        $\eta_{\mathrm{int}}\,g \cdot pc$: even a strong-marker signal
        for high $g$ is partly cancelled when controversy is high.

  \item \textbf{Credibility discounting}: for strong-marker conditions
        presented in a high-pc\textsubscript{prag} context (condition~3),
        $\widehat{A}$ should be \emph{lower} than for the same marker in
        a low-pc\textsubscript{prag} context (condition~1).
        This is the model's signature prediction distinguishing Bayesian
        inversion from a simpler marker-strength heuristic: a sceptical
        listener who knows controversy is high will discount the adoption
        implication of even the strongest marker.
        The $g_{\mathrm{implied}} \times \textsc{pc\textsubscript{prag}}$
        interaction on $\widehat{A}$ provides the empirical test.
\end{enumerate}

\end{document}
.
%
% Compile stand-alone:
%   pdflatex comprehension_exp.tex
%
% Embed in main.tex:
%   % comprehension_exp.tex
% Stand-alone plan for the comprehension experiment.
% Can be embedded in main.tex via % comprehension_exp.tex
% Stand-alone plan for the comprehension experiment.
% Can be embedded in main.tex via \input{comprehension_exp}.
%
% Compile stand-alone:
%   pdflatex comprehension_exp.tex
%
% Embed in main.tex:
%   \input{comprehension_exp}   % (inside a \section{Experiment} environment)

\documentclass[12pt]{article}
\usepackage{CSPArticleStyle}
\usepackage{csquotes}
\usepackage[margin=1in]{geometry}
\usepackage{booktabs}
\usepackage{parskip}
\usepackage{amsmath}
\usepackage{times}
\usepackage[T1]{fontenc}
\usepackage[utf8]{inputenc}
\usepackage{xcolor}
\usepackage{tabularray}
\usepackage{gb4e}
\noautomath

\title{Comprehension Experiment: Inferring Speaker Goal and Action Adoption
       from Consensus-Marker Use}
\author{}
\date{}

\begin{document}

%% ============================================================
%% When embedding in main.tex, delete from \documentclass to
%% \begin{document} and delete \end{document} at the bottom.
%% ============================================================

\section{Experiment: Consensus-Marker Comprehension}

\subsection{Overview}

The comprehension experiment is the listener-side counterpart of the
production experiment and tests whether listeners can invert the
speaker's marker choice to recover the latent context that motivated it.
Where the production experiment gave participants a context (fixing
\textsc{pc\textsubscript{prag}} and $g$) and asked which marker they
would use, the comprehension experiment gives participants a context
(fixing \textsc{pc\textsubscript{prag}}) together with a marker already
embedded in the sentence and asks what goal strength the speaker likely
had and how likely the listener is to adopt $q$.

The specific marker presented in each vignette is not chosen arbitrarily
from the three candidate markers. Our preregistered plan was to assign,
for each item$\times$condition cell, the \emph{Bayesian best-choice marker}:
the marker with the highest posterior probability under the cost-augmented
noisy-threshold RSA model fitted to the production data, given the context's
\textsc{pc\textsubscript{prag}} level and a specified $g$ level. This was
intended to mirror the production design closely and to ensure that the
comprehension stimuli reflect model-based marker use rather than arbitrary
marker$\times$item pairings.

Applying this strict modal-assignment rule to the fitted production model,
however, yielded a highly imbalanced stimulus set: most critical cells were
assigned the strongest marker, with very little variation across marker
categories. In addition, many of these modal decisions were supported only by
very small probability margins over the next-best alternative. Because such a
stimulus set would make the comprehension experiment substantially less
informative about marker-based listener inferences, we introduced a documented,
exploratory balancing step for stimulus construction. Starting from the same
posterior predictive probabilities, we selected a near-balanced assignment of
the three markers across the 32 critical cells while maximizing total posterior
support under the fitted production model. The strict modal assignment is
retained as the preregistered baseline, but the balanced assignment is used for
the actual comprehension materials.

\textsc{pc\textsubscript{prag}} is retained as an explicit context
manipulation (low vs.\ high), so that listeners can be expected to
have a well-defined prior on pragmatic controversy when interpreting
the marker. \textsc{pc\textsubscript{prop}} is not manipulated within
the experiment; it enters as a continuous topic-level covariate from
the norming study.

After reading each vignette, participants provide two slider ratings
operationalising the key listener-side quantities:
\begin{enumerate}
  \item \textbf{Inferred speaker goal strength} ($\widehat{g}$):
        how strongly does the speaker seem to want the listener to
        accept and act on $q$?
  \item \textbf{Adoption likelihood} ($\widehat{A}$):
        how likely is the listener to follow the recommended action $q$?
\end{enumerate}

\noindent Perceived controversy is \emph{not} elicited as a rating:
because \textsc{pc\textsubscript{prag}} is already manipulated as a
context factor, asking participants to rate controversy would be
redundant with respect to the manipulation check and circular with
respect to the research question. The theoretically central inference
to test is whether listeners can decode the speaker's \emph{goal} from
marker use, and whether that inference translates into adoption
behaviour.

The two slider measures, together with the \textsc{pc\textsubscript{prag}}
context and the \textsc{pc\textsubscript{prop}} covariate, provide the
data needed to fit the model's listener parameters
$(\eta_0, \eta_g, \eta_{pc}, \eta_{\mathrm{int}})$.

\subsection{Participants}

Participants will be recruited via Prolific and will be required to be
native speakers of German (language-at-home filter).
Data collection is planned after completion of the production experiment
(model fitting required to determine the best-choice markers).
A target of $N = 80$ participants is planned ($n = 20$ per list across
4 counterbalancing lists), each completing 16 trials (8 critical + 8
filler), matching the structure of the production experiment.

\subsection{Materials}

\subsubsection{Markers}

Each critical vignette is presented with one specific consensus marker
already embedded in the Mittelfeld of the sentence frame (see
example~(\ref{ex:frame_comp})):

\begin{exe}
  \ex \label{ex:frame_comp}
  Digitale Kompetenzen sind \textbf{[MARKER]}
  in der heutigen Arbeitswelt unverzichtbar.\\
  `Digital skills are \textbf{[MARKER]} indispensable in today's
  working world.'
\end{exe}

Unlike the production experiment, participants do \emph{not} choose the
marker; they read the sentence with the marker already in place and respond
to the two rating probes.

Critically, the marker assigned to each condition is not drawn arbitrarily
from the three candidates. Under the preregistered plan, for each of the four
production conditions (pc\textsubscript{prag} $\times$ $g$), the marker used in
comprehension was to be the \textbf{Bayesian best-choice marker}: the marker
$u^*$ with the highest posterior probability under the fitted
cost-augmented noisy-threshold RSA production model,
\[
u^* = \arg\max_{u} \; P(u \mid pc_{\mathrm{prag}}, g,
                           pc_{\mathrm{prop}\_c}).
\]
Because $pc_{\mathrm{prop}}$ varies continuously across items,
$u^*$ may differ per item even within the same
(pc\textsubscript{prag}, $g$) cell; the final stimulus list is therefore
constructed item-by-item once the production model has been fitted.
In practice, the fitted production model yielded a nearly degenerate modal map,
with most cells assigned to \textit{bekanntlich}. We therefore retained the
model-based logic of the procedure but replaced strict cellwise modal selection
with an uncertainty-aware balancing step: among the marker assignments with
high posterior support, we selected a near-balanced configuration across the 32
critical cells. This preserves the production--comprehension link while also
ensuring that the comprehension experiment contains sufficient marker variation
to test listener-side inferences.

All three markers remain candidates; the mapping from conditions to
markers is an empirical output of the production analysis, not a design
pre-commitment. In the final comprehension materials, we use the
\textbf{balanced model-based assignment}, rather than the strict modal
assignment, in order to maintain informative variation across marker types.

\subsubsection{Topic Items}

The same eight topic domains used in the norming study and the production
experiment are used here (Table~\ref{tab:items_comp}).
The 4-condition $\times$ 8-item Latin square (Section~\ref{subsub:counterbalancing_comp})
requires no additional items.

\begin{table}[ht]
\centering
\caption{Topic items: sentence frame for $p$ (with marker position [M])
         and consequent $q$.}
\label{tab:items_comp}
\small
\begin{tabular}{clll}
\toprule
\textbf{\#} & \textbf{Topic} & \textbf{Sentence frame} & \textbf{Consequent $q$} \\
\midrule
1 & Klimawandel &
  \textit{Der Klimawandel wird [M]} &
  auf Flugreisen verzichten \\
  & & \textit{durch menschliche Aktivit\"aten verursacht.} & \\[3pt]
2 & Ern\"ahrung &
  \textit{Eine \"uberwiegend pflanzliche Ern\"ahrung} &
  mehr pflanzliche Lebens- \\
  & & \textit{f\"ordert [M] die Gesundheit.} &
  mittel einbauen \\[3pt]
3 & Stadtverkehr &
  \textit{Ein ausgebautes Nahverkehrsnetz} &
  h\"aufiger auf \\
  & & \textit{entlastet [M] den Stadtverkehr} &
  \"offentliche Verkehrs- \\
  & & \textit{deutlich.} &
  mittel umsteigen \\[3pt]
4 & Digitalisierung &
  \textit{Digitale Kompetenzen sind [M] in} &
  regelm\"a\ss{}ig Weiter- \\
  & & \textit{der heutigen Arbeitswelt unverzichtbar.} &
  bildungsangebote nutzen \\[3pt]
5 & Schlaf &
  \textit{Ausreichend Schlaf ist [M] entscheidend} &
  auf feste Schlafzeiten \\
  & & \textit{f\"ur die kognitive Leistungsf\"ahigkeit.} &
  achten \\[3pt]
6 & Plastik &
  \textit{Einwegplastik schadet [M] den} &
  konsequent auf Einweg- \\
  & & \textit{Meeres\"okosystemen erheblich.} &
  plastik verzichten \\[3pt]
7 & Sport &
  \textit{Regelm\"a\ss{}ige k\"orperliche Bewegung} &
  regelm\"a\ss{}ig Sport in \\
  & & \textit{senkt [M] das Risiko f\"ur} &
  den Alltag integrieren \\
  & & \textit{Herz-Kreislauf-Erkrankungen deutlich.} & \\[3pt]
8 & Lokalkauf &
  \textit{Das Kaufen bei lokalen H\"andlern} &
  \"ofter bei lokalen \\
  & & \textit{st\"arkt [M] die regionale Wirtschaft} &
  Gesch\"aften einkaufen \\
  & & \textit{nachhaltig.} & \\
\bottomrule
\end{tabular}
\end{table}

\noindent [M] denotes the position of the pre-filled best-choice marker.

\subsubsection{Vignettes}

\paragraph{Structure.}
Each trial presents a vignette consisting of three parts:
(a) a context paragraph establishing the communicative situation and
fixing the \textsc{pc\textsubscript{prag}} level;
(b) the target sentence with the best-choice marker already embedded
in the Mittelfeld;
(c) the consequent action clause.
The two rating probes follow immediately below the vignette.

\paragraph{Propositional perceived controversy (\textsc{pc\textsubscript{prop}}).}
As in the production experiment, \textsc{pc\textsubscript{prop}} is not
manipulated within the comprehension experiment.
Topic-level propositional acceptance ratings from the norming study
(0--100 scale; e.g., \textit{Digitalisierung}: 85/100,
\textit{Lokalkauf}: 58/100) serve as a continuous, mean-centred
covariate \textsc{pc\textsubscript{prop}}$_c$ in analysis.
Context paragraphs present $p$ without editorial evaluation; variation in
listeners' individual assessments of propositional controversy is
captured by the norming-derived covariate.

\paragraph{Pragmatic perceived controversy (\textsc{pc\textsubscript{prag}}).}
\textsc{pc\textsubscript{prag}} is manipulated between item instances
using the same two context versions as in the production experiment:
the speaker's social circle either \emph{largely agrees} (low) or
\emph{holds divergent views} (high) about $p$.
For example, for the \textit{Digitalisierung} item:

\begin{itemize}
  \item $\textsc{pc}_\text{prag} = \text{low}$: \textit{Deine Kollegin spricht
        mit dir \"uber Berufsausbildung und Weiterbildung.
        In eurem Team teilen die meisten die Ansicht, dass digitale Kompetenzen
        in der heutigen Arbeitswelt unverzichtbar sind.}
  \item $\textsc{pc}_\text{prag} = \text{high}$: \textit{Deine Kollegin spricht
        mit dir \"uber Berufsausbildung und Weiterbildung.
        In eurem Team gehen die Meinungen dazu stark auseinander.}
\end{itemize}

The context is framed from the listener's perspective as a direct addressee,
mirroring the production study where participants adopt the speaker's
perspective.

\paragraph{Marker manipulation.}
The dependent measures of the production experiment (which marker is chosen)
become the independent variable here.
One of the three consensus markers is inserted into the Mittelfeld of the
target sentence.
A complete sample item for the \textit{Digitalisierung} topic
($\textsc{pc}_\text{prag} = \text{high}$, marker = \textit{bekanntlich};
norming $\textsc{pc}_\text{prop}$ rating: 85/100) is given in
(\ref{ex:comp_sample}):

\begin{exe}
  \ex \label{ex:comp_sample}
  \textit{Deine Kollegin spricht mit dir \"uber Berufsausbildung und
  Weiterbildung.
  In eurem Team gehen die Meinungen dazu stark auseinander.}

  \textit{Sie sagt zu dir:} Digitale Kompetenzen sind \textbf{bekanntlich}
  in der heutigen Arbeitswelt unverzichtbar.
  Deshalb solltest du regelm\"a\ss{}ig digitale Weiterbildungsangebote nutzen.

  \noindent\textit{[Rating probes follow.]}
\end{exe}

\subsection{Design}

\subsubsection{Factorial structure}

The comprehension experiment mirrors the production experiment's
$\mathbf{2}$~\textbf{(pc\textsubscript{prag}: low vs.\ high)}
$\mathbf{\times}$~$\mathbf{2}$~\textbf{($g_{\mathrm{implied}}$: low
vs.\ high) design}, yielding \textbf{4 conditions}.
The $g_{\mathrm{implied}}$ dimension is not manipulated via an explicit
goal instruction (as in production) but is instead conveyed entirely
through the choice of marker: for each condition, the marker presented
is the best-choice marker for the corresponding
(\textsc{pc\textsubscript{prag}},~$g$) cell from the production model.
\textsc{pc\textsubscript{prop}} enters as a continuous norming covariate.
With 8 topic items and a Latin square rotation across 4 lists, each
participant sees exactly \textbf{8 critical trials}---one item per
condition, each condition appearing twice---plus 8 filler trials,
for \textbf{16 trials} in total (Table~\ref{tab:design_comp}).

\begin{table}[ht]
\centering
\caption{%
  The $2\times2$ design (4 conditions), mirroring the production experiment.
  The marker presented in each condition is the Bayesian best-choice
  marker from the fitted production model; specific markers are
  determined after production data collection.}
\label{tab:design_comp}
\small
\begin{tabular}{ccll}
\toprule
\textbf{pc\textsubscript{prag}} &
\textbf{$g_{\mathrm{implied}}$} &
\textbf{Marker (TBD from production)} &
\textbf{Condition label} \\
\midrule
low  & low  & best-choice for $(\text{low}, g_{\text{low}})$  &
  (0) community agrees, weak signal \\
low  & high & best-choice for $(\text{low}, g_{\text{high}})$ &
  (1) community agrees, strong signal \\
high & low  & best-choice for $(\text{high}, g_{\text{low}})$ &
  (2) community divided, weak signal \\
high & high & best-choice for $(\text{high}, g_{\text{high}})$ &
  (3) community divided, strong signal \\
\bottomrule
\end{tabular}
\end{table}

\subsubsection{Counterbalancing}
\label{subsub:counterbalancing_comp}

The experiment uses the \textbf{same $4 \times 8$ Latin square} as the
production experiment: each of the 4 lists assigns item $i$
($i = 0,\ldots,7$) to condition $(i + l)\bmod 4$, where
$l \in \{0,1,2,3\}$ is the list number.
This guarantees that:
\begin{itemize}
  \item each participant sees every item exactly once;
  \item each participant sees each of the 4 conditions exactly twice
        (two items per condition), ensuring balanced coverage of the
        $2\times2$ design;
  \item across all 4 lists, each item appears equally often in each
        condition.
\end{itemize}

List assignment follows the same Prolific-and-URL strategy as the production
experiment: each Prolific study embeds \texttt{?list=N} (0--3) in the
experiment URL, capping recruitment at $n = 20$ participants per list
for a total of $N = 80$.

In addition to the 8 critical trials, each participant completed 8
\textbf{filler trials} using the same filler items as the production
experiment (train travel, reading light, cooking, coffee, vacation,
hydration, fresh air, noise), each presented with a randomly assigned
marker.
Fillers are excluded from analysis; they serve to dilute the critical
items and reduce demand characteristics.

\subsubsection{Dependent variables}
\label{subsub:dvs_comp}

Two slider ratings are collected after each vignette (0--100 continuous
scale with no visible numbers, labelled at the endpoints as shown):

\begin{enumerate}
  \item \textbf{Inferred speaker goal strength} (operationalises $\widehat{g}$):
  \begin{quote}
    \textit{Wie stark hast du den Eindruck, dass die Sprecherin m\"ochte,
    dass du ihre Empfehlung befolgst?}\\
    (`How strongly do you get the impression that the speaker wants you to
    follow their recommendation?')\\
    Endpoints: 0 = \textit{\"uberhaupt nicht}
    (`not at all') ---
    100 = \textit{sehr stark} (`very strongly').
  \end{quote}

  \item \textbf{Adoption likelihood of $q$} (operationalises $\widehat{A}$):
  \begin{quote}
    \textit{Wie wahrscheinlich ist es, dass du die empfohlene Handlung
    dann auch tats\"achlich umsetzt?}\\
    (`How likely is it that you will actually carry out the recommended
    action?')\\
    Endpoints: 0 = \textit{sehr unwahrscheinlich}
    (`very unlikely') ---
    100 = \textit{sehr wahrscheinlich} (`very likely').
  \end{quote}
\end{enumerate}

\noindent Perceived controversy is deliberately \emph{not} included as a
rating probe: \textsc{pc\textsubscript{prag}} is experimentally
manipulated via the context paragraph, so a controversy rating would
be redundant and risks obscuring the marker's contribution to inferred
goal strength.

Recorded per trial: \texttt{trial\_id}, \texttt{topic},
\texttt{is\_filler}, \texttt{condition\_index}, \texttt{marker},
\texttt{pc\_prag}, \texttt{g\_implied}, \texttt{rating\_g},
\texttt{rating\_adopt}, \texttt{rt} (ms); and once per participant:
\texttt{list\_num}.
The continuous \texttt{pc\_prop\_rating} covariate (topic-level norming
mean) is joined to the trial data during analysis.

\subsection{Procedure}

Participants complete the experiment online via the magpie platform
\citep{kochetkova2024magpie} in a browser.
After reading instructions in German, participants are assigned to one of
4 counterbalancing lists (via URL \texttt{?list=N} parameter assigned by
Prolific; no random fallback in deployment).
They then proceed through 16 trials---8 critical and 8 fillers---in a
fully randomised order, followed by a post-test questionnaire eliciting
native language, age, gender, education level, and optional comments.

Each trial displays:
(a) a context paragraph establishing the communicative situation and the
\textsc{pc\textsubscript{prag}} level;
(b) the target sentence with the best-choice marker already in place in
the Mittelfeld;
(c) the consequent action clause;
(d) the two rating sliders in fixed order ($\widehat{g}$, $\widehat{A}$),
each labelled in German.
The ``Weiter'' (Next) button is disabled until both sliders have been
moved at least once, ensuring a complete response on every trial.

\subsection{Planned Analysis}

\subsubsection{Empirical plots}

For each dependent variable ($\widehat{g}$, $\widehat{A}$):

\begin{itemize}
  \item \textbf{Condition means}: a dot-and-whisker plot shows mean
        ratings per condition (the $2\times2$ of pc\textsubscript{prag}
        $\times$ $g_{\mathrm{implied}}$) with 95\% bootstrapped CIs.
        The expected pattern: $\widehat{g}$ and $\widehat{A}$ should
        be higher in the high-$g_{\mathrm{implied}}$ conditions
        (where the model assigns a stronger marker), and lower in the
        high-pc\textsubscript{prag} conditions (credibility discounting).
  \item \textbf{Marker identity}: because the best-choice marker may
        differ across items within the same condition, a secondary plot
        shows ratings grouped by actual marker identity, allowing
        inspection of whether marker type (rather than just condition)
        drives the effects.
  \item \textbf{pc\textsubscript{prop} covariate}: scatter plots display
        each DV as a function of the continuous norming rating per
        condition, testing whether higher propositional acceptance
        amplifies the $\widehat{g}$ and $\widehat{A}$ signals for
        strong markers.
\end{itemize}

\subsubsection{Bayesian mixed-effects regression}

Two Bayesian mixed-effects regression models are fitted via
\texttt{brms} \citep{burkner2017}, one per DV ($\widehat{g}$ and
$\widehat{A}$):

\begin{align}
y &\sim \mathrm{Beta}(\mu, \phi)
\quad \text{(ratings rescaled to }(0,1)\text{)} \\
\mathrm{logit}(\mu) &= \alpha
+ \beta_{g} \cdot \mathrm{g\_implied\_c}
+ \beta_{\mathrm{prag}} \cdot \mathrm{pc\_prag\_c}
+ \beta_{\mathrm{prop}} \cdot \mathrm{pc\_prop\_c}
+ \beta_{\mathrm{int}} \cdot \mathrm{g\_implied\_c}
  \times \mathrm{pc\_prag\_c}
+ (1 \mid \mathrm{submission\_id})
+ (1 \mid \mathrm{topic}),
\end{align}

where \texttt{g\_implied\_c} is the $g$ level from the corresponding
production condition ($-0.5$ = low, $+0.5$ = high), effectively
coding the marker's position on the production utility scale.
A secondary specification replaces \texttt{g\_implied\_c} with
actualy marker identity as an unordered factor, to examine
non-monotonic patterns and verify that marker rather than
condition label drives effects.

Model comparison via approximate leave-one-out cross-validation (LOO)
selects among additive, interaction, and marker-as-categorical
formulations.

\subsubsection{RSA listener parameter estimation}
\label{subsub:rsa_fit}

The two empirical rating means (per condition cell: 4 conditions) constrain
the listener parameters of the noisy-threshold RSA model.
Specifically, the model predicts:

\begin{align}
\mathbb{E}[\widehat{g} \mid u, \text{context}]
  &= \int g \cdot P(g \mid u, pc_{\mathrm{prag}})\, dg,
\label{eq:pred_g} \\
P(\widehat{A} \mid u, \text{context})
  &= \iiint \sigma(\eta_0 + \eta_g g - \eta_{pc}\, pc
     - \eta_{\mathrm{int}}\, g \cdot pc)\;
     P(pc_{\mathrm{prop}}, pc_{\mathrm{prag}}, g \mid u)\,
     d{pc_{\mathrm{prop}}}\,d{pc_{\mathrm{prag}}}\,dg,
\label{eq:adoption}
\end{align}

\noindent where the posterior $P(pc_{\mathrm{prop}}, pc_{\mathrm{prag}}, g
\mid u)$ is conditioned on the context's \textsc{pc\textsubscript{prag}}
level as a known prior (since it is explicitly given to the listener).

The threshold parameters $(\mu_1,\mu_2, \sigma_1,\sigma_2,
\lambda)$ together with the marker-specific costs $(c_u)$ are estimated
jointly from the production experiment data via the Stan model.
Given those threshold estimates, the listener parameters
$(\eta_0, \eta_g, \eta_{pc}, \eta_{\mathrm{int}})$ are then estimated by
minimising the weighted squared deviation between the model-predicted
means (Eqs.~\ref{eq:pred_g}--\ref{eq:adoption}) and the empirical rating
means from the comprehension experiment, or via a full joint Bayesian
fit over both experiments simultaneously.

\subsection{Model Predictions}

The noisy-threshold RSA model makes the following predictions for
$\widehat{g}$ and $\widehat{A}$, which serve as pre-registered hypotheses.
In all cases, ``strong-marker condition'' refers to conditions 1 and 3
(high $g_{\mathrm{implied}}$, i.e.\ the marker selected by the production
model under a high-$g$ context), and ``weak-marker condition'' to conditions
0 and 2.

\begin{enumerate}
  \item \textbf{Marker encodes inferred goal $\widehat{g}$}: $\widehat{g}$
        ratings should be higher in strong-marker conditions than in
        weak-marker conditions, across both pc\textsubscript{prag} levels.
        This is the core inversion prediction: if the production model
        correctly captures speaker behaviour, listeners should be able to
        decode $g$ from the marker.

  \item \textbf{Marker encodes adoption $\widehat{A}$}: $\widehat{A}$
        ratings should likewise be higher in strong-marker conditions,
        reflecting that higher inferred $g$ translates into greater
        adoption pressure downstream.

  \item \textbf{pc\textsubscript{prag} modulates $\widehat{A}$ but not
        $\widehat{g}$}: because $g$ is orthogonal to pc\textsubscript{prag}
        in the model's utility function, the effect of the marker on
        $\widehat{g}$ should be largely the same in low- and high-
        pc\textsubscript{prag} contexts.
        By contrast, $\widehat{A}$ should be reduced in high-
        pc\textsubscript{prag} contexts via the interaction term
        $\eta_{\mathrm{int}}\,g \cdot pc$: even a strong-marker signal
        for high $g$ is partly cancelled when controversy is high.

  \item \textbf{Credibility discounting}: for strong-marker conditions
        presented in a high-pc\textsubscript{prag} context (condition~3),
        $\widehat{A}$ should be \emph{lower} than for the same marker in
        a low-pc\textsubscript{prag} context (condition~1).
        This is the model's signature prediction distinguishing Bayesian
        inversion from a simpler marker-strength heuristic: a sceptical
        listener who knows controversy is high will discount the adoption
        implication of even the strongest marker.
        The $g_{\mathrm{implied}} \times \textsc{pc\textsubscript{prag}}$
        interaction on $\widehat{A}$ provides the empirical test.
\end{enumerate}

\end{document}
.
%
% Compile stand-alone:
%   pdflatex comprehension_exp.tex
%
% Embed in main.tex:
%   % comprehension_exp.tex
% Stand-alone plan for the comprehension experiment.
% Can be embedded in main.tex via \input{comprehension_exp}.
%
% Compile stand-alone:
%   pdflatex comprehension_exp.tex
%
% Embed in main.tex:
%   \input{comprehension_exp}   % (inside a \section{Experiment} environment)

\documentclass[12pt]{article}
\usepackage{CSPArticleStyle}
\usepackage{csquotes}
\usepackage[margin=1in]{geometry}
\usepackage{booktabs}
\usepackage{parskip}
\usepackage{amsmath}
\usepackage{times}
\usepackage[T1]{fontenc}
\usepackage[utf8]{inputenc}
\usepackage{xcolor}
\usepackage{tabularray}
\usepackage{gb4e}
\noautomath

\title{Comprehension Experiment: Inferring Speaker Goal and Action Adoption
       from Consensus-Marker Use}
\author{}
\date{}

\begin{document}

%% ============================================================
%% When embedding in main.tex, delete from \documentclass to
%% \begin{document} and delete \end{document} at the bottom.
%% ============================================================

\section{Experiment: Consensus-Marker Comprehension}

\subsection{Overview}

The comprehension experiment is the listener-side counterpart of the
production experiment and tests whether listeners can invert the
speaker's marker choice to recover the latent context that motivated it.
Where the production experiment gave participants a context (fixing
\textsc{pc\textsubscript{prag}} and $g$) and asked which marker they
would use, the comprehension experiment gives participants a context
(fixing \textsc{pc\textsubscript{prag}}) together with a marker already
embedded in the sentence and asks what goal strength the speaker likely
had and how likely the listener is to adopt $q$.

The specific marker presented in each vignette is not chosen arbitrarily
from the three candidate markers. Our preregistered plan was to assign,
for each item$\times$condition cell, the \emph{Bayesian best-choice marker}:
the marker with the highest posterior probability under the cost-augmented
noisy-threshold RSA model fitted to the production data, given the context's
\textsc{pc\textsubscript{prag}} level and a specified $g$ level. This was
intended to mirror the production design closely and to ensure that the
comprehension stimuli reflect model-based marker use rather than arbitrary
marker$\times$item pairings.

Applying this strict modal-assignment rule to the fitted production model,
however, yielded a highly imbalanced stimulus set: most critical cells were
assigned the strongest marker, with very little variation across marker
categories. In addition, many of these modal decisions were supported only by
very small probability margins over the next-best alternative. Because such a
stimulus set would make the comprehension experiment substantially less
informative about marker-based listener inferences, we introduced a documented,
exploratory balancing step for stimulus construction. Starting from the same
posterior predictive probabilities, we selected a near-balanced assignment of
the three markers across the 32 critical cells while maximizing total posterior
support under the fitted production model. The strict modal assignment is
retained as the preregistered baseline, but the balanced assignment is used for
the actual comprehension materials.

\textsc{pc\textsubscript{prag}} is retained as an explicit context
manipulation (low vs.\ high), so that listeners can be expected to
have a well-defined prior on pragmatic controversy when interpreting
the marker. \textsc{pc\textsubscript{prop}} is not manipulated within
the experiment; it enters as a continuous topic-level covariate from
the norming study.

After reading each vignette, participants provide two slider ratings
operationalising the key listener-side quantities:
\begin{enumerate}
  \item \textbf{Inferred speaker goal strength} ($\widehat{g}$):
        how strongly does the speaker seem to want the listener to
        accept and act on $q$?
  \item \textbf{Adoption likelihood} ($\widehat{A}$):
        how likely is the listener to follow the recommended action $q$?
\end{enumerate}

\noindent Perceived controversy is \emph{not} elicited as a rating:
because \textsc{pc\textsubscript{prag}} is already manipulated as a
context factor, asking participants to rate controversy would be
redundant with respect to the manipulation check and circular with
respect to the research question. The theoretically central inference
to test is whether listeners can decode the speaker's \emph{goal} from
marker use, and whether that inference translates into adoption
behaviour.

The two slider measures, together with the \textsc{pc\textsubscript{prag}}
context and the \textsc{pc\textsubscript{prop}} covariate, provide the
data needed to fit the model's listener parameters
$(\eta_0, \eta_g, \eta_{pc}, \eta_{\mathrm{int}})$.

\subsection{Participants}

Participants will be recruited via Prolific and will be required to be
native speakers of German (language-at-home filter).
Data collection is planned after completion of the production experiment
(model fitting required to determine the best-choice markers).
A target of $N = 80$ participants is planned ($n = 20$ per list across
4 counterbalancing lists), each completing 16 trials (8 critical + 8
filler), matching the structure of the production experiment.

\subsection{Materials}

\subsubsection{Markers}

Each critical vignette is presented with one specific consensus marker
already embedded in the Mittelfeld of the sentence frame (see
example~(\ref{ex:frame_comp})):

\begin{exe}
  \ex \label{ex:frame_comp}
  Digitale Kompetenzen sind \textbf{[MARKER]}
  in der heutigen Arbeitswelt unverzichtbar.\\
  `Digital skills are \textbf{[MARKER]} indispensable in today's
  working world.'
\end{exe}

Unlike the production experiment, participants do \emph{not} choose the
marker; they read the sentence with the marker already in place and respond
to the two rating probes.

Critically, the marker assigned to each condition is not drawn arbitrarily
from the three candidates. Under the preregistered plan, for each of the four
production conditions (pc\textsubscript{prag} $\times$ $g$), the marker used in
comprehension was to be the \textbf{Bayesian best-choice marker}: the marker
$u^*$ with the highest posterior probability under the fitted
cost-augmented noisy-threshold RSA production model,
\[
u^* = \arg\max_{u} \; P(u \mid pc_{\mathrm{prag}}, g,
                           pc_{\mathrm{prop}\_c}).
\]
Because $pc_{\mathrm{prop}}$ varies continuously across items,
$u^*$ may differ per item even within the same
(pc\textsubscript{prag}, $g$) cell; the final stimulus list is therefore
constructed item-by-item once the production model has been fitted.
In practice, the fitted production model yielded a nearly degenerate modal map,
with most cells assigned to \textit{bekanntlich}. We therefore retained the
model-based logic of the procedure but replaced strict cellwise modal selection
with an uncertainty-aware balancing step: among the marker assignments with
high posterior support, we selected a near-balanced configuration across the 32
critical cells. This preserves the production--comprehension link while also
ensuring that the comprehension experiment contains sufficient marker variation
to test listener-side inferences.

All three markers remain candidates; the mapping from conditions to
markers is an empirical output of the production analysis, not a design
pre-commitment. In the final comprehension materials, we use the
\textbf{balanced model-based assignment}, rather than the strict modal
assignment, in order to maintain informative variation across marker types.

\subsubsection{Topic Items}

The same eight topic domains used in the norming study and the production
experiment are used here (Table~\ref{tab:items_comp}).
The 4-condition $\times$ 8-item Latin square (Section~\ref{subsub:counterbalancing_comp})
requires no additional items.

\begin{table}[ht]
\centering
\caption{Topic items: sentence frame for $p$ (with marker position [M])
         and consequent $q$.}
\label{tab:items_comp}
\small
\begin{tabular}{clll}
\toprule
\textbf{\#} & \textbf{Topic} & \textbf{Sentence frame} & \textbf{Consequent $q$} \\
\midrule
1 & Klimawandel &
  \textit{Der Klimawandel wird [M]} &
  auf Flugreisen verzichten \\
  & & \textit{durch menschliche Aktivit\"aten verursacht.} & \\[3pt]
2 & Ern\"ahrung &
  \textit{Eine \"uberwiegend pflanzliche Ern\"ahrung} &
  mehr pflanzliche Lebens- \\
  & & \textit{f\"ordert [M] die Gesundheit.} &
  mittel einbauen \\[3pt]
3 & Stadtverkehr &
  \textit{Ein ausgebautes Nahverkehrsnetz} &
  h\"aufiger auf \\
  & & \textit{entlastet [M] den Stadtverkehr} &
  \"offentliche Verkehrs- \\
  & & \textit{deutlich.} &
  mittel umsteigen \\[3pt]
4 & Digitalisierung &
  \textit{Digitale Kompetenzen sind [M] in} &
  regelm\"a\ss{}ig Weiter- \\
  & & \textit{der heutigen Arbeitswelt unverzichtbar.} &
  bildungsangebote nutzen \\[3pt]
5 & Schlaf &
  \textit{Ausreichend Schlaf ist [M] entscheidend} &
  auf feste Schlafzeiten \\
  & & \textit{f\"ur die kognitive Leistungsf\"ahigkeit.} &
  achten \\[3pt]
6 & Plastik &
  \textit{Einwegplastik schadet [M] den} &
  konsequent auf Einweg- \\
  & & \textit{Meeres\"okosystemen erheblich.} &
  plastik verzichten \\[3pt]
7 & Sport &
  \textit{Regelm\"a\ss{}ige k\"orperliche Bewegung} &
  regelm\"a\ss{}ig Sport in \\
  & & \textit{senkt [M] das Risiko f\"ur} &
  den Alltag integrieren \\
  & & \textit{Herz-Kreislauf-Erkrankungen deutlich.} & \\[3pt]
8 & Lokalkauf &
  \textit{Das Kaufen bei lokalen H\"andlern} &
  \"ofter bei lokalen \\
  & & \textit{st\"arkt [M] die regionale Wirtschaft} &
  Gesch\"aften einkaufen \\
  & & \textit{nachhaltig.} & \\
\bottomrule
\end{tabular}
\end{table}

\noindent [M] denotes the position of the pre-filled best-choice marker.

\subsubsection{Vignettes}

\paragraph{Structure.}
Each trial presents a vignette consisting of three parts:
(a) a context paragraph establishing the communicative situation and
fixing the \textsc{pc\textsubscript{prag}} level;
(b) the target sentence with the best-choice marker already embedded
in the Mittelfeld;
(c) the consequent action clause.
The two rating probes follow immediately below the vignette.

\paragraph{Propositional perceived controversy (\textsc{pc\textsubscript{prop}}).}
As in the production experiment, \textsc{pc\textsubscript{prop}} is not
manipulated within the comprehension experiment.
Topic-level propositional acceptance ratings from the norming study
(0--100 scale; e.g., \textit{Digitalisierung}: 85/100,
\textit{Lokalkauf}: 58/100) serve as a continuous, mean-centred
covariate \textsc{pc\textsubscript{prop}}$_c$ in analysis.
Context paragraphs present $p$ without editorial evaluation; variation in
listeners' individual assessments of propositional controversy is
captured by the norming-derived covariate.

\paragraph{Pragmatic perceived controversy (\textsc{pc\textsubscript{prag}}).}
\textsc{pc\textsubscript{prag}} is manipulated between item instances
using the same two context versions as in the production experiment:
the speaker's social circle either \emph{largely agrees} (low) or
\emph{holds divergent views} (high) about $p$.
For example, for the \textit{Digitalisierung} item:

\begin{itemize}
  \item $\textsc{pc}_\text{prag} = \text{low}$: \textit{Deine Kollegin spricht
        mit dir \"uber Berufsausbildung und Weiterbildung.
        In eurem Team teilen die meisten die Ansicht, dass digitale Kompetenzen
        in der heutigen Arbeitswelt unverzichtbar sind.}
  \item $\textsc{pc}_\text{prag} = \text{high}$: \textit{Deine Kollegin spricht
        mit dir \"uber Berufsausbildung und Weiterbildung.
        In eurem Team gehen die Meinungen dazu stark auseinander.}
\end{itemize}

The context is framed from the listener's perspective as a direct addressee,
mirroring the production study where participants adopt the speaker's
perspective.

\paragraph{Marker manipulation.}
The dependent measures of the production experiment (which marker is chosen)
become the independent variable here.
One of the three consensus markers is inserted into the Mittelfeld of the
target sentence.
A complete sample item for the \textit{Digitalisierung} topic
($\textsc{pc}_\text{prag} = \text{high}$, marker = \textit{bekanntlich};
norming $\textsc{pc}_\text{prop}$ rating: 85/100) is given in
(\ref{ex:comp_sample}):

\begin{exe}
  \ex \label{ex:comp_sample}
  \textit{Deine Kollegin spricht mit dir \"uber Berufsausbildung und
  Weiterbildung.
  In eurem Team gehen die Meinungen dazu stark auseinander.}

  \textit{Sie sagt zu dir:} Digitale Kompetenzen sind \textbf{bekanntlich}
  in der heutigen Arbeitswelt unverzichtbar.
  Deshalb solltest du regelm\"a\ss{}ig digitale Weiterbildungsangebote nutzen.

  \noindent\textit{[Rating probes follow.]}
\end{exe}

\subsection{Design}

\subsubsection{Factorial structure}

The comprehension experiment mirrors the production experiment's
$\mathbf{2}$~\textbf{(pc\textsubscript{prag}: low vs.\ high)}
$\mathbf{\times}$~$\mathbf{2}$~\textbf{($g_{\mathrm{implied}}$: low
vs.\ high) design}, yielding \textbf{4 conditions}.
The $g_{\mathrm{implied}}$ dimension is not manipulated via an explicit
goal instruction (as in production) but is instead conveyed entirely
through the choice of marker: for each condition, the marker presented
is the best-choice marker for the corresponding
(\textsc{pc\textsubscript{prag}},~$g$) cell from the production model.
\textsc{pc\textsubscript{prop}} enters as a continuous norming covariate.
With 8 topic items and a Latin square rotation across 4 lists, each
participant sees exactly \textbf{8 critical trials}---one item per
condition, each condition appearing twice---plus 8 filler trials,
for \textbf{16 trials} in total (Table~\ref{tab:design_comp}).

\begin{table}[ht]
\centering
\caption{%
  The $2\times2$ design (4 conditions), mirroring the production experiment.
  The marker presented in each condition is the Bayesian best-choice
  marker from the fitted production model; specific markers are
  determined after production data collection.}
\label{tab:design_comp}
\small
\begin{tabular}{ccll}
\toprule
\textbf{pc\textsubscript{prag}} &
\textbf{$g_{\mathrm{implied}}$} &
\textbf{Marker (TBD from production)} &
\textbf{Condition label} \\
\midrule
low  & low  & best-choice for $(\text{low}, g_{\text{low}})$  &
  (0) community agrees, weak signal \\
low  & high & best-choice for $(\text{low}, g_{\text{high}})$ &
  (1) community agrees, strong signal \\
high & low  & best-choice for $(\text{high}, g_{\text{low}})$ &
  (2) community divided, weak signal \\
high & high & best-choice for $(\text{high}, g_{\text{high}})$ &
  (3) community divided, strong signal \\
\bottomrule
\end{tabular}
\end{table}

\subsubsection{Counterbalancing}
\label{subsub:counterbalancing_comp}

The experiment uses the \textbf{same $4 \times 8$ Latin square} as the
production experiment: each of the 4 lists assigns item $i$
($i = 0,\ldots,7$) to condition $(i + l)\bmod 4$, where
$l \in \{0,1,2,3\}$ is the list number.
This guarantees that:
\begin{itemize}
  \item each participant sees every item exactly once;
  \item each participant sees each of the 4 conditions exactly twice
        (two items per condition), ensuring balanced coverage of the
        $2\times2$ design;
  \item across all 4 lists, each item appears equally often in each
        condition.
\end{itemize}

List assignment follows the same Prolific-and-URL strategy as the production
experiment: each Prolific study embeds \texttt{?list=N} (0--3) in the
experiment URL, capping recruitment at $n = 20$ participants per list
for a total of $N = 80$.

In addition to the 8 critical trials, each participant completed 8
\textbf{filler trials} using the same filler items as the production
experiment (train travel, reading light, cooking, coffee, vacation,
hydration, fresh air, noise), each presented with a randomly assigned
marker.
Fillers are excluded from analysis; they serve to dilute the critical
items and reduce demand characteristics.

\subsubsection{Dependent variables}
\label{subsub:dvs_comp}

Two slider ratings are collected after each vignette (0--100 continuous
scale with no visible numbers, labelled at the endpoints as shown):

\begin{enumerate}
  \item \textbf{Inferred speaker goal strength} (operationalises $\widehat{g}$):
  \begin{quote}
    \textit{Wie stark hast du den Eindruck, dass die Sprecherin m\"ochte,
    dass du ihre Empfehlung befolgst?}\\
    (`How strongly do you get the impression that the speaker wants you to
    follow their recommendation?')\\
    Endpoints: 0 = \textit{\"uberhaupt nicht}
    (`not at all') ---
    100 = \textit{sehr stark} (`very strongly').
  \end{quote}

  \item \textbf{Adoption likelihood of $q$} (operationalises $\widehat{A}$):
  \begin{quote}
    \textit{Wie wahrscheinlich ist es, dass du die empfohlene Handlung
    dann auch tats\"achlich umsetzt?}\\
    (`How likely is it that you will actually carry out the recommended
    action?')\\
    Endpoints: 0 = \textit{sehr unwahrscheinlich}
    (`very unlikely') ---
    100 = \textit{sehr wahrscheinlich} (`very likely').
  \end{quote}
\end{enumerate}

\noindent Perceived controversy is deliberately \emph{not} included as a
rating probe: \textsc{pc\textsubscript{prag}} is experimentally
manipulated via the context paragraph, so a controversy rating would
be redundant and risks obscuring the marker's contribution to inferred
goal strength.

Recorded per trial: \texttt{trial\_id}, \texttt{topic},
\texttt{is\_filler}, \texttt{condition\_index}, \texttt{marker},
\texttt{pc\_prag}, \texttt{g\_implied}, \texttt{rating\_g},
\texttt{rating\_adopt}, \texttt{rt} (ms); and once per participant:
\texttt{list\_num}.
The continuous \texttt{pc\_prop\_rating} covariate (topic-level norming
mean) is joined to the trial data during analysis.

\subsection{Procedure}

Participants complete the experiment online via the magpie platform
\citep{kochetkova2024magpie} in a browser.
After reading instructions in German, participants are assigned to one of
4 counterbalancing lists (via URL \texttt{?list=N} parameter assigned by
Prolific; no random fallback in deployment).
They then proceed through 16 trials---8 critical and 8 fillers---in a
fully randomised order, followed by a post-test questionnaire eliciting
native language, age, gender, education level, and optional comments.

Each trial displays:
(a) a context paragraph establishing the communicative situation and the
\textsc{pc\textsubscript{prag}} level;
(b) the target sentence with the best-choice marker already in place in
the Mittelfeld;
(c) the consequent action clause;
(d) the two rating sliders in fixed order ($\widehat{g}$, $\widehat{A}$),
each labelled in German.
The ``Weiter'' (Next) button is disabled until both sliders have been
moved at least once, ensuring a complete response on every trial.

\subsection{Planned Analysis}

\subsubsection{Empirical plots}

For each dependent variable ($\widehat{g}$, $\widehat{A}$):

\begin{itemize}
  \item \textbf{Condition means}: a dot-and-whisker plot shows mean
        ratings per condition (the $2\times2$ of pc\textsubscript{prag}
        $\times$ $g_{\mathrm{implied}}$) with 95\% bootstrapped CIs.
        The expected pattern: $\widehat{g}$ and $\widehat{A}$ should
        be higher in the high-$g_{\mathrm{implied}}$ conditions
        (where the model assigns a stronger marker), and lower in the
        high-pc\textsubscript{prag} conditions (credibility discounting).
  \item \textbf{Marker identity}: because the best-choice marker may
        differ across items within the same condition, a secondary plot
        shows ratings grouped by actual marker identity, allowing
        inspection of whether marker type (rather than just condition)
        drives the effects.
  \item \textbf{pc\textsubscript{prop} covariate}: scatter plots display
        each DV as a function of the continuous norming rating per
        condition, testing whether higher propositional acceptance
        amplifies the $\widehat{g}$ and $\widehat{A}$ signals for
        strong markers.
\end{itemize}

\subsubsection{Bayesian mixed-effects regression}

Two Bayesian mixed-effects regression models are fitted via
\texttt{brms} \citep{burkner2017}, one per DV ($\widehat{g}$ and
$\widehat{A}$):

\begin{align}
y &\sim \mathrm{Beta}(\mu, \phi)
\quad \text{(ratings rescaled to }(0,1)\text{)} \\
\mathrm{logit}(\mu) &= \alpha
+ \beta_{g} \cdot \mathrm{g\_implied\_c}
+ \beta_{\mathrm{prag}} \cdot \mathrm{pc\_prag\_c}
+ \beta_{\mathrm{prop}} \cdot \mathrm{pc\_prop\_c}
+ \beta_{\mathrm{int}} \cdot \mathrm{g\_implied\_c}
  \times \mathrm{pc\_prag\_c}
+ (1 \mid \mathrm{submission\_id})
+ (1 \mid \mathrm{topic}),
\end{align}

where \texttt{g\_implied\_c} is the $g$ level from the corresponding
production condition ($-0.5$ = low, $+0.5$ = high), effectively
coding the marker's position on the production utility scale.
A secondary specification replaces \texttt{g\_implied\_c} with
actualy marker identity as an unordered factor, to examine
non-monotonic patterns and verify that marker rather than
condition label drives effects.

Model comparison via approximate leave-one-out cross-validation (LOO)
selects among additive, interaction, and marker-as-categorical
formulations.

\subsubsection{RSA listener parameter estimation}
\label{subsub:rsa_fit}

The two empirical rating means (per condition cell: 4 conditions) constrain
the listener parameters of the noisy-threshold RSA model.
Specifically, the model predicts:

\begin{align}
\mathbb{E}[\widehat{g} \mid u, \text{context}]
  &= \int g \cdot P(g \mid u, pc_{\mathrm{prag}})\, dg,
\label{eq:pred_g} \\
P(\widehat{A} \mid u, \text{context})
  &= \iiint \sigma(\eta_0 + \eta_g g - \eta_{pc}\, pc
     - \eta_{\mathrm{int}}\, g \cdot pc)\;
     P(pc_{\mathrm{prop}}, pc_{\mathrm{prag}}, g \mid u)\,
     d{pc_{\mathrm{prop}}}\,d{pc_{\mathrm{prag}}}\,dg,
\label{eq:adoption}
\end{align}

\noindent where the posterior $P(pc_{\mathrm{prop}}, pc_{\mathrm{prag}}, g
\mid u)$ is conditioned on the context's \textsc{pc\textsubscript{prag}}
level as a known prior (since it is explicitly given to the listener).

The threshold parameters $(\mu_1,\mu_2, \sigma_1,\sigma_2,
\lambda)$ together with the marker-specific costs $(c_u)$ are estimated
jointly from the production experiment data via the Stan model.
Given those threshold estimates, the listener parameters
$(\eta_0, \eta_g, \eta_{pc}, \eta_{\mathrm{int}})$ are then estimated by
minimising the weighted squared deviation between the model-predicted
means (Eqs.~\ref{eq:pred_g}--\ref{eq:adoption}) and the empirical rating
means from the comprehension experiment, or via a full joint Bayesian
fit over both experiments simultaneously.

\subsection{Model Predictions}

The noisy-threshold RSA model makes the following predictions for
$\widehat{g}$ and $\widehat{A}$, which serve as pre-registered hypotheses.
In all cases, ``strong-marker condition'' refers to conditions 1 and 3
(high $g_{\mathrm{implied}}$, i.e.\ the marker selected by the production
model under a high-$g$ context), and ``weak-marker condition'' to conditions
0 and 2.

\begin{enumerate}
  \item \textbf{Marker encodes inferred goal $\widehat{g}$}: $\widehat{g}$
        ratings should be higher in strong-marker conditions than in
        weak-marker conditions, across both pc\textsubscript{prag} levels.
        This is the core inversion prediction: if the production model
        correctly captures speaker behaviour, listeners should be able to
        decode $g$ from the marker.

  \item \textbf{Marker encodes adoption $\widehat{A}$}: $\widehat{A}$
        ratings should likewise be higher in strong-marker conditions,
        reflecting that higher inferred $g$ translates into greater
        adoption pressure downstream.

  \item \textbf{pc\textsubscript{prag} modulates $\widehat{A}$ but not
        $\widehat{g}$}: because $g$ is orthogonal to pc\textsubscript{prag}
        in the model's utility function, the effect of the marker on
        $\widehat{g}$ should be largely the same in low- and high-
        pc\textsubscript{prag} contexts.
        By contrast, $\widehat{A}$ should be reduced in high-
        pc\textsubscript{prag} contexts via the interaction term
        $\eta_{\mathrm{int}}\,g \cdot pc$: even a strong-marker signal
        for high $g$ is partly cancelled when controversy is high.

  \item \textbf{Credibility discounting}: for strong-marker conditions
        presented in a high-pc\textsubscript{prag} context (condition~3),
        $\widehat{A}$ should be \emph{lower} than for the same marker in
        a low-pc\textsubscript{prag} context (condition~1).
        This is the model's signature prediction distinguishing Bayesian
        inversion from a simpler marker-strength heuristic: a sceptical
        listener who knows controversy is high will discount the adoption
        implication of even the strongest marker.
        The $g_{\mathrm{implied}} \times \textsc{pc\textsubscript{prag}}$
        interaction on $\widehat{A}$ provides the empirical test.
\end{enumerate}

\end{document}
   % (inside a \section{Experiment} environment)

\documentclass[12pt]{article}
\usepackage{CSPArticleStyle}
\usepackage{csquotes}
\usepackage[margin=1in]{geometry}
\usepackage{booktabs}
\usepackage{parskip}
\usepackage{amsmath}
\usepackage{times}
\usepackage[T1]{fontenc}
\usepackage[utf8]{inputenc}
\usepackage{xcolor}
\usepackage{tabularray}
\usepackage{gb4e}
\noautomath

\title{Comprehension Experiment: Inferring Speaker Goal and Action Adoption
       from Consensus-Marker Use}
\author{}
\date{}

\begin{document}

%% ============================================================
%% When embedding in main.tex, delete from \documentclass to
%% \begin{document} and delete \end{document} at the bottom.
%% ============================================================

\section{Experiment: Consensus-Marker Comprehension}

\subsection{Overview}

The comprehension experiment is the listener-side counterpart of the
production experiment and tests whether listeners can invert the
speaker's marker choice to recover the latent context that motivated it.
Where the production experiment gave participants a context (fixing
\textsc{pc\textsubscript{prag}} and $g$) and asked which marker they
would use, the comprehension experiment gives participants a context
(fixing \textsc{pc\textsubscript{prag}}) together with a marker already
embedded in the sentence and asks what goal strength the speaker likely
had and how likely the listener is to adopt $q$.

The specific marker presented in each vignette is not chosen arbitrarily
from the three candidate markers. Our preregistered plan was to assign,
for each item$\times$condition cell, the \emph{Bayesian best-choice marker}:
the marker with the highest posterior probability under the cost-augmented
noisy-threshold RSA model fitted to the production data, given the context's
\textsc{pc\textsubscript{prag}} level and a specified $g$ level. This was
intended to mirror the production design closely and to ensure that the
comprehension stimuli reflect model-based marker use rather than arbitrary
marker$\times$item pairings.

Applying this strict modal-assignment rule to the fitted production model,
however, yielded a highly imbalanced stimulus set: most critical cells were
assigned the strongest marker, with very little variation across marker
categories. In addition, many of these modal decisions were supported only by
very small probability margins over the next-best alternative. Because such a
stimulus set would make the comprehension experiment substantially less
informative about marker-based listener inferences, we introduced a documented,
exploratory balancing step for stimulus construction. Starting from the same
posterior predictive probabilities, we selected a near-balanced assignment of
the three markers across the 32 critical cells while maximizing total posterior
support under the fitted production model. The strict modal assignment is
retained as the preregistered baseline, but the balanced assignment is used for
the actual comprehension materials.

\textsc{pc\textsubscript{prag}} is retained as an explicit context
manipulation (low vs.\ high), so that listeners can be expected to
have a well-defined prior on pragmatic controversy when interpreting
the marker. \textsc{pc\textsubscript{prop}} is not manipulated within
the experiment; it enters as a continuous topic-level covariate from
the norming study.

After reading each vignette, participants provide two slider ratings
operationalising the key listener-side quantities:
\begin{enumerate}
  \item \textbf{Inferred speaker goal strength} ($\widehat{g}$):
        how strongly does the speaker seem to want the listener to
        accept and act on $q$?
  \item \textbf{Adoption likelihood} ($\widehat{A}$):
        how likely is the listener to follow the recommended action $q$?
\end{enumerate}

\noindent Perceived controversy is \emph{not} elicited as a rating:
because \textsc{pc\textsubscript{prag}} is already manipulated as a
context factor, asking participants to rate controversy would be
redundant with respect to the manipulation check and circular with
respect to the research question. The theoretically central inference
to test is whether listeners can decode the speaker's \emph{goal} from
marker use, and whether that inference translates into adoption
behaviour.

The two slider measures, together with the \textsc{pc\textsubscript{prag}}
context and the \textsc{pc\textsubscript{prop}} covariate, provide the
data needed to fit the model's listener parameters
$(\eta_0, \eta_g, \eta_{pc}, \eta_{\mathrm{int}})$.

\subsection{Participants}

Participants will be recruited via Prolific and will be required to be
native speakers of German (language-at-home filter).
Data collection is planned after completion of the production experiment
(model fitting required to determine the best-choice markers).
A target of $N = 80$ participants is planned ($n = 20$ per list across
4 counterbalancing lists), each completing 16 trials (8 critical + 8
filler), matching the structure of the production experiment.

\subsection{Materials}

\subsubsection{Markers}

Each critical vignette is presented with one specific consensus marker
already embedded in the Mittelfeld of the sentence frame (see
example~(\ref{ex:frame_comp})):

\begin{exe}
  \ex \label{ex:frame_comp}
  Digitale Kompetenzen sind \textbf{[MARKER]}
  in der heutigen Arbeitswelt unverzichtbar.\\
  `Digital skills are \textbf{[MARKER]} indispensable in today's
  working world.'
\end{exe}

Unlike the production experiment, participants do \emph{not} choose the
marker; they read the sentence with the marker already in place and respond
to the two rating probes.

Critically, the marker assigned to each condition is not drawn arbitrarily
from the three candidates. Under the preregistered plan, for each of the four
production conditions (pc\textsubscript{prag} $\times$ $g$), the marker used in
comprehension was to be the \textbf{Bayesian best-choice marker}: the marker
$u^*$ with the highest posterior probability under the fitted
cost-augmented noisy-threshold RSA production model,
\[
u^* = \arg\max_{u} \; P(u \mid pc_{\mathrm{prag}}, g,
                           pc_{\mathrm{prop}\_c}).
\]
Because $pc_{\mathrm{prop}}$ varies continuously across items,
$u^*$ may differ per item even within the same
(pc\textsubscript{prag}, $g$) cell; the final stimulus list is therefore
constructed item-by-item once the production model has been fitted.
In practice, the fitted production model yielded a nearly degenerate modal map,
with most cells assigned to \textit{bekanntlich}. We therefore retained the
model-based logic of the procedure but replaced strict cellwise modal selection
with an uncertainty-aware balancing step: among the marker assignments with
high posterior support, we selected a near-balanced configuration across the 32
critical cells. This preserves the production--comprehension link while also
ensuring that the comprehension experiment contains sufficient marker variation
to test listener-side inferences.

All three markers remain candidates; the mapping from conditions to
markers is an empirical output of the production analysis, not a design
pre-commitment. In the final comprehension materials, we use the
\textbf{balanced model-based assignment}, rather than the strict modal
assignment, in order to maintain informative variation across marker types.

\subsubsection{Topic Items}

The same eight topic domains used in the norming study and the production
experiment are used here (Table~\ref{tab:items_comp}).
The 4-condition $\times$ 8-item Latin square (Section~\ref{subsub:counterbalancing_comp})
requires no additional items.

\begin{table}[ht]
\centering
\caption{Topic items: sentence frame for $p$ (with marker position [M])
         and consequent $q$.}
\label{tab:items_comp}
\small
\begin{tabular}{clll}
\toprule
\textbf{\#} & \textbf{Topic} & \textbf{Sentence frame} & \textbf{Consequent $q$} \\
\midrule
1 & Klimawandel &
  \textit{Der Klimawandel wird [M]} &
  auf Flugreisen verzichten \\
  & & \textit{durch menschliche Aktivit\"aten verursacht.} & \\[3pt]
2 & Ern\"ahrung &
  \textit{Eine \"uberwiegend pflanzliche Ern\"ahrung} &
  mehr pflanzliche Lebens- \\
  & & \textit{f\"ordert [M] die Gesundheit.} &
  mittel einbauen \\[3pt]
3 & Stadtverkehr &
  \textit{Ein ausgebautes Nahverkehrsnetz} &
  h\"aufiger auf \\
  & & \textit{entlastet [M] den Stadtverkehr} &
  \"offentliche Verkehrs- \\
  & & \textit{deutlich.} &
  mittel umsteigen \\[3pt]
4 & Digitalisierung &
  \textit{Digitale Kompetenzen sind [M] in} &
  regelm\"a\ss{}ig Weiter- \\
  & & \textit{der heutigen Arbeitswelt unverzichtbar.} &
  bildungsangebote nutzen \\[3pt]
5 & Schlaf &
  \textit{Ausreichend Schlaf ist [M] entscheidend} &
  auf feste Schlafzeiten \\
  & & \textit{f\"ur die kognitive Leistungsf\"ahigkeit.} &
  achten \\[3pt]
6 & Plastik &
  \textit{Einwegplastik schadet [M] den} &
  konsequent auf Einweg- \\
  & & \textit{Meeres\"okosystemen erheblich.} &
  plastik verzichten \\[3pt]
7 & Sport &
  \textit{Regelm\"a\ss{}ige k\"orperliche Bewegung} &
  regelm\"a\ss{}ig Sport in \\
  & & \textit{senkt [M] das Risiko f\"ur} &
  den Alltag integrieren \\
  & & \textit{Herz-Kreislauf-Erkrankungen deutlich.} & \\[3pt]
8 & Lokalkauf &
  \textit{Das Kaufen bei lokalen H\"andlern} &
  \"ofter bei lokalen \\
  & & \textit{st\"arkt [M] die regionale Wirtschaft} &
  Gesch\"aften einkaufen \\
  & & \textit{nachhaltig.} & \\
\bottomrule
\end{tabular}
\end{table}

\noindent [M] denotes the position of the pre-filled best-choice marker.

\subsubsection{Vignettes}

\paragraph{Structure.}
Each trial presents a vignette consisting of three parts:
(a) a context paragraph establishing the communicative situation and
fixing the \textsc{pc\textsubscript{prag}} level;
(b) the target sentence with the best-choice marker already embedded
in the Mittelfeld;
(c) the consequent action clause.
The two rating probes follow immediately below the vignette.

\paragraph{Propositional perceived controversy (\textsc{pc\textsubscript{prop}}).}
As in the production experiment, \textsc{pc\textsubscript{prop}} is not
manipulated within the comprehension experiment.
Topic-level propositional acceptance ratings from the norming study
(0--100 scale; e.g., \textit{Digitalisierung}: 85/100,
\textit{Lokalkauf}: 58/100) serve as a continuous, mean-centred
covariate \textsc{pc\textsubscript{prop}}$_c$ in analysis.
Context paragraphs present $p$ without editorial evaluation; variation in
listeners' individual assessments of propositional controversy is
captured by the norming-derived covariate.

\paragraph{Pragmatic perceived controversy (\textsc{pc\textsubscript{prag}}).}
\textsc{pc\textsubscript{prag}} is manipulated between item instances
using the same two context versions as in the production experiment:
the speaker's social circle either \emph{largely agrees} (low) or
\emph{holds divergent views} (high) about $p$.
For example, for the \textit{Digitalisierung} item:

\begin{itemize}
  \item $\textsc{pc}_\text{prag} = \text{low}$: \textit{Deine Kollegin spricht
        mit dir \"uber Berufsausbildung und Weiterbildung.
        In eurem Team teilen die meisten die Ansicht, dass digitale Kompetenzen
        in der heutigen Arbeitswelt unverzichtbar sind.}
  \item $\textsc{pc}_\text{prag} = \text{high}$: \textit{Deine Kollegin spricht
        mit dir \"uber Berufsausbildung und Weiterbildung.
        In eurem Team gehen die Meinungen dazu stark auseinander.}
\end{itemize}

The context is framed from the listener's perspective as a direct addressee,
mirroring the production study where participants adopt the speaker's
perspective.

\paragraph{Marker manipulation.}
The dependent measures of the production experiment (which marker is chosen)
become the independent variable here.
One of the three consensus markers is inserted into the Mittelfeld of the
target sentence.
A complete sample item for the \textit{Digitalisierung} topic
($\textsc{pc}_\text{prag} = \text{high}$, marker = \textit{bekanntlich};
norming $\textsc{pc}_\text{prop}$ rating: 85/100) is given in
(\ref{ex:comp_sample}):

\begin{exe}
  \ex \label{ex:comp_sample}
  \textit{Deine Kollegin spricht mit dir \"uber Berufsausbildung und
  Weiterbildung.
  In eurem Team gehen die Meinungen dazu stark auseinander.}

  \textit{Sie sagt zu dir:} Digitale Kompetenzen sind \textbf{bekanntlich}
  in der heutigen Arbeitswelt unverzichtbar.
  Deshalb solltest du regelm\"a\ss{}ig digitale Weiterbildungsangebote nutzen.

  \noindent\textit{[Rating probes follow.]}
\end{exe}

\subsection{Design}

\subsubsection{Factorial structure}

The comprehension experiment mirrors the production experiment's
$\mathbf{2}$~\textbf{(pc\textsubscript{prag}: low vs.\ high)}
$\mathbf{\times}$~$\mathbf{2}$~\textbf{($g_{\mathrm{implied}}$: low
vs.\ high) design}, yielding \textbf{4 conditions}.
The $g_{\mathrm{implied}}$ dimension is not manipulated via an explicit
goal instruction (as in production) but is instead conveyed entirely
through the choice of marker: for each condition, the marker presented
is the best-choice marker for the corresponding
(\textsc{pc\textsubscript{prag}},~$g$) cell from the production model.
\textsc{pc\textsubscript{prop}} enters as a continuous norming covariate.
With 8 topic items and a Latin square rotation across 4 lists, each
participant sees exactly \textbf{8 critical trials}---one item per
condition, each condition appearing twice---plus 8 filler trials,
for \textbf{16 trials} in total (Table~\ref{tab:design_comp}).

\begin{table}[ht]
\centering
\caption{%
  The $2\times2$ design (4 conditions), mirroring the production experiment.
  The marker presented in each condition is the Bayesian best-choice
  marker from the fitted production model; specific markers are
  determined after production data collection.}
\label{tab:design_comp}
\small
\begin{tabular}{ccll}
\toprule
\textbf{pc\textsubscript{prag}} &
\textbf{$g_{\mathrm{implied}}$} &
\textbf{Marker (TBD from production)} &
\textbf{Condition label} \\
\midrule
low  & low  & best-choice for $(\text{low}, g_{\text{low}})$  &
  (0) community agrees, weak signal \\
low  & high & best-choice for $(\text{low}, g_{\text{high}})$ &
  (1) community agrees, strong signal \\
high & low  & best-choice for $(\text{high}, g_{\text{low}})$ &
  (2) community divided, weak signal \\
high & high & best-choice for $(\text{high}, g_{\text{high}})$ &
  (3) community divided, strong signal \\
\bottomrule
\end{tabular}
\end{table}

\subsubsection{Counterbalancing}
\label{subsub:counterbalancing_comp}

The experiment uses the \textbf{same $4 \times 8$ Latin square} as the
production experiment: each of the 4 lists assigns item $i$
($i = 0,\ldots,7$) to condition $(i + l)\bmod 4$, where
$l \in \{0,1,2,3\}$ is the list number.
This guarantees that:
\begin{itemize}
  \item each participant sees every item exactly once;
  \item each participant sees each of the 4 conditions exactly twice
        (two items per condition), ensuring balanced coverage of the
        $2\times2$ design;
  \item across all 4 lists, each item appears equally often in each
        condition.
\end{itemize}

List assignment follows the same Prolific-and-URL strategy as the production
experiment: each Prolific study embeds \texttt{?list=N} (0--3) in the
experiment URL, capping recruitment at $n = 20$ participants per list
for a total of $N = 80$.

In addition to the 8 critical trials, each participant completed 8
\textbf{filler trials} using the same filler items as the production
experiment (train travel, reading light, cooking, coffee, vacation,
hydration, fresh air, noise), each presented with a randomly assigned
marker.
Fillers are excluded from analysis; they serve to dilute the critical
items and reduce demand characteristics.

\subsubsection{Dependent variables}
\label{subsub:dvs_comp}

Two slider ratings are collected after each vignette (0--100 continuous
scale with no visible numbers, labelled at the endpoints as shown):

\begin{enumerate}
  \item \textbf{Inferred speaker goal strength} (operationalises $\widehat{g}$):
  \begin{quote}
    \textit{Wie stark hast du den Eindruck, dass die Sprecherin m\"ochte,
    dass du ihre Empfehlung befolgst?}\\
    (`How strongly do you get the impression that the speaker wants you to
    follow their recommendation?')\\
    Endpoints: 0 = \textit{\"uberhaupt nicht}
    (`not at all') ---
    100 = \textit{sehr stark} (`very strongly').
  \end{quote}

  \item \textbf{Adoption likelihood of $q$} (operationalises $\widehat{A}$):
  \begin{quote}
    \textit{Wie wahrscheinlich ist es, dass du die empfohlene Handlung
    dann auch tats\"achlich umsetzt?}\\
    (`How likely is it that you will actually carry out the recommended
    action?')\\
    Endpoints: 0 = \textit{sehr unwahrscheinlich}
    (`very unlikely') ---
    100 = \textit{sehr wahrscheinlich} (`very likely').
  \end{quote}
\end{enumerate}

\noindent Perceived controversy is deliberately \emph{not} included as a
rating probe: \textsc{pc\textsubscript{prag}} is experimentally
manipulated via the context paragraph, so a controversy rating would
be redundant and risks obscuring the marker's contribution to inferred
goal strength.

Recorded per trial: \texttt{trial\_id}, \texttt{topic},
\texttt{is\_filler}, \texttt{condition\_index}, \texttt{marker},
\texttt{pc\_prag}, \texttt{g\_implied}, \texttt{rating\_g},
\texttt{rating\_adopt}, \texttt{rt} (ms); and once per participant:
\texttt{list\_num}.
The continuous \texttt{pc\_prop\_rating} covariate (topic-level norming
mean) is joined to the trial data during analysis.

\subsection{Procedure}

Participants complete the experiment online via the magpie platform
\citep{kochetkova2024magpie} in a browser.
After reading instructions in German, participants are assigned to one of
4 counterbalancing lists (via URL \texttt{?list=N} parameter assigned by
Prolific; no random fallback in deployment).
They then proceed through 16 trials---8 critical and 8 fillers---in a
fully randomised order, followed by a post-test questionnaire eliciting
native language, age, gender, education level, and optional comments.

Each trial displays:
(a) a context paragraph establishing the communicative situation and the
\textsc{pc\textsubscript{prag}} level;
(b) the target sentence with the best-choice marker already in place in
the Mittelfeld;
(c) the consequent action clause;
(d) the two rating sliders in fixed order ($\widehat{g}$, $\widehat{A}$),
each labelled in German.
The ``Weiter'' (Next) button is disabled until both sliders have been
moved at least once, ensuring a complete response on every trial.

\subsection{Planned Analysis}

\subsubsection{Empirical plots}

For each dependent variable ($\widehat{g}$, $\widehat{A}$):

\begin{itemize}
  \item \textbf{Condition means}: a dot-and-whisker plot shows mean
        ratings per condition (the $2\times2$ of pc\textsubscript{prag}
        $\times$ $g_{\mathrm{implied}}$) with 95\% bootstrapped CIs.
        The expected pattern: $\widehat{g}$ and $\widehat{A}$ should
        be higher in the high-$g_{\mathrm{implied}}$ conditions
        (where the model assigns a stronger marker), and lower in the
        high-pc\textsubscript{prag} conditions (credibility discounting).
  \item \textbf{Marker identity}: because the best-choice marker may
        differ across items within the same condition, a secondary plot
        shows ratings grouped by actual marker identity, allowing
        inspection of whether marker type (rather than just condition)
        drives the effects.
  \item \textbf{pc\textsubscript{prop} covariate}: scatter plots display
        each DV as a function of the continuous norming rating per
        condition, testing whether higher propositional acceptance
        amplifies the $\widehat{g}$ and $\widehat{A}$ signals for
        strong markers.
\end{itemize}

\subsubsection{Bayesian mixed-effects regression}

Two Bayesian mixed-effects regression models are fitted via
\texttt{brms} \citep{burkner2017}, one per DV ($\widehat{g}$ and
$\widehat{A}$):

\begin{align}
y &\sim \mathrm{Beta}(\mu, \phi)
\quad \text{(ratings rescaled to }(0,1)\text{)} \\
\mathrm{logit}(\mu) &= \alpha
+ \beta_{g} \cdot \mathrm{g\_implied\_c}
+ \beta_{\mathrm{prag}} \cdot \mathrm{pc\_prag\_c}
+ \beta_{\mathrm{prop}} \cdot \mathrm{pc\_prop\_c}
+ \beta_{\mathrm{int}} \cdot \mathrm{g\_implied\_c}
  \times \mathrm{pc\_prag\_c}
+ (1 \mid \mathrm{submission\_id})
+ (1 \mid \mathrm{topic}),
\end{align}

where \texttt{g\_implied\_c} is the $g$ level from the corresponding
production condition ($-0.5$ = low, $+0.5$ = high), effectively
coding the marker's position on the production utility scale.
A secondary specification replaces \texttt{g\_implied\_c} with
actualy marker identity as an unordered factor, to examine
non-monotonic patterns and verify that marker rather than
condition label drives effects.

Model comparison via approximate leave-one-out cross-validation (LOO)
selects among additive, interaction, and marker-as-categorical
formulations.

\subsubsection{RSA listener parameter estimation}
\label{subsub:rsa_fit}

The two empirical rating means (per condition cell: 4 conditions) constrain
the listener parameters of the noisy-threshold RSA model.
Specifically, the model predicts:

\begin{align}
\mathbb{E}[\widehat{g} \mid u, \text{context}]
  &= \int g \cdot P(g \mid u, pc_{\mathrm{prag}})\, dg,
\label{eq:pred_g} \\
P(\widehat{A} \mid u, \text{context})
  &= \iiint \sigma(\eta_0 + \eta_g g - \eta_{pc}\, pc
     - \eta_{\mathrm{int}}\, g \cdot pc)\;
     P(pc_{\mathrm{prop}}, pc_{\mathrm{prag}}, g \mid u)\,
     d{pc_{\mathrm{prop}}}\,d{pc_{\mathrm{prag}}}\,dg,
\label{eq:adoption}
\end{align}

\noindent where the posterior $P(pc_{\mathrm{prop}}, pc_{\mathrm{prag}}, g
\mid u)$ is conditioned on the context's \textsc{pc\textsubscript{prag}}
level as a known prior (since it is explicitly given to the listener).

The threshold parameters $(\mu_1,\mu_2, \sigma_1,\sigma_2,
\lambda)$ together with the marker-specific costs $(c_u)$ are estimated
jointly from the production experiment data via the Stan model.
Given those threshold estimates, the listener parameters
$(\eta_0, \eta_g, \eta_{pc}, \eta_{\mathrm{int}})$ are then estimated by
minimising the weighted squared deviation between the model-predicted
means (Eqs.~\ref{eq:pred_g}--\ref{eq:adoption}) and the empirical rating
means from the comprehension experiment, or via a full joint Bayesian
fit over both experiments simultaneously.

\subsection{Model Predictions}

The noisy-threshold RSA model makes the following predictions for
$\widehat{g}$ and $\widehat{A}$, which serve as pre-registered hypotheses.
In all cases, ``strong-marker condition'' refers to conditions 1 and 3
(high $g_{\mathrm{implied}}$, i.e.\ the marker selected by the production
model under a high-$g$ context), and ``weak-marker condition'' to conditions
0 and 2.

\begin{enumerate}
  \item \textbf{Marker encodes inferred goal $\widehat{g}$}: $\widehat{g}$
        ratings should be higher in strong-marker conditions than in
        weak-marker conditions, across both pc\textsubscript{prag} levels.
        This is the core inversion prediction: if the production model
        correctly captures speaker behaviour, listeners should be able to
        decode $g$ from the marker.

  \item \textbf{Marker encodes adoption $\widehat{A}$}: $\widehat{A}$
        ratings should likewise be higher in strong-marker conditions,
        reflecting that higher inferred $g$ translates into greater
        adoption pressure downstream.

  \item \textbf{pc\textsubscript{prag} modulates $\widehat{A}$ but not
        $\widehat{g}$}: because $g$ is orthogonal to pc\textsubscript{prag}
        in the model's utility function, the effect of the marker on
        $\widehat{g}$ should be largely the same in low- and high-
        pc\textsubscript{prag} contexts.
        By contrast, $\widehat{A}$ should be reduced in high-
        pc\textsubscript{prag} contexts via the interaction term
        $\eta_{\mathrm{int}}\,g \cdot pc$: even a strong-marker signal
        for high $g$ is partly cancelled when controversy is high.

  \item \textbf{Credibility discounting}: for strong-marker conditions
        presented in a high-pc\textsubscript{prag} context (condition~3),
        $\widehat{A}$ should be \emph{lower} than for the same marker in
        a low-pc\textsubscript{prag} context (condition~1).
        This is the model's signature prediction distinguishing Bayesian
        inversion from a simpler marker-strength heuristic: a sceptical
        listener who knows controversy is high will discount the adoption
        implication of even the strongest marker.
        The $g_{\mathrm{implied}} \times \textsc{pc\textsubscript{prag}}$
        interaction on $\widehat{A}$ provides the empirical test.
\end{enumerate}

\end{document}
   % (inside a \section{Experiment} environment)

\documentclass[12pt]{article}
\usepackage{CSPArticleStyle}
\usepackage{csquotes}
\usepackage[margin=1in]{geometry}
\usepackage{booktabs}
\usepackage{parskip}
\usepackage{amsmath}
\usepackage{times}
\usepackage[T1]{fontenc}
\usepackage[utf8]{inputenc}
\usepackage{xcolor}
\usepackage{tabularray}
\usepackage{gb4e}
\noautomath

\title{Comprehension Experiment: Inferring Speaker Goal and Action Adoption
       from Consensus-Marker Use}
\author{}
\date{}

\begin{document}

%% ============================================================
%% When embedding in main.tex, delete from \documentclass to
%% \begin{document} and delete \end{document} at the bottom.
%% ============================================================

\section{Experiment: Consensus-Marker Comprehension}

\subsection{Overview}

The comprehension experiment is the listener-side counterpart of the
production experiment and tests whether listeners can invert the
speaker's marker choice to recover the latent context that motivated it.
Where the production experiment gave participants a context (fixing
\textsc{pc\textsubscript{prag}} and $g$) and asked which marker they
would use, the comprehension experiment gives participants a context
(fixing \textsc{pc\textsubscript{prag}}) together with a marker already
embedded in the sentence and asks what goal strength the speaker likely
had and how likely the listener is to adopt $q$.

The specific marker presented in each vignette is not chosen arbitrarily
from the three candidate markers. Our preregistered plan was to assign,
for each item$\times$condition cell, the \emph{Bayesian best-choice marker}:
the marker with the highest posterior probability under the cost-augmented
noisy-threshold RSA model fitted to the production data, given the context's
\textsc{pc\textsubscript{prag}} level and a specified $g$ level. This was
intended to mirror the production design closely and to ensure that the
comprehension stimuli reflect model-based marker use rather than arbitrary
marker$\times$item pairings.

Applying this strict modal-assignment rule to the fitted production model,
however, yielded a highly imbalanced stimulus set: most critical cells were
assigned the strongest marker, with very little variation across marker
categories. In addition, many of these modal decisions were supported only by
very small probability margins over the next-best alternative. Because such a
stimulus set would make the comprehension experiment substantially less
informative about marker-based listener inferences, we introduced a documented,
exploratory balancing step for stimulus construction. Starting from the same
posterior predictive probabilities, we selected a near-balanced assignment of
the three markers across the 32 critical cells while maximizing total posterior
support under the fitted production model. The strict modal assignment is
retained as the preregistered baseline, but the balanced assignment is used for
the actual comprehension materials.

\textsc{pc\textsubscript{prag}} is retained as an explicit context
manipulation (low vs.\ high), so that listeners can be expected to
have a well-defined prior on pragmatic controversy when interpreting
the marker. \textsc{pc\textsubscript{prop}} is not manipulated within
the experiment; it enters as a continuous topic-level covariate from
the norming study.

After reading each vignette, participants provide two slider ratings
operationalising the key listener-side quantities:
\begin{enumerate}
  \item \textbf{Inferred speaker goal strength} ($\widehat{g}$):
        how strongly does the speaker seem to want the listener to
        accept and act on $q$?
  \item \textbf{Adoption likelihood} ($\widehat{A}$):
        how likely is the listener to follow the recommended action $q$?
\end{enumerate}

\noindent Perceived controversy is \emph{not} elicited as a rating:
because \textsc{pc\textsubscript{prag}} is already manipulated as a
context factor, asking participants to rate controversy would be
redundant with respect to the manipulation check and circular with
respect to the research question. The theoretically central inference
to test is whether listeners can decode the speaker's \emph{goal} from
marker use, and whether that inference translates into adoption
behaviour.

The two slider measures, together with the \textsc{pc\textsubscript{prag}}
context and the \textsc{pc\textsubscript{prop}} covariate, provide the
data needed to fit the model's listener parameters
$(\eta_0, \eta_g, \eta_{pc}, \eta_{\mathrm{int}})$.

\subsection{Participants}

Participants will be recruited via Prolific and will be required to be
native speakers of German (language-at-home filter).
Data collection is planned after completion of the production experiment
(model fitting required to determine the best-choice markers).
A target of $N = 80$ participants is planned ($n = 20$ per list across
4 counterbalancing lists), each completing 16 trials (8 critical + 8
filler), matching the structure of the production experiment.

\subsection{Materials}

\subsubsection{Markers}

Each critical vignette is presented with one specific consensus marker
already embedded in the Mittelfeld of the sentence frame (see
example~(\ref{ex:frame_comp})):

\begin{exe}
  \ex \label{ex:frame_comp}
  Digitale Kompetenzen sind \textbf{[MARKER]}
  in der heutigen Arbeitswelt unverzichtbar.\\
  `Digital skills are \textbf{[MARKER]} indispensable in today's
  working world.'
\end{exe}

Unlike the production experiment, participants do \emph{not} choose the
marker; they read the sentence with the marker already in place and respond
to the two rating probes.

Critically, the marker assigned to each condition is not drawn arbitrarily
from the three candidates. Under the preregistered plan, for each of the four
production conditions (pc\textsubscript{prag} $\times$ $g$), the marker used in
comprehension was to be the \textbf{Bayesian best-choice marker}: the marker
$u^*$ with the highest posterior probability under the fitted
cost-augmented noisy-threshold RSA production model,
\[
u^* = \arg\max_{u} \; P(u \mid pc_{\mathrm{prag}}, g,
                           pc_{\mathrm{prop}\_c}).
\]
Because $pc_{\mathrm{prop}}$ varies continuously across items,
$u^*$ may differ per item even within the same
(pc\textsubscript{prag}, $g$) cell; the final stimulus list is therefore
constructed item-by-item once the production model has been fitted.
In practice, the fitted production model yielded a nearly degenerate modal map,
with most cells assigned to \textit{bekanntlich}. We therefore retained the
model-based logic of the procedure but replaced strict cellwise modal selection
with an uncertainty-aware balancing step: among the marker assignments with
high posterior support, we selected a near-balanced configuration across the 32
critical cells. This preserves the production--comprehension link while also
ensuring that the comprehension experiment contains sufficient marker variation
to test listener-side inferences.

All three markers remain candidates; the mapping from conditions to
markers is an empirical output of the production analysis, not a design
pre-commitment. In the final comprehension materials, we use the
\textbf{balanced model-based assignment}, rather than the strict modal
assignment, in order to maintain informative variation across marker types.

\subsubsection{Topic Items}

The same eight topic domains used in the norming study and the production
experiment are used here (Table~\ref{tab:items_comp}).
The 4-condition $\times$ 8-item Latin square (Section~\ref{subsub:counterbalancing_comp})
requires no additional items.

\begin{table}[ht]
\centering
\caption{Topic items: sentence frame for $p$ (with marker position [M])
         and consequent $q$.}
\label{tab:items_comp}
\small
\begin{tabular}{clll}
\toprule
\textbf{\#} & \textbf{Topic} & \textbf{Sentence frame} & \textbf{Consequent $q$} \\
\midrule
1 & Klimawandel &
  \textit{Der Klimawandel wird [M]} &
  auf Flugreisen verzichten \\
  & & \textit{durch menschliche Aktivit\"aten verursacht.} & \\[3pt]
2 & Ern\"ahrung &
  \textit{Eine \"uberwiegend pflanzliche Ern\"ahrung} &
  mehr pflanzliche Lebens- \\
  & & \textit{f\"ordert [M] die Gesundheit.} &
  mittel einbauen \\[3pt]
3 & Stadtverkehr &
  \textit{Ein ausgebautes Nahverkehrsnetz} &
  h\"aufiger auf \\
  & & \textit{entlastet [M] den Stadtverkehr} &
  \"offentliche Verkehrs- \\
  & & \textit{deutlich.} &
  mittel umsteigen \\[3pt]
4 & Digitalisierung &
  \textit{Digitale Kompetenzen sind [M] in} &
  regelm\"a\ss{}ig Weiter- \\
  & & \textit{der heutigen Arbeitswelt unverzichtbar.} &
  bildungsangebote nutzen \\[3pt]
5 & Schlaf &
  \textit{Ausreichend Schlaf ist [M] entscheidend} &
  auf feste Schlafzeiten \\
  & & \textit{f\"ur die kognitive Leistungsf\"ahigkeit.} &
  achten \\[3pt]
6 & Plastik &
  \textit{Einwegplastik schadet [M] den} &
  konsequent auf Einweg- \\
  & & \textit{Meeres\"okosystemen erheblich.} &
  plastik verzichten \\[3pt]
7 & Sport &
  \textit{Regelm\"a\ss{}ige k\"orperliche Bewegung} &
  regelm\"a\ss{}ig Sport in \\
  & & \textit{senkt [M] das Risiko f\"ur} &
  den Alltag integrieren \\
  & & \textit{Herz-Kreislauf-Erkrankungen deutlich.} & \\[3pt]
8 & Lokalkauf &
  \textit{Das Kaufen bei lokalen H\"andlern} &
  \"ofter bei lokalen \\
  & & \textit{st\"arkt [M] die regionale Wirtschaft} &
  Gesch\"aften einkaufen \\
  & & \textit{nachhaltig.} & \\
\bottomrule
\end{tabular}
\end{table}

\noindent [M] denotes the position of the pre-filled best-choice marker.

\subsubsection{Vignettes}

\paragraph{Structure.}
Each trial presents a vignette consisting of three parts:
(a) a context paragraph establishing the communicative situation and
fixing the \textsc{pc\textsubscript{prag}} level;
(b) the target sentence with the best-choice marker already embedded
in the Mittelfeld;
(c) the consequent action clause.
The two rating probes follow immediately below the vignette.

\paragraph{Propositional perceived controversy (\textsc{pc\textsubscript{prop}}).}
As in the production experiment, \textsc{pc\textsubscript{prop}} is not
manipulated within the comprehension experiment.
Topic-level propositional acceptance ratings from the norming study
(0--100 scale; e.g., \textit{Digitalisierung}: 85/100,
\textit{Lokalkauf}: 58/100) serve as a continuous, mean-centred
covariate \textsc{pc\textsubscript{prop}}$_c$ in analysis.
Context paragraphs present $p$ without editorial evaluation; variation in
listeners' individual assessments of propositional controversy is
captured by the norming-derived covariate.

\paragraph{Pragmatic perceived controversy (\textsc{pc\textsubscript{prag}}).}
\textsc{pc\textsubscript{prag}} is manipulated between item instances
using the same two context versions as in the production experiment:
the speaker's social circle either \emph{largely agrees} (low) or
\emph{holds divergent views} (high) about $p$.
For example, for the \textit{Digitalisierung} item:

\begin{itemize}
  \item $\textsc{pc}_\text{prag} = \text{low}$: \textit{Deine Kollegin spricht
        mit dir \"uber Berufsausbildung und Weiterbildung.
        In eurem Team teilen die meisten die Ansicht, dass digitale Kompetenzen
        in der heutigen Arbeitswelt unverzichtbar sind.}
  \item $\textsc{pc}_\text{prag} = \text{high}$: \textit{Deine Kollegin spricht
        mit dir \"uber Berufsausbildung und Weiterbildung.
        In eurem Team gehen die Meinungen dazu stark auseinander.}
\end{itemize}

The context is framed from the listener's perspective as a direct addressee,
mirroring the production study where participants adopt the speaker's
perspective.

\paragraph{Marker manipulation.}
The dependent measures of the production experiment (which marker is chosen)
become the independent variable here.
One of the three consensus markers is inserted into the Mittelfeld of the
target sentence.
A complete sample item for the \textit{Digitalisierung} topic
($\textsc{pc}_\text{prag} = \text{high}$, marker = \textit{bekanntlich};
norming $\textsc{pc}_\text{prop}$ rating: 85/100) is given in
(\ref{ex:comp_sample}):

\begin{exe}
  \ex \label{ex:comp_sample}
  \textit{Deine Kollegin spricht mit dir \"uber Berufsausbildung und
  Weiterbildung.
  In eurem Team gehen die Meinungen dazu stark auseinander.}

  \textit{Sie sagt zu dir:} Digitale Kompetenzen sind \textbf{bekanntlich}
  in der heutigen Arbeitswelt unverzichtbar.
  Deshalb solltest du regelm\"a\ss{}ig digitale Weiterbildungsangebote nutzen.

  \noindent\textit{[Rating probes follow.]}
\end{exe}

\subsection{Design}

\subsubsection{Factorial structure}

The comprehension experiment mirrors the production experiment's
$\mathbf{2}$~\textbf{(pc\textsubscript{prag}: low vs.\ high)}
$\mathbf{\times}$~$\mathbf{2}$~\textbf{($g_{\mathrm{implied}}$: low
vs.\ high) design}, yielding \textbf{4 conditions}.
The $g_{\mathrm{implied}}$ dimension is not manipulated via an explicit
goal instruction (as in production) but is instead conveyed entirely
through the choice of marker: for each condition, the marker presented
is the best-choice marker for the corresponding
(\textsc{pc\textsubscript{prag}},~$g$) cell from the production model.
\textsc{pc\textsubscript{prop}} enters as a continuous norming covariate.
With 8 topic items and a Latin square rotation across 4 lists, each
participant sees exactly \textbf{8 critical trials}---one item per
condition, each condition appearing twice---plus 8 filler trials,
for \textbf{16 trials} in total (Table~\ref{tab:design_comp}).

\begin{table}[ht]
\centering
\caption{%
  The $2\times2$ design (4 conditions), mirroring the production experiment.
  The marker presented in each condition is the Bayesian best-choice
  marker from the fitted production model; specific markers are
  determined after production data collection.}
\label{tab:design_comp}
\small
\begin{tabular}{ccll}
\toprule
\textbf{pc\textsubscript{prag}} &
\textbf{$g_{\mathrm{implied}}$} &
\textbf{Marker (TBD from production)} &
\textbf{Condition label} \\
\midrule
low  & low  & best-choice for $(\text{low}, g_{\text{low}})$  &
  (0) community agrees, weak signal \\
low  & high & best-choice for $(\text{low}, g_{\text{high}})$ &
  (1) community agrees, strong signal \\
high & low  & best-choice for $(\text{high}, g_{\text{low}})$ &
  (2) community divided, weak signal \\
high & high & best-choice for $(\text{high}, g_{\text{high}})$ &
  (3) community divided, strong signal \\
\bottomrule
\end{tabular}
\end{table}

\subsubsection{Counterbalancing}
\label{subsub:counterbalancing_comp}

The experiment uses the \textbf{same $4 \times 8$ Latin square} as the
production experiment: each of the 4 lists assigns item $i$
($i = 0,\ldots,7$) to condition $(i + l)\bmod 4$, where
$l \in \{0,1,2,3\}$ is the list number.
This guarantees that:
\begin{itemize}
  \item each participant sees every item exactly once;
  \item each participant sees each of the 4 conditions exactly twice
        (two items per condition), ensuring balanced coverage of the
        $2\times2$ design;
  \item across all 4 lists, each item appears equally often in each
        condition.
\end{itemize}

List assignment follows the same Prolific-and-URL strategy as the production
experiment: each Prolific study embeds \texttt{?list=N} (0--3) in the
experiment URL, capping recruitment at $n = 20$ participants per list
for a total of $N = 80$.

In addition to the 8 critical trials, each participant completed 8
\textbf{filler trials} using the same filler items as the production
experiment (train travel, reading light, cooking, coffee, vacation,
hydration, fresh air, noise), each presented with a randomly assigned
marker.
Fillers are excluded from analysis; they serve to dilute the critical
items and reduce demand characteristics.

\subsubsection{Dependent variables}
\label{subsub:dvs_comp}

Two slider ratings are collected after each vignette (0--100 continuous
scale with no visible numbers, labelled at the endpoints as shown):

\begin{enumerate}
  \item \textbf{Inferred speaker goal strength} (operationalises $\widehat{g}$):
  \begin{quote}
    \textit{Wie stark hast du den Eindruck, dass die Sprecherin m\"ochte,
    dass du ihre Empfehlung befolgst?}\\
    (`How strongly do you get the impression that the speaker wants you to
    follow their recommendation?')\\
    Endpoints: 0 = \textit{\"uberhaupt nicht}
    (`not at all') ---
    100 = \textit{sehr stark} (`very strongly').
  \end{quote}

  \item \textbf{Adoption likelihood of $q$} (operationalises $\widehat{A}$):
  \begin{quote}
    \textit{Wie wahrscheinlich ist es, dass du die empfohlene Handlung
    dann auch tats\"achlich umsetzt?}\\
    (`How likely is it that you will actually carry out the recommended
    action?')\\
    Endpoints: 0 = \textit{sehr unwahrscheinlich}
    (`very unlikely') ---
    100 = \textit{sehr wahrscheinlich} (`very likely').
  \end{quote}
\end{enumerate}

\noindent Perceived controversy is deliberately \emph{not} included as a
rating probe: \textsc{pc\textsubscript{prag}} is experimentally
manipulated via the context paragraph, so a controversy rating would
be redundant and risks obscuring the marker's contribution to inferred
goal strength.

Recorded per trial: \texttt{trial\_id}, \texttt{topic},
\texttt{is\_filler}, \texttt{condition\_index}, \texttt{marker},
\texttt{pc\_prag}, \texttt{g\_implied}, \texttt{rating\_g},
\texttt{rating\_adopt}, \texttt{rt} (ms); and once per participant:
\texttt{list\_num}.
The continuous \texttt{pc\_prop\_rating} covariate (topic-level norming
mean) is joined to the trial data during analysis.

\subsection{Procedure}

Participants complete the experiment online via the magpie platform
\citep{kochetkova2024magpie} in a browser.
After reading instructions in German, participants are assigned to one of
4 counterbalancing lists (via URL \texttt{?list=N} parameter assigned by
Prolific; no random fallback in deployment).
They then proceed through 16 trials---8 critical and 8 fillers---in a
fully randomised order, followed by a post-test questionnaire eliciting
native language, age, gender, education level, and optional comments.

Each trial displays:
(a) a context paragraph establishing the communicative situation and the
\textsc{pc\textsubscript{prag}} level;
(b) the target sentence with the best-choice marker already in place in
the Mittelfeld;
(c) the consequent action clause;
(d) the two rating sliders in fixed order ($\widehat{g}$, $\widehat{A}$),
each labelled in German.
The ``Weiter'' (Next) button is disabled until both sliders have been
moved at least once, ensuring a complete response on every trial.

\subsection{Planned Analysis}

\subsubsection{Empirical plots}

For each dependent variable ($\widehat{g}$, $\widehat{A}$):

\begin{itemize}
  \item \textbf{Condition means}: a dot-and-whisker plot shows mean
        ratings per condition (the $2\times2$ of pc\textsubscript{prag}
        $\times$ $g_{\mathrm{implied}}$) with 95\% bootstrapped CIs.
        The expected pattern: $\widehat{g}$ and $\widehat{A}$ should
        be higher in the high-$g_{\mathrm{implied}}$ conditions
        (where the model assigns a stronger marker), and lower in the
        high-pc\textsubscript{prag} conditions (credibility discounting).
  \item \textbf{Marker identity}: because the best-choice marker may
        differ across items within the same condition, a secondary plot
        shows ratings grouped by actual marker identity, allowing
        inspection of whether marker type (rather than just condition)
        drives the effects.
  \item \textbf{pc\textsubscript{prop} covariate}: scatter plots display
        each DV as a function of the continuous norming rating per
        condition, testing whether higher propositional acceptance
        amplifies the $\widehat{g}$ and $\widehat{A}$ signals for
        strong markers.
\end{itemize}

\subsubsection{Bayesian mixed-effects regression}

Two Bayesian mixed-effects regression models are fitted via
\texttt{brms} \citep{burkner2017}, one per DV ($\widehat{g}$ and
$\widehat{A}$):

\begin{align}
y &\sim \mathrm{Beta}(\mu, \phi)
\quad \text{(ratings rescaled to }(0,1)\text{)} \\
\mathrm{logit}(\mu) &= \alpha
+ \beta_{g} \cdot \mathrm{g\_implied\_c}
+ \beta_{\mathrm{prag}} \cdot \mathrm{pc\_prag\_c}
+ \beta_{\mathrm{prop}} \cdot \mathrm{pc\_prop\_c}
+ \beta_{\mathrm{int}} \cdot \mathrm{g\_implied\_c}
  \times \mathrm{pc\_prag\_c}
+ (1 \mid \mathrm{submission\_id})
+ (1 \mid \mathrm{topic}),
\end{align}

where \texttt{g\_implied\_c} is the $g$ level from the corresponding
production condition ($-0.5$ = low, $+0.5$ = high), effectively
coding the marker's position on the production utility scale.
A secondary specification replaces \texttt{g\_implied\_c} with
actualy marker identity as an unordered factor, to examine
non-monotonic patterns and verify that marker rather than
condition label drives effects.

Model comparison via approximate leave-one-out cross-validation (LOO)
selects among additive, interaction, and marker-as-categorical
formulations.

\subsubsection{RSA listener parameter estimation}
\label{subsub:rsa_fit}

The two empirical rating means (per condition cell: 4 conditions) constrain
the listener parameters of the noisy-threshold RSA model.
Specifically, the model predicts:

\begin{align}
\mathbb{E}[\widehat{g} \mid u, \text{context}]
  &= \int g \cdot P(g \mid u, pc_{\mathrm{prag}})\, dg,
\label{eq:pred_g} \\
P(\widehat{A} \mid u, \text{context})
  &= \iiint \sigma(\eta_0 + \eta_g g - \eta_{pc}\, pc
     - \eta_{\mathrm{int}}\, g \cdot pc)\;
     P(pc_{\mathrm{prop}}, pc_{\mathrm{prag}}, g \mid u)\,
     d{pc_{\mathrm{prop}}}\,d{pc_{\mathrm{prag}}}\,dg,
\label{eq:adoption}
\end{align}

\noindent where the posterior $P(pc_{\mathrm{prop}}, pc_{\mathrm{prag}}, g
\mid u)$ is conditioned on the context's \textsc{pc\textsubscript{prag}}
level as a known prior (since it is explicitly given to the listener).

The threshold parameters $(\mu_1,\mu_2, \sigma_1,\sigma_2,
\lambda)$ together with the marker-specific costs $(c_u)$ are estimated
jointly from the production experiment data via the Stan model.
Given those threshold estimates, the listener parameters
$(\eta_0, \eta_g, \eta_{pc}, \eta_{\mathrm{int}})$ are then estimated by
minimising the weighted squared deviation between the model-predicted
means (Eqs.~\ref{eq:pred_g}--\ref{eq:adoption}) and the empirical rating
means from the comprehension experiment, or via a full joint Bayesian
fit over both experiments simultaneously.

\subsection{Model Predictions}

The noisy-threshold RSA model makes the following predictions for
$\widehat{g}$ and $\widehat{A}$, which serve as pre-registered hypotheses.
In all cases, ``strong-marker condition'' refers to conditions 1 and 3
(high $g_{\mathrm{implied}}$, i.e.\ the marker selected by the production
model under a high-$g$ context), and ``weak-marker condition'' to conditions
0 and 2.

\begin{enumerate}
  \item \textbf{Marker encodes inferred goal $\widehat{g}$}: $\widehat{g}$
        ratings should be higher in strong-marker conditions than in
        weak-marker conditions, across both pc\textsubscript{prag} levels.
        This is the core inversion prediction: if the production model
        correctly captures speaker behaviour, listeners should be able to
        decode $g$ from the marker.

  \item \textbf{Marker encodes adoption $\widehat{A}$}: $\widehat{A}$
        ratings should likewise be higher in strong-marker conditions,
        reflecting that higher inferred $g$ translates into greater
        adoption pressure downstream.

  \item \textbf{pc\textsubscript{prag} modulates $\widehat{A}$ but not
        $\widehat{g}$}: because $g$ is orthogonal to pc\textsubscript{prag}
        in the model's utility function, the effect of the marker on
        $\widehat{g}$ should be largely the same in low- and high-
        pc\textsubscript{prag} contexts.
        By contrast, $\widehat{A}$ should be reduced in high-
        pc\textsubscript{prag} contexts via the interaction term
        $\eta_{\mathrm{int}}\,g \cdot pc$: even a strong-marker signal
        for high $g$ is partly cancelled when controversy is high.

  \item \textbf{Credibility discounting}: for strong-marker conditions
        presented in a high-pc\textsubscript{prag} context (condition~3),
        $\widehat{A}$ should be \emph{lower} than for the same marker in
        a low-pc\textsubscript{prag} context (condition~1).
        This is the model's signature prediction distinguishing Bayesian
        inversion from a simpler marker-strength heuristic: a sceptical
        listener who knows controversy is high will discount the adoption
        implication of even the strongest marker.
        The $g_{\mathrm{implied}} \times \textsc{pc\textsubscript{prag}}$
        interaction on $\widehat{A}$ provides the empirical test.
\end{enumerate}

\end{document}
.
%
% Compile stand-alone:
%   pdflatex comprehension_exp.tex
%
% Embed in main.tex:
%   % comprehension_exp.tex
% Stand-alone plan for the comprehension experiment.
% Can be embedded in main.tex via % comprehension_exp.tex
% Stand-alone plan for the comprehension experiment.
% Can be embedded in main.tex via % comprehension_exp.tex
% Stand-alone plan for the comprehension experiment.
% Can be embedded in main.tex via \input{comprehension_exp}.
%
% Compile stand-alone:
%   pdflatex comprehension_exp.tex
%
% Embed in main.tex:
%   \input{comprehension_exp}   % (inside a \section{Experiment} environment)

\documentclass[12pt]{article}
\usepackage{CSPArticleStyle}
\usepackage{csquotes}
\usepackage[margin=1in]{geometry}
\usepackage{booktabs}
\usepackage{parskip}
\usepackage{amsmath}
\usepackage{times}
\usepackage[T1]{fontenc}
\usepackage[utf8]{inputenc}
\usepackage{xcolor}
\usepackage{tabularray}
\usepackage{gb4e}
\noautomath

\title{Comprehension Experiment: Inferring Speaker Goal and Action Adoption
       from Consensus-Marker Use}
\author{}
\date{}

\begin{document}

%% ============================================================
%% When embedding in main.tex, delete from \documentclass to
%% \begin{document} and delete \end{document} at the bottom.
%% ============================================================

\section{Experiment: Consensus-Marker Comprehension}

\subsection{Overview}

The comprehension experiment is the listener-side counterpart of the
production experiment and tests whether listeners can invert the
speaker's marker choice to recover the latent context that motivated it.
Where the production experiment gave participants a context (fixing
\textsc{pc\textsubscript{prag}} and $g$) and asked which marker they
would use, the comprehension experiment gives participants a context
(fixing \textsc{pc\textsubscript{prag}}) together with a marker already
embedded in the sentence and asks what goal strength the speaker likely
had and how likely the listener is to adopt $q$.

The specific marker presented in each vignette is not chosen arbitrarily
from the three candidate markers. Our preregistered plan was to assign,
for each item$\times$condition cell, the \emph{Bayesian best-choice marker}:
the marker with the highest posterior probability under the cost-augmented
noisy-threshold RSA model fitted to the production data, given the context's
\textsc{pc\textsubscript{prag}} level and a specified $g$ level. This was
intended to mirror the production design closely and to ensure that the
comprehension stimuli reflect model-based marker use rather than arbitrary
marker$\times$item pairings.

Applying this strict modal-assignment rule to the fitted production model,
however, yielded a highly imbalanced stimulus set: most critical cells were
assigned the strongest marker, with very little variation across marker
categories. In addition, many of these modal decisions were supported only by
very small probability margins over the next-best alternative. Because such a
stimulus set would make the comprehension experiment substantially less
informative about marker-based listener inferences, we introduced a documented,
exploratory balancing step for stimulus construction. Starting from the same
posterior predictive probabilities, we selected a near-balanced assignment of
the three markers across the 32 critical cells while maximizing total posterior
support under the fitted production model. The strict modal assignment is
retained as the preregistered baseline, but the balanced assignment is used for
the actual comprehension materials.

\textsc{pc\textsubscript{prag}} is retained as an explicit context
manipulation (low vs.\ high), so that listeners can be expected to
have a well-defined prior on pragmatic controversy when interpreting
the marker. \textsc{pc\textsubscript{prop}} is not manipulated within
the experiment; it enters as a continuous topic-level covariate from
the norming study.

After reading each vignette, participants provide two slider ratings
operationalising the key listener-side quantities:
\begin{enumerate}
  \item \textbf{Inferred speaker goal strength} ($\widehat{g}$):
        how strongly does the speaker seem to want the listener to
        accept and act on $q$?
  \item \textbf{Adoption likelihood} ($\widehat{A}$):
        how likely is the listener to follow the recommended action $q$?
\end{enumerate}

\noindent Perceived controversy is \emph{not} elicited as a rating:
because \textsc{pc\textsubscript{prag}} is already manipulated as a
context factor, asking participants to rate controversy would be
redundant with respect to the manipulation check and circular with
respect to the research question. The theoretically central inference
to test is whether listeners can decode the speaker's \emph{goal} from
marker use, and whether that inference translates into adoption
behaviour.

The two slider measures, together with the \textsc{pc\textsubscript{prag}}
context and the \textsc{pc\textsubscript{prop}} covariate, provide the
data needed to fit the model's listener parameters
$(\eta_0, \eta_g, \eta_{pc}, \eta_{\mathrm{int}})$.

\subsection{Participants}

Participants will be recruited via Prolific and will be required to be
native speakers of German (language-at-home filter).
Data collection is planned after completion of the production experiment
(model fitting required to determine the best-choice markers).
A target of $N = 80$ participants is planned ($n = 20$ per list across
4 counterbalancing lists), each completing 16 trials (8 critical + 8
filler), matching the structure of the production experiment.

\subsection{Materials}

\subsubsection{Markers}

Each critical vignette is presented with one specific consensus marker
already embedded in the Mittelfeld of the sentence frame (see
example~(\ref{ex:frame_comp})):

\begin{exe}
  \ex \label{ex:frame_comp}
  Digitale Kompetenzen sind \textbf{[MARKER]}
  in der heutigen Arbeitswelt unverzichtbar.\\
  `Digital skills are \textbf{[MARKER]} indispensable in today's
  working world.'
\end{exe}

Unlike the production experiment, participants do \emph{not} choose the
marker; they read the sentence with the marker already in place and respond
to the two rating probes.

Critically, the marker assigned to each condition is not drawn arbitrarily
from the three candidates. Under the preregistered plan, for each of the four
production conditions (pc\textsubscript{prag} $\times$ $g$), the marker used in
comprehension was to be the \textbf{Bayesian best-choice marker}: the marker
$u^*$ with the highest posterior probability under the fitted
cost-augmented noisy-threshold RSA production model,
\[
u^* = \arg\max_{u} \; P(u \mid pc_{\mathrm{prag}}, g,
                           pc_{\mathrm{prop}\_c}).
\]
Because $pc_{\mathrm{prop}}$ varies continuously across items,
$u^*$ may differ per item even within the same
(pc\textsubscript{prag}, $g$) cell; the final stimulus list is therefore
constructed item-by-item once the production model has been fitted.
In practice, the fitted production model yielded a nearly degenerate modal map,
with most cells assigned to \textit{bekanntlich}. We therefore retained the
model-based logic of the procedure but replaced strict cellwise modal selection
with an uncertainty-aware balancing step: among the marker assignments with
high posterior support, we selected a near-balanced configuration across the 32
critical cells. This preserves the production--comprehension link while also
ensuring that the comprehension experiment contains sufficient marker variation
to test listener-side inferences.

All three markers remain candidates; the mapping from conditions to
markers is an empirical output of the production analysis, not a design
pre-commitment. In the final comprehension materials, we use the
\textbf{balanced model-based assignment}, rather than the strict modal
assignment, in order to maintain informative variation across marker types.

\subsubsection{Topic Items}

The same eight topic domains used in the norming study and the production
experiment are used here (Table~\ref{tab:items_comp}).
The 4-condition $\times$ 8-item Latin square (Section~\ref{subsub:counterbalancing_comp})
requires no additional items.

\begin{table}[ht]
\centering
\caption{Topic items: sentence frame for $p$ (with marker position [M])
         and consequent $q$.}
\label{tab:items_comp}
\small
\begin{tabular}{clll}
\toprule
\textbf{\#} & \textbf{Topic} & \textbf{Sentence frame} & \textbf{Consequent $q$} \\
\midrule
1 & Klimawandel &
  \textit{Der Klimawandel wird [M]} &
  auf Flugreisen verzichten \\
  & & \textit{durch menschliche Aktivit\"aten verursacht.} & \\[3pt]
2 & Ern\"ahrung &
  \textit{Eine \"uberwiegend pflanzliche Ern\"ahrung} &
  mehr pflanzliche Lebens- \\
  & & \textit{f\"ordert [M] die Gesundheit.} &
  mittel einbauen \\[3pt]
3 & Stadtverkehr &
  \textit{Ein ausgebautes Nahverkehrsnetz} &
  h\"aufiger auf \\
  & & \textit{entlastet [M] den Stadtverkehr} &
  \"offentliche Verkehrs- \\
  & & \textit{deutlich.} &
  mittel umsteigen \\[3pt]
4 & Digitalisierung &
  \textit{Digitale Kompetenzen sind [M] in} &
  regelm\"a\ss{}ig Weiter- \\
  & & \textit{der heutigen Arbeitswelt unverzichtbar.} &
  bildungsangebote nutzen \\[3pt]
5 & Schlaf &
  \textit{Ausreichend Schlaf ist [M] entscheidend} &
  auf feste Schlafzeiten \\
  & & \textit{f\"ur die kognitive Leistungsf\"ahigkeit.} &
  achten \\[3pt]
6 & Plastik &
  \textit{Einwegplastik schadet [M] den} &
  konsequent auf Einweg- \\
  & & \textit{Meeres\"okosystemen erheblich.} &
  plastik verzichten \\[3pt]
7 & Sport &
  \textit{Regelm\"a\ss{}ige k\"orperliche Bewegung} &
  regelm\"a\ss{}ig Sport in \\
  & & \textit{senkt [M] das Risiko f\"ur} &
  den Alltag integrieren \\
  & & \textit{Herz-Kreislauf-Erkrankungen deutlich.} & \\[3pt]
8 & Lokalkauf &
  \textit{Das Kaufen bei lokalen H\"andlern} &
  \"ofter bei lokalen \\
  & & \textit{st\"arkt [M] die regionale Wirtschaft} &
  Gesch\"aften einkaufen \\
  & & \textit{nachhaltig.} & \\
\bottomrule
\end{tabular}
\end{table}

\noindent [M] denotes the position of the pre-filled best-choice marker.

\subsubsection{Vignettes}

\paragraph{Structure.}
Each trial presents a vignette consisting of three parts:
(a) a context paragraph establishing the communicative situation and
fixing the \textsc{pc\textsubscript{prag}} level;
(b) the target sentence with the best-choice marker already embedded
in the Mittelfeld;
(c) the consequent action clause.
The two rating probes follow immediately below the vignette.

\paragraph{Propositional perceived controversy (\textsc{pc\textsubscript{prop}}).}
As in the production experiment, \textsc{pc\textsubscript{prop}} is not
manipulated within the comprehension experiment.
Topic-level propositional acceptance ratings from the norming study
(0--100 scale; e.g., \textit{Digitalisierung}: 85/100,
\textit{Lokalkauf}: 58/100) serve as a continuous, mean-centred
covariate \textsc{pc\textsubscript{prop}}$_c$ in analysis.
Context paragraphs present $p$ without editorial evaluation; variation in
listeners' individual assessments of propositional controversy is
captured by the norming-derived covariate.

\paragraph{Pragmatic perceived controversy (\textsc{pc\textsubscript{prag}}).}
\textsc{pc\textsubscript{prag}} is manipulated between item instances
using the same two context versions as in the production experiment:
the speaker's social circle either \emph{largely agrees} (low) or
\emph{holds divergent views} (high) about $p$.
For example, for the \textit{Digitalisierung} item:

\begin{itemize}
  \item $\textsc{pc}_\text{prag} = \text{low}$: \textit{Deine Kollegin spricht
        mit dir \"uber Berufsausbildung und Weiterbildung.
        In eurem Team teilen die meisten die Ansicht, dass digitale Kompetenzen
        in der heutigen Arbeitswelt unverzichtbar sind.}
  \item $\textsc{pc}_\text{prag} = \text{high}$: \textit{Deine Kollegin spricht
        mit dir \"uber Berufsausbildung und Weiterbildung.
        In eurem Team gehen die Meinungen dazu stark auseinander.}
\end{itemize}

The context is framed from the listener's perspective as a direct addressee,
mirroring the production study where participants adopt the speaker's
perspective.

\paragraph{Marker manipulation.}
The dependent measures of the production experiment (which marker is chosen)
become the independent variable here.
One of the three consensus markers is inserted into the Mittelfeld of the
target sentence.
A complete sample item for the \textit{Digitalisierung} topic
($\textsc{pc}_\text{prag} = \text{high}$, marker = \textit{bekanntlich};
norming $\textsc{pc}_\text{prop}$ rating: 85/100) is given in
(\ref{ex:comp_sample}):

\begin{exe}
  \ex \label{ex:comp_sample}
  \textit{Deine Kollegin spricht mit dir \"uber Berufsausbildung und
  Weiterbildung.
  In eurem Team gehen die Meinungen dazu stark auseinander.}

  \textit{Sie sagt zu dir:} Digitale Kompetenzen sind \textbf{bekanntlich}
  in der heutigen Arbeitswelt unverzichtbar.
  Deshalb solltest du regelm\"a\ss{}ig digitale Weiterbildungsangebote nutzen.

  \noindent\textit{[Rating probes follow.]}
\end{exe}

\subsection{Design}

\subsubsection{Factorial structure}

The comprehension experiment mirrors the production experiment's
$\mathbf{2}$~\textbf{(pc\textsubscript{prag}: low vs.\ high)}
$\mathbf{\times}$~$\mathbf{2}$~\textbf{($g_{\mathrm{implied}}$: low
vs.\ high) design}, yielding \textbf{4 conditions}.
The $g_{\mathrm{implied}}$ dimension is not manipulated via an explicit
goal instruction (as in production) but is instead conveyed entirely
through the choice of marker: for each condition, the marker presented
is the best-choice marker for the corresponding
(\textsc{pc\textsubscript{prag}},~$g$) cell from the production model.
\textsc{pc\textsubscript{prop}} enters as a continuous norming covariate.
With 8 topic items and a Latin square rotation across 4 lists, each
participant sees exactly \textbf{8 critical trials}---one item per
condition, each condition appearing twice---plus 8 filler trials,
for \textbf{16 trials} in total (Table~\ref{tab:design_comp}).

\begin{table}[ht]
\centering
\caption{%
  The $2\times2$ design (4 conditions), mirroring the production experiment.
  The marker presented in each condition is the Bayesian best-choice
  marker from the fitted production model; specific markers are
  determined after production data collection.}
\label{tab:design_comp}
\small
\begin{tabular}{ccll}
\toprule
\textbf{pc\textsubscript{prag}} &
\textbf{$g_{\mathrm{implied}}$} &
\textbf{Marker (TBD from production)} &
\textbf{Condition label} \\
\midrule
low  & low  & best-choice for $(\text{low}, g_{\text{low}})$  &
  (0) community agrees, weak signal \\
low  & high & best-choice for $(\text{low}, g_{\text{high}})$ &
  (1) community agrees, strong signal \\
high & low  & best-choice for $(\text{high}, g_{\text{low}})$ &
  (2) community divided, weak signal \\
high & high & best-choice for $(\text{high}, g_{\text{high}})$ &
  (3) community divided, strong signal \\
\bottomrule
\end{tabular}
\end{table}

\subsubsection{Counterbalancing}
\label{subsub:counterbalancing_comp}

The experiment uses the \textbf{same $4 \times 8$ Latin square} as the
production experiment: each of the 4 lists assigns item $i$
($i = 0,\ldots,7$) to condition $(i + l)\bmod 4$, where
$l \in \{0,1,2,3\}$ is the list number.
This guarantees that:
\begin{itemize}
  \item each participant sees every item exactly once;
  \item each participant sees each of the 4 conditions exactly twice
        (two items per condition), ensuring balanced coverage of the
        $2\times2$ design;
  \item across all 4 lists, each item appears equally often in each
        condition.
\end{itemize}

List assignment follows the same Prolific-and-URL strategy as the production
experiment: each Prolific study embeds \texttt{?list=N} (0--3) in the
experiment URL, capping recruitment at $n = 20$ participants per list
for a total of $N = 80$.

In addition to the 8 critical trials, each participant completed 8
\textbf{filler trials} using the same filler items as the production
experiment (train travel, reading light, cooking, coffee, vacation,
hydration, fresh air, noise), each presented with a randomly assigned
marker.
Fillers are excluded from analysis; they serve to dilute the critical
items and reduce demand characteristics.

\subsubsection{Dependent variables}
\label{subsub:dvs_comp}

Two slider ratings are collected after each vignette (0--100 continuous
scale with no visible numbers, labelled at the endpoints as shown):

\begin{enumerate}
  \item \textbf{Inferred speaker goal strength} (operationalises $\widehat{g}$):
  \begin{quote}
    \textit{Wie stark hast du den Eindruck, dass die Sprecherin m\"ochte,
    dass du ihre Empfehlung befolgst?}\\
    (`How strongly do you get the impression that the speaker wants you to
    follow their recommendation?')\\
    Endpoints: 0 = \textit{\"uberhaupt nicht}
    (`not at all') ---
    100 = \textit{sehr stark} (`very strongly').
  \end{quote}

  \item \textbf{Adoption likelihood of $q$} (operationalises $\widehat{A}$):
  \begin{quote}
    \textit{Wie wahrscheinlich ist es, dass du die empfohlene Handlung
    dann auch tats\"achlich umsetzt?}\\
    (`How likely is it that you will actually carry out the recommended
    action?')\\
    Endpoints: 0 = \textit{sehr unwahrscheinlich}
    (`very unlikely') ---
    100 = \textit{sehr wahrscheinlich} (`very likely').
  \end{quote}
\end{enumerate}

\noindent Perceived controversy is deliberately \emph{not} included as a
rating probe: \textsc{pc\textsubscript{prag}} is experimentally
manipulated via the context paragraph, so a controversy rating would
be redundant and risks obscuring the marker's contribution to inferred
goal strength.

Recorded per trial: \texttt{trial\_id}, \texttt{topic},
\texttt{is\_filler}, \texttt{condition\_index}, \texttt{marker},
\texttt{pc\_prag}, \texttt{g\_implied}, \texttt{rating\_g},
\texttt{rating\_adopt}, \texttt{rt} (ms); and once per participant:
\texttt{list\_num}.
The continuous \texttt{pc\_prop\_rating} covariate (topic-level norming
mean) is joined to the trial data during analysis.

\subsection{Procedure}

Participants complete the experiment online via the magpie platform
\citep{kochetkova2024magpie} in a browser.
After reading instructions in German, participants are assigned to one of
4 counterbalancing lists (via URL \texttt{?list=N} parameter assigned by
Prolific; no random fallback in deployment).
They then proceed through 16 trials---8 critical and 8 fillers---in a
fully randomised order, followed by a post-test questionnaire eliciting
native language, age, gender, education level, and optional comments.

Each trial displays:
(a) a context paragraph establishing the communicative situation and the
\textsc{pc\textsubscript{prag}} level;
(b) the target sentence with the best-choice marker already in place in
the Mittelfeld;
(c) the consequent action clause;
(d) the two rating sliders in fixed order ($\widehat{g}$, $\widehat{A}$),
each labelled in German.
The ``Weiter'' (Next) button is disabled until both sliders have been
moved at least once, ensuring a complete response on every trial.

\subsection{Planned Analysis}

\subsubsection{Empirical plots}

For each dependent variable ($\widehat{g}$, $\widehat{A}$):

\begin{itemize}
  \item \textbf{Condition means}: a dot-and-whisker plot shows mean
        ratings per condition (the $2\times2$ of pc\textsubscript{prag}
        $\times$ $g_{\mathrm{implied}}$) with 95\% bootstrapped CIs.
        The expected pattern: $\widehat{g}$ and $\widehat{A}$ should
        be higher in the high-$g_{\mathrm{implied}}$ conditions
        (where the model assigns a stronger marker), and lower in the
        high-pc\textsubscript{prag} conditions (credibility discounting).
  \item \textbf{Marker identity}: because the best-choice marker may
        differ across items within the same condition, a secondary plot
        shows ratings grouped by actual marker identity, allowing
        inspection of whether marker type (rather than just condition)
        drives the effects.
  \item \textbf{pc\textsubscript{prop} covariate}: scatter plots display
        each DV as a function of the continuous norming rating per
        condition, testing whether higher propositional acceptance
        amplifies the $\widehat{g}$ and $\widehat{A}$ signals for
        strong markers.
\end{itemize}

\subsubsection{Bayesian mixed-effects regression}

Two Bayesian mixed-effects regression models are fitted via
\texttt{brms} \citep{burkner2017}, one per DV ($\widehat{g}$ and
$\widehat{A}$):

\begin{align}
y &\sim \mathrm{Beta}(\mu, \phi)
\quad \text{(ratings rescaled to }(0,1)\text{)} \\
\mathrm{logit}(\mu) &= \alpha
+ \beta_{g} \cdot \mathrm{g\_implied\_c}
+ \beta_{\mathrm{prag}} \cdot \mathrm{pc\_prag\_c}
+ \beta_{\mathrm{prop}} \cdot \mathrm{pc\_prop\_c}
+ \beta_{\mathrm{int}} \cdot \mathrm{g\_implied\_c}
  \times \mathrm{pc\_prag\_c}
+ (1 \mid \mathrm{submission\_id})
+ (1 \mid \mathrm{topic}),
\end{align}

where \texttt{g\_implied\_c} is the $g$ level from the corresponding
production condition ($-0.5$ = low, $+0.5$ = high), effectively
coding the marker's position on the production utility scale.
A secondary specification replaces \texttt{g\_implied\_c} with
actualy marker identity as an unordered factor, to examine
non-monotonic patterns and verify that marker rather than
condition label drives effects.

Model comparison via approximate leave-one-out cross-validation (LOO)
selects among additive, interaction, and marker-as-categorical
formulations.

\subsubsection{RSA listener parameter estimation}
\label{subsub:rsa_fit}

The two empirical rating means (per condition cell: 4 conditions) constrain
the listener parameters of the noisy-threshold RSA model.
Specifically, the model predicts:

\begin{align}
\mathbb{E}[\widehat{g} \mid u, \text{context}]
  &= \int g \cdot P(g \mid u, pc_{\mathrm{prag}})\, dg,
\label{eq:pred_g} \\
P(\widehat{A} \mid u, \text{context})
  &= \iiint \sigma(\eta_0 + \eta_g g - \eta_{pc}\, pc
     - \eta_{\mathrm{int}}\, g \cdot pc)\;
     P(pc_{\mathrm{prop}}, pc_{\mathrm{prag}}, g \mid u)\,
     d{pc_{\mathrm{prop}}}\,d{pc_{\mathrm{prag}}}\,dg,
\label{eq:adoption}
\end{align}

\noindent where the posterior $P(pc_{\mathrm{prop}}, pc_{\mathrm{prag}}, g
\mid u)$ is conditioned on the context's \textsc{pc\textsubscript{prag}}
level as a known prior (since it is explicitly given to the listener).

The threshold parameters $(\mu_1,\mu_2, \sigma_1,\sigma_2,
\lambda)$ together with the marker-specific costs $(c_u)$ are estimated
jointly from the production experiment data via the Stan model.
Given those threshold estimates, the listener parameters
$(\eta_0, \eta_g, \eta_{pc}, \eta_{\mathrm{int}})$ are then estimated by
minimising the weighted squared deviation between the model-predicted
means (Eqs.~\ref{eq:pred_g}--\ref{eq:adoption}) and the empirical rating
means from the comprehension experiment, or via a full joint Bayesian
fit over both experiments simultaneously.

\subsection{Model Predictions}

The noisy-threshold RSA model makes the following predictions for
$\widehat{g}$ and $\widehat{A}$, which serve as pre-registered hypotheses.
In all cases, ``strong-marker condition'' refers to conditions 1 and 3
(high $g_{\mathrm{implied}}$, i.e.\ the marker selected by the production
model under a high-$g$ context), and ``weak-marker condition'' to conditions
0 and 2.

\begin{enumerate}
  \item \textbf{Marker encodes inferred goal $\widehat{g}$}: $\widehat{g}$
        ratings should be higher in strong-marker conditions than in
        weak-marker conditions, across both pc\textsubscript{prag} levels.
        This is the core inversion prediction: if the production model
        correctly captures speaker behaviour, listeners should be able to
        decode $g$ from the marker.

  \item \textbf{Marker encodes adoption $\widehat{A}$}: $\widehat{A}$
        ratings should likewise be higher in strong-marker conditions,
        reflecting that higher inferred $g$ translates into greater
        adoption pressure downstream.

  \item \textbf{pc\textsubscript{prag} modulates $\widehat{A}$ but not
        $\widehat{g}$}: because $g$ is orthogonal to pc\textsubscript{prag}
        in the model's utility function, the effect of the marker on
        $\widehat{g}$ should be largely the same in low- and high-
        pc\textsubscript{prag} contexts.
        By contrast, $\widehat{A}$ should be reduced in high-
        pc\textsubscript{prag} contexts via the interaction term
        $\eta_{\mathrm{int}}\,g \cdot pc$: even a strong-marker signal
        for high $g$ is partly cancelled when controversy is high.

  \item \textbf{Credibility discounting}: for strong-marker conditions
        presented in a high-pc\textsubscript{prag} context (condition~3),
        $\widehat{A}$ should be \emph{lower} than for the same marker in
        a low-pc\textsubscript{prag} context (condition~1).
        This is the model's signature prediction distinguishing Bayesian
        inversion from a simpler marker-strength heuristic: a sceptical
        listener who knows controversy is high will discount the adoption
        implication of even the strongest marker.
        The $g_{\mathrm{implied}} \times \textsc{pc\textsubscript{prag}}$
        interaction on $\widehat{A}$ provides the empirical test.
\end{enumerate}

\end{document}
.
%
% Compile stand-alone:
%   pdflatex comprehension_exp.tex
%
% Embed in main.tex:
%   % comprehension_exp.tex
% Stand-alone plan for the comprehension experiment.
% Can be embedded in main.tex via \input{comprehension_exp}.
%
% Compile stand-alone:
%   pdflatex comprehension_exp.tex
%
% Embed in main.tex:
%   \input{comprehension_exp}   % (inside a \section{Experiment} environment)

\documentclass[12pt]{article}
\usepackage{CSPArticleStyle}
\usepackage{csquotes}
\usepackage[margin=1in]{geometry}
\usepackage{booktabs}
\usepackage{parskip}
\usepackage{amsmath}
\usepackage{times}
\usepackage[T1]{fontenc}
\usepackage[utf8]{inputenc}
\usepackage{xcolor}
\usepackage{tabularray}
\usepackage{gb4e}
\noautomath

\title{Comprehension Experiment: Inferring Speaker Goal and Action Adoption
       from Consensus-Marker Use}
\author{}
\date{}

\begin{document}

%% ============================================================
%% When embedding in main.tex, delete from \documentclass to
%% \begin{document} and delete \end{document} at the bottom.
%% ============================================================

\section{Experiment: Consensus-Marker Comprehension}

\subsection{Overview}

The comprehension experiment is the listener-side counterpart of the
production experiment and tests whether listeners can invert the
speaker's marker choice to recover the latent context that motivated it.
Where the production experiment gave participants a context (fixing
\textsc{pc\textsubscript{prag}} and $g$) and asked which marker they
would use, the comprehension experiment gives participants a context
(fixing \textsc{pc\textsubscript{prag}}) together with a marker already
embedded in the sentence and asks what goal strength the speaker likely
had and how likely the listener is to adopt $q$.

The specific marker presented in each vignette is not chosen arbitrarily
from the three candidate markers. Our preregistered plan was to assign,
for each item$\times$condition cell, the \emph{Bayesian best-choice marker}:
the marker with the highest posterior probability under the cost-augmented
noisy-threshold RSA model fitted to the production data, given the context's
\textsc{pc\textsubscript{prag}} level and a specified $g$ level. This was
intended to mirror the production design closely and to ensure that the
comprehension stimuli reflect model-based marker use rather than arbitrary
marker$\times$item pairings.

Applying this strict modal-assignment rule to the fitted production model,
however, yielded a highly imbalanced stimulus set: most critical cells were
assigned the strongest marker, with very little variation across marker
categories. In addition, many of these modal decisions were supported only by
very small probability margins over the next-best alternative. Because such a
stimulus set would make the comprehension experiment substantially less
informative about marker-based listener inferences, we introduced a documented,
exploratory balancing step for stimulus construction. Starting from the same
posterior predictive probabilities, we selected a near-balanced assignment of
the three markers across the 32 critical cells while maximizing total posterior
support under the fitted production model. The strict modal assignment is
retained as the preregistered baseline, but the balanced assignment is used for
the actual comprehension materials.

\textsc{pc\textsubscript{prag}} is retained as an explicit context
manipulation (low vs.\ high), so that listeners can be expected to
have a well-defined prior on pragmatic controversy when interpreting
the marker. \textsc{pc\textsubscript{prop}} is not manipulated within
the experiment; it enters as a continuous topic-level covariate from
the norming study.

After reading each vignette, participants provide two slider ratings
operationalising the key listener-side quantities:
\begin{enumerate}
  \item \textbf{Inferred speaker goal strength} ($\widehat{g}$):
        how strongly does the speaker seem to want the listener to
        accept and act on $q$?
  \item \textbf{Adoption likelihood} ($\widehat{A}$):
        how likely is the listener to follow the recommended action $q$?
\end{enumerate}

\noindent Perceived controversy is \emph{not} elicited as a rating:
because \textsc{pc\textsubscript{prag}} is already manipulated as a
context factor, asking participants to rate controversy would be
redundant with respect to the manipulation check and circular with
respect to the research question. The theoretically central inference
to test is whether listeners can decode the speaker's \emph{goal} from
marker use, and whether that inference translates into adoption
behaviour.

The two slider measures, together with the \textsc{pc\textsubscript{prag}}
context and the \textsc{pc\textsubscript{prop}} covariate, provide the
data needed to fit the model's listener parameters
$(\eta_0, \eta_g, \eta_{pc}, \eta_{\mathrm{int}})$.

\subsection{Participants}

Participants will be recruited via Prolific and will be required to be
native speakers of German (language-at-home filter).
Data collection is planned after completion of the production experiment
(model fitting required to determine the best-choice markers).
A target of $N = 80$ participants is planned ($n = 20$ per list across
4 counterbalancing lists), each completing 16 trials (8 critical + 8
filler), matching the structure of the production experiment.

\subsection{Materials}

\subsubsection{Markers}

Each critical vignette is presented with one specific consensus marker
already embedded in the Mittelfeld of the sentence frame (see
example~(\ref{ex:frame_comp})):

\begin{exe}
  \ex \label{ex:frame_comp}
  Digitale Kompetenzen sind \textbf{[MARKER]}
  in der heutigen Arbeitswelt unverzichtbar.\\
  `Digital skills are \textbf{[MARKER]} indispensable in today's
  working world.'
\end{exe}

Unlike the production experiment, participants do \emph{not} choose the
marker; they read the sentence with the marker already in place and respond
to the two rating probes.

Critically, the marker assigned to each condition is not drawn arbitrarily
from the three candidates. Under the preregistered plan, for each of the four
production conditions (pc\textsubscript{prag} $\times$ $g$), the marker used in
comprehension was to be the \textbf{Bayesian best-choice marker}: the marker
$u^*$ with the highest posterior probability under the fitted
cost-augmented noisy-threshold RSA production model,
\[
u^* = \arg\max_{u} \; P(u \mid pc_{\mathrm{prag}}, g,
                           pc_{\mathrm{prop}\_c}).
\]
Because $pc_{\mathrm{prop}}$ varies continuously across items,
$u^*$ may differ per item even within the same
(pc\textsubscript{prag}, $g$) cell; the final stimulus list is therefore
constructed item-by-item once the production model has been fitted.
In practice, the fitted production model yielded a nearly degenerate modal map,
with most cells assigned to \textit{bekanntlich}. We therefore retained the
model-based logic of the procedure but replaced strict cellwise modal selection
with an uncertainty-aware balancing step: among the marker assignments with
high posterior support, we selected a near-balanced configuration across the 32
critical cells. This preserves the production--comprehension link while also
ensuring that the comprehension experiment contains sufficient marker variation
to test listener-side inferences.

All three markers remain candidates; the mapping from conditions to
markers is an empirical output of the production analysis, not a design
pre-commitment. In the final comprehension materials, we use the
\textbf{balanced model-based assignment}, rather than the strict modal
assignment, in order to maintain informative variation across marker types.

\subsubsection{Topic Items}

The same eight topic domains used in the norming study and the production
experiment are used here (Table~\ref{tab:items_comp}).
The 4-condition $\times$ 8-item Latin square (Section~\ref{subsub:counterbalancing_comp})
requires no additional items.

\begin{table}[ht]
\centering
\caption{Topic items: sentence frame for $p$ (with marker position [M])
         and consequent $q$.}
\label{tab:items_comp}
\small
\begin{tabular}{clll}
\toprule
\textbf{\#} & \textbf{Topic} & \textbf{Sentence frame} & \textbf{Consequent $q$} \\
\midrule
1 & Klimawandel &
  \textit{Der Klimawandel wird [M]} &
  auf Flugreisen verzichten \\
  & & \textit{durch menschliche Aktivit\"aten verursacht.} & \\[3pt]
2 & Ern\"ahrung &
  \textit{Eine \"uberwiegend pflanzliche Ern\"ahrung} &
  mehr pflanzliche Lebens- \\
  & & \textit{f\"ordert [M] die Gesundheit.} &
  mittel einbauen \\[3pt]
3 & Stadtverkehr &
  \textit{Ein ausgebautes Nahverkehrsnetz} &
  h\"aufiger auf \\
  & & \textit{entlastet [M] den Stadtverkehr} &
  \"offentliche Verkehrs- \\
  & & \textit{deutlich.} &
  mittel umsteigen \\[3pt]
4 & Digitalisierung &
  \textit{Digitale Kompetenzen sind [M] in} &
  regelm\"a\ss{}ig Weiter- \\
  & & \textit{der heutigen Arbeitswelt unverzichtbar.} &
  bildungsangebote nutzen \\[3pt]
5 & Schlaf &
  \textit{Ausreichend Schlaf ist [M] entscheidend} &
  auf feste Schlafzeiten \\
  & & \textit{f\"ur die kognitive Leistungsf\"ahigkeit.} &
  achten \\[3pt]
6 & Plastik &
  \textit{Einwegplastik schadet [M] den} &
  konsequent auf Einweg- \\
  & & \textit{Meeres\"okosystemen erheblich.} &
  plastik verzichten \\[3pt]
7 & Sport &
  \textit{Regelm\"a\ss{}ige k\"orperliche Bewegung} &
  regelm\"a\ss{}ig Sport in \\
  & & \textit{senkt [M] das Risiko f\"ur} &
  den Alltag integrieren \\
  & & \textit{Herz-Kreislauf-Erkrankungen deutlich.} & \\[3pt]
8 & Lokalkauf &
  \textit{Das Kaufen bei lokalen H\"andlern} &
  \"ofter bei lokalen \\
  & & \textit{st\"arkt [M] die regionale Wirtschaft} &
  Gesch\"aften einkaufen \\
  & & \textit{nachhaltig.} & \\
\bottomrule
\end{tabular}
\end{table}

\noindent [M] denotes the position of the pre-filled best-choice marker.

\subsubsection{Vignettes}

\paragraph{Structure.}
Each trial presents a vignette consisting of three parts:
(a) a context paragraph establishing the communicative situation and
fixing the \textsc{pc\textsubscript{prag}} level;
(b) the target sentence with the best-choice marker already embedded
in the Mittelfeld;
(c) the consequent action clause.
The two rating probes follow immediately below the vignette.

\paragraph{Propositional perceived controversy (\textsc{pc\textsubscript{prop}}).}
As in the production experiment, \textsc{pc\textsubscript{prop}} is not
manipulated within the comprehension experiment.
Topic-level propositional acceptance ratings from the norming study
(0--100 scale; e.g., \textit{Digitalisierung}: 85/100,
\textit{Lokalkauf}: 58/100) serve as a continuous, mean-centred
covariate \textsc{pc\textsubscript{prop}}$_c$ in analysis.
Context paragraphs present $p$ without editorial evaluation; variation in
listeners' individual assessments of propositional controversy is
captured by the norming-derived covariate.

\paragraph{Pragmatic perceived controversy (\textsc{pc\textsubscript{prag}}).}
\textsc{pc\textsubscript{prag}} is manipulated between item instances
using the same two context versions as in the production experiment:
the speaker's social circle either \emph{largely agrees} (low) or
\emph{holds divergent views} (high) about $p$.
For example, for the \textit{Digitalisierung} item:

\begin{itemize}
  \item $\textsc{pc}_\text{prag} = \text{low}$: \textit{Deine Kollegin spricht
        mit dir \"uber Berufsausbildung und Weiterbildung.
        In eurem Team teilen die meisten die Ansicht, dass digitale Kompetenzen
        in der heutigen Arbeitswelt unverzichtbar sind.}
  \item $\textsc{pc}_\text{prag} = \text{high}$: \textit{Deine Kollegin spricht
        mit dir \"uber Berufsausbildung und Weiterbildung.
        In eurem Team gehen die Meinungen dazu stark auseinander.}
\end{itemize}

The context is framed from the listener's perspective as a direct addressee,
mirroring the production study where participants adopt the speaker's
perspective.

\paragraph{Marker manipulation.}
The dependent measures of the production experiment (which marker is chosen)
become the independent variable here.
One of the three consensus markers is inserted into the Mittelfeld of the
target sentence.
A complete sample item for the \textit{Digitalisierung} topic
($\textsc{pc}_\text{prag} = \text{high}$, marker = \textit{bekanntlich};
norming $\textsc{pc}_\text{prop}$ rating: 85/100) is given in
(\ref{ex:comp_sample}):

\begin{exe}
  \ex \label{ex:comp_sample}
  \textit{Deine Kollegin spricht mit dir \"uber Berufsausbildung und
  Weiterbildung.
  In eurem Team gehen die Meinungen dazu stark auseinander.}

  \textit{Sie sagt zu dir:} Digitale Kompetenzen sind \textbf{bekanntlich}
  in der heutigen Arbeitswelt unverzichtbar.
  Deshalb solltest du regelm\"a\ss{}ig digitale Weiterbildungsangebote nutzen.

  \noindent\textit{[Rating probes follow.]}
\end{exe}

\subsection{Design}

\subsubsection{Factorial structure}

The comprehension experiment mirrors the production experiment's
$\mathbf{2}$~\textbf{(pc\textsubscript{prag}: low vs.\ high)}
$\mathbf{\times}$~$\mathbf{2}$~\textbf{($g_{\mathrm{implied}}$: low
vs.\ high) design}, yielding \textbf{4 conditions}.
The $g_{\mathrm{implied}}$ dimension is not manipulated via an explicit
goal instruction (as in production) but is instead conveyed entirely
through the choice of marker: for each condition, the marker presented
is the best-choice marker for the corresponding
(\textsc{pc\textsubscript{prag}},~$g$) cell from the production model.
\textsc{pc\textsubscript{prop}} enters as a continuous norming covariate.
With 8 topic items and a Latin square rotation across 4 lists, each
participant sees exactly \textbf{8 critical trials}---one item per
condition, each condition appearing twice---plus 8 filler trials,
for \textbf{16 trials} in total (Table~\ref{tab:design_comp}).

\begin{table}[ht]
\centering
\caption{%
  The $2\times2$ design (4 conditions), mirroring the production experiment.
  The marker presented in each condition is the Bayesian best-choice
  marker from the fitted production model; specific markers are
  determined after production data collection.}
\label{tab:design_comp}
\small
\begin{tabular}{ccll}
\toprule
\textbf{pc\textsubscript{prag}} &
\textbf{$g_{\mathrm{implied}}$} &
\textbf{Marker (TBD from production)} &
\textbf{Condition label} \\
\midrule
low  & low  & best-choice for $(\text{low}, g_{\text{low}})$  &
  (0) community agrees, weak signal \\
low  & high & best-choice for $(\text{low}, g_{\text{high}})$ &
  (1) community agrees, strong signal \\
high & low  & best-choice for $(\text{high}, g_{\text{low}})$ &
  (2) community divided, weak signal \\
high & high & best-choice for $(\text{high}, g_{\text{high}})$ &
  (3) community divided, strong signal \\
\bottomrule
\end{tabular}
\end{table}

\subsubsection{Counterbalancing}
\label{subsub:counterbalancing_comp}

The experiment uses the \textbf{same $4 \times 8$ Latin square} as the
production experiment: each of the 4 lists assigns item $i$
($i = 0,\ldots,7$) to condition $(i + l)\bmod 4$, where
$l \in \{0,1,2,3\}$ is the list number.
This guarantees that:
\begin{itemize}
  \item each participant sees every item exactly once;
  \item each participant sees each of the 4 conditions exactly twice
        (two items per condition), ensuring balanced coverage of the
        $2\times2$ design;
  \item across all 4 lists, each item appears equally often in each
        condition.
\end{itemize}

List assignment follows the same Prolific-and-URL strategy as the production
experiment: each Prolific study embeds \texttt{?list=N} (0--3) in the
experiment URL, capping recruitment at $n = 20$ participants per list
for a total of $N = 80$.

In addition to the 8 critical trials, each participant completed 8
\textbf{filler trials} using the same filler items as the production
experiment (train travel, reading light, cooking, coffee, vacation,
hydration, fresh air, noise), each presented with a randomly assigned
marker.
Fillers are excluded from analysis; they serve to dilute the critical
items and reduce demand characteristics.

\subsubsection{Dependent variables}
\label{subsub:dvs_comp}

Two slider ratings are collected after each vignette (0--100 continuous
scale with no visible numbers, labelled at the endpoints as shown):

\begin{enumerate}
  \item \textbf{Inferred speaker goal strength} (operationalises $\widehat{g}$):
  \begin{quote}
    \textit{Wie stark hast du den Eindruck, dass die Sprecherin m\"ochte,
    dass du ihre Empfehlung befolgst?}\\
    (`How strongly do you get the impression that the speaker wants you to
    follow their recommendation?')\\
    Endpoints: 0 = \textit{\"uberhaupt nicht}
    (`not at all') ---
    100 = \textit{sehr stark} (`very strongly').
  \end{quote}

  \item \textbf{Adoption likelihood of $q$} (operationalises $\widehat{A}$):
  \begin{quote}
    \textit{Wie wahrscheinlich ist es, dass du die empfohlene Handlung
    dann auch tats\"achlich umsetzt?}\\
    (`How likely is it that you will actually carry out the recommended
    action?')\\
    Endpoints: 0 = \textit{sehr unwahrscheinlich}
    (`very unlikely') ---
    100 = \textit{sehr wahrscheinlich} (`very likely').
  \end{quote}
\end{enumerate}

\noindent Perceived controversy is deliberately \emph{not} included as a
rating probe: \textsc{pc\textsubscript{prag}} is experimentally
manipulated via the context paragraph, so a controversy rating would
be redundant and risks obscuring the marker's contribution to inferred
goal strength.

Recorded per trial: \texttt{trial\_id}, \texttt{topic},
\texttt{is\_filler}, \texttt{condition\_index}, \texttt{marker},
\texttt{pc\_prag}, \texttt{g\_implied}, \texttt{rating\_g},
\texttt{rating\_adopt}, \texttt{rt} (ms); and once per participant:
\texttt{list\_num}.
The continuous \texttt{pc\_prop\_rating} covariate (topic-level norming
mean) is joined to the trial data during analysis.

\subsection{Procedure}

Participants complete the experiment online via the magpie platform
\citep{kochetkova2024magpie} in a browser.
After reading instructions in German, participants are assigned to one of
4 counterbalancing lists (via URL \texttt{?list=N} parameter assigned by
Prolific; no random fallback in deployment).
They then proceed through 16 trials---8 critical and 8 fillers---in a
fully randomised order, followed by a post-test questionnaire eliciting
native language, age, gender, education level, and optional comments.

Each trial displays:
(a) a context paragraph establishing the communicative situation and the
\textsc{pc\textsubscript{prag}} level;
(b) the target sentence with the best-choice marker already in place in
the Mittelfeld;
(c) the consequent action clause;
(d) the two rating sliders in fixed order ($\widehat{g}$, $\widehat{A}$),
each labelled in German.
The ``Weiter'' (Next) button is disabled until both sliders have been
moved at least once, ensuring a complete response on every trial.

\subsection{Planned Analysis}

\subsubsection{Empirical plots}

For each dependent variable ($\widehat{g}$, $\widehat{A}$):

\begin{itemize}
  \item \textbf{Condition means}: a dot-and-whisker plot shows mean
        ratings per condition (the $2\times2$ of pc\textsubscript{prag}
        $\times$ $g_{\mathrm{implied}}$) with 95\% bootstrapped CIs.
        The expected pattern: $\widehat{g}$ and $\widehat{A}$ should
        be higher in the high-$g_{\mathrm{implied}}$ conditions
        (where the model assigns a stronger marker), and lower in the
        high-pc\textsubscript{prag} conditions (credibility discounting).
  \item \textbf{Marker identity}: because the best-choice marker may
        differ across items within the same condition, a secondary plot
        shows ratings grouped by actual marker identity, allowing
        inspection of whether marker type (rather than just condition)
        drives the effects.
  \item \textbf{pc\textsubscript{prop} covariate}: scatter plots display
        each DV as a function of the continuous norming rating per
        condition, testing whether higher propositional acceptance
        amplifies the $\widehat{g}$ and $\widehat{A}$ signals for
        strong markers.
\end{itemize}

\subsubsection{Bayesian mixed-effects regression}

Two Bayesian mixed-effects regression models are fitted via
\texttt{brms} \citep{burkner2017}, one per DV ($\widehat{g}$ and
$\widehat{A}$):

\begin{align}
y &\sim \mathrm{Beta}(\mu, \phi)
\quad \text{(ratings rescaled to }(0,1)\text{)} \\
\mathrm{logit}(\mu) &= \alpha
+ \beta_{g} \cdot \mathrm{g\_implied\_c}
+ \beta_{\mathrm{prag}} \cdot \mathrm{pc\_prag\_c}
+ \beta_{\mathrm{prop}} \cdot \mathrm{pc\_prop\_c}
+ \beta_{\mathrm{int}} \cdot \mathrm{g\_implied\_c}
  \times \mathrm{pc\_prag\_c}
+ (1 \mid \mathrm{submission\_id})
+ (1 \mid \mathrm{topic}),
\end{align}

where \texttt{g\_implied\_c} is the $g$ level from the corresponding
production condition ($-0.5$ = low, $+0.5$ = high), effectively
coding the marker's position on the production utility scale.
A secondary specification replaces \texttt{g\_implied\_c} with
actualy marker identity as an unordered factor, to examine
non-monotonic patterns and verify that marker rather than
condition label drives effects.

Model comparison via approximate leave-one-out cross-validation (LOO)
selects among additive, interaction, and marker-as-categorical
formulations.

\subsubsection{RSA listener parameter estimation}
\label{subsub:rsa_fit}

The two empirical rating means (per condition cell: 4 conditions) constrain
the listener parameters of the noisy-threshold RSA model.
Specifically, the model predicts:

\begin{align}
\mathbb{E}[\widehat{g} \mid u, \text{context}]
  &= \int g \cdot P(g \mid u, pc_{\mathrm{prag}})\, dg,
\label{eq:pred_g} \\
P(\widehat{A} \mid u, \text{context})
  &= \iiint \sigma(\eta_0 + \eta_g g - \eta_{pc}\, pc
     - \eta_{\mathrm{int}}\, g \cdot pc)\;
     P(pc_{\mathrm{prop}}, pc_{\mathrm{prag}}, g \mid u)\,
     d{pc_{\mathrm{prop}}}\,d{pc_{\mathrm{prag}}}\,dg,
\label{eq:adoption}
\end{align}

\noindent where the posterior $P(pc_{\mathrm{prop}}, pc_{\mathrm{prag}}, g
\mid u)$ is conditioned on the context's \textsc{pc\textsubscript{prag}}
level as a known prior (since it is explicitly given to the listener).

The threshold parameters $(\mu_1,\mu_2, \sigma_1,\sigma_2,
\lambda)$ together with the marker-specific costs $(c_u)$ are estimated
jointly from the production experiment data via the Stan model.
Given those threshold estimates, the listener parameters
$(\eta_0, \eta_g, \eta_{pc}, \eta_{\mathrm{int}})$ are then estimated by
minimising the weighted squared deviation between the model-predicted
means (Eqs.~\ref{eq:pred_g}--\ref{eq:adoption}) and the empirical rating
means from the comprehension experiment, or via a full joint Bayesian
fit over both experiments simultaneously.

\subsection{Model Predictions}

The noisy-threshold RSA model makes the following predictions for
$\widehat{g}$ and $\widehat{A}$, which serve as pre-registered hypotheses.
In all cases, ``strong-marker condition'' refers to conditions 1 and 3
(high $g_{\mathrm{implied}}$, i.e.\ the marker selected by the production
model under a high-$g$ context), and ``weak-marker condition'' to conditions
0 and 2.

\begin{enumerate}
  \item \textbf{Marker encodes inferred goal $\widehat{g}$}: $\widehat{g}$
        ratings should be higher in strong-marker conditions than in
        weak-marker conditions, across both pc\textsubscript{prag} levels.
        This is the core inversion prediction: if the production model
        correctly captures speaker behaviour, listeners should be able to
        decode $g$ from the marker.

  \item \textbf{Marker encodes adoption $\widehat{A}$}: $\widehat{A}$
        ratings should likewise be higher in strong-marker conditions,
        reflecting that higher inferred $g$ translates into greater
        adoption pressure downstream.

  \item \textbf{pc\textsubscript{prag} modulates $\widehat{A}$ but not
        $\widehat{g}$}: because $g$ is orthogonal to pc\textsubscript{prag}
        in the model's utility function, the effect of the marker on
        $\widehat{g}$ should be largely the same in low- and high-
        pc\textsubscript{prag} contexts.
        By contrast, $\widehat{A}$ should be reduced in high-
        pc\textsubscript{prag} contexts via the interaction term
        $\eta_{\mathrm{int}}\,g \cdot pc$: even a strong-marker signal
        for high $g$ is partly cancelled when controversy is high.

  \item \textbf{Credibility discounting}: for strong-marker conditions
        presented in a high-pc\textsubscript{prag} context (condition~3),
        $\widehat{A}$ should be \emph{lower} than for the same marker in
        a low-pc\textsubscript{prag} context (condition~1).
        This is the model's signature prediction distinguishing Bayesian
        inversion from a simpler marker-strength heuristic: a sceptical
        listener who knows controversy is high will discount the adoption
        implication of even the strongest marker.
        The $g_{\mathrm{implied}} \times \textsc{pc\textsubscript{prag}}$
        interaction on $\widehat{A}$ provides the empirical test.
\end{enumerate}

\end{document}
   % (inside a \section{Experiment} environment)

\documentclass[12pt]{article}
\usepackage{CSPArticleStyle}
\usepackage{csquotes}
\usepackage[margin=1in]{geometry}
\usepackage{booktabs}
\usepackage{parskip}
\usepackage{amsmath}
\usepackage{times}
\usepackage[T1]{fontenc}
\usepackage[utf8]{inputenc}
\usepackage{xcolor}
\usepackage{tabularray}
\usepackage{gb4e}
\noautomath

\title{Comprehension Experiment: Inferring Speaker Goal and Action Adoption
       from Consensus-Marker Use}
\author{}
\date{}

\begin{document}

%% ============================================================
%% When embedding in main.tex, delete from \documentclass to
%% \begin{document} and delete \end{document} at the bottom.
%% ============================================================

\section{Experiment: Consensus-Marker Comprehension}

\subsection{Overview}

The comprehension experiment is the listener-side counterpart of the
production experiment and tests whether listeners can invert the
speaker's marker choice to recover the latent context that motivated it.
Where the production experiment gave participants a context (fixing
\textsc{pc\textsubscript{prag}} and $g$) and asked which marker they
would use, the comprehension experiment gives participants a context
(fixing \textsc{pc\textsubscript{prag}}) together with a marker already
embedded in the sentence and asks what goal strength the speaker likely
had and how likely the listener is to adopt $q$.

The specific marker presented in each vignette is not chosen arbitrarily
from the three candidate markers. Our preregistered plan was to assign,
for each item$\times$condition cell, the \emph{Bayesian best-choice marker}:
the marker with the highest posterior probability under the cost-augmented
noisy-threshold RSA model fitted to the production data, given the context's
\textsc{pc\textsubscript{prag}} level and a specified $g$ level. This was
intended to mirror the production design closely and to ensure that the
comprehension stimuli reflect model-based marker use rather than arbitrary
marker$\times$item pairings.

Applying this strict modal-assignment rule to the fitted production model,
however, yielded a highly imbalanced stimulus set: most critical cells were
assigned the strongest marker, with very little variation across marker
categories. In addition, many of these modal decisions were supported only by
very small probability margins over the next-best alternative. Because such a
stimulus set would make the comprehension experiment substantially less
informative about marker-based listener inferences, we introduced a documented,
exploratory balancing step for stimulus construction. Starting from the same
posterior predictive probabilities, we selected a near-balanced assignment of
the three markers across the 32 critical cells while maximizing total posterior
support under the fitted production model. The strict modal assignment is
retained as the preregistered baseline, but the balanced assignment is used for
the actual comprehension materials.

\textsc{pc\textsubscript{prag}} is retained as an explicit context
manipulation (low vs.\ high), so that listeners can be expected to
have a well-defined prior on pragmatic controversy when interpreting
the marker. \textsc{pc\textsubscript{prop}} is not manipulated within
the experiment; it enters as a continuous topic-level covariate from
the norming study.

After reading each vignette, participants provide two slider ratings
operationalising the key listener-side quantities:
\begin{enumerate}
  \item \textbf{Inferred speaker goal strength} ($\widehat{g}$):
        how strongly does the speaker seem to want the listener to
        accept and act on $q$?
  \item \textbf{Adoption likelihood} ($\widehat{A}$):
        how likely is the listener to follow the recommended action $q$?
\end{enumerate}

\noindent Perceived controversy is \emph{not} elicited as a rating:
because \textsc{pc\textsubscript{prag}} is already manipulated as a
context factor, asking participants to rate controversy would be
redundant with respect to the manipulation check and circular with
respect to the research question. The theoretically central inference
to test is whether listeners can decode the speaker's \emph{goal} from
marker use, and whether that inference translates into adoption
behaviour.

The two slider measures, together with the \textsc{pc\textsubscript{prag}}
context and the \textsc{pc\textsubscript{prop}} covariate, provide the
data needed to fit the model's listener parameters
$(\eta_0, \eta_g, \eta_{pc}, \eta_{\mathrm{int}})$.

\subsection{Participants}

Participants will be recruited via Prolific and will be required to be
native speakers of German (language-at-home filter).
Data collection is planned after completion of the production experiment
(model fitting required to determine the best-choice markers).
A target of $N = 80$ participants is planned ($n = 20$ per list across
4 counterbalancing lists), each completing 16 trials (8 critical + 8
filler), matching the structure of the production experiment.

\subsection{Materials}

\subsubsection{Markers}

Each critical vignette is presented with one specific consensus marker
already embedded in the Mittelfeld of the sentence frame (see
example~(\ref{ex:frame_comp})):

\begin{exe}
  \ex \label{ex:frame_comp}
  Digitale Kompetenzen sind \textbf{[MARKER]}
  in der heutigen Arbeitswelt unverzichtbar.\\
  `Digital skills are \textbf{[MARKER]} indispensable in today's
  working world.'
\end{exe}

Unlike the production experiment, participants do \emph{not} choose the
marker; they read the sentence with the marker already in place and respond
to the two rating probes.

Critically, the marker assigned to each condition is not drawn arbitrarily
from the three candidates. Under the preregistered plan, for each of the four
production conditions (pc\textsubscript{prag} $\times$ $g$), the marker used in
comprehension was to be the \textbf{Bayesian best-choice marker}: the marker
$u^*$ with the highest posterior probability under the fitted
cost-augmented noisy-threshold RSA production model,
\[
u^* = \arg\max_{u} \; P(u \mid pc_{\mathrm{prag}}, g,
                           pc_{\mathrm{prop}\_c}).
\]
Because $pc_{\mathrm{prop}}$ varies continuously across items,
$u^*$ may differ per item even within the same
(pc\textsubscript{prag}, $g$) cell; the final stimulus list is therefore
constructed item-by-item once the production model has been fitted.
In practice, the fitted production model yielded a nearly degenerate modal map,
with most cells assigned to \textit{bekanntlich}. We therefore retained the
model-based logic of the procedure but replaced strict cellwise modal selection
with an uncertainty-aware balancing step: among the marker assignments with
high posterior support, we selected a near-balanced configuration across the 32
critical cells. This preserves the production--comprehension link while also
ensuring that the comprehension experiment contains sufficient marker variation
to test listener-side inferences.

All three markers remain candidates; the mapping from conditions to
markers is an empirical output of the production analysis, not a design
pre-commitment. In the final comprehension materials, we use the
\textbf{balanced model-based assignment}, rather than the strict modal
assignment, in order to maintain informative variation across marker types.

\subsubsection{Topic Items}

The same eight topic domains used in the norming study and the production
experiment are used here (Table~\ref{tab:items_comp}).
The 4-condition $\times$ 8-item Latin square (Section~\ref{subsub:counterbalancing_comp})
requires no additional items.

\begin{table}[ht]
\centering
\caption{Topic items: sentence frame for $p$ (with marker position [M])
         and consequent $q$.}
\label{tab:items_comp}
\small
\begin{tabular}{clll}
\toprule
\textbf{\#} & \textbf{Topic} & \textbf{Sentence frame} & \textbf{Consequent $q$} \\
\midrule
1 & Klimawandel &
  \textit{Der Klimawandel wird [M]} &
  auf Flugreisen verzichten \\
  & & \textit{durch menschliche Aktivit\"aten verursacht.} & \\[3pt]
2 & Ern\"ahrung &
  \textit{Eine \"uberwiegend pflanzliche Ern\"ahrung} &
  mehr pflanzliche Lebens- \\
  & & \textit{f\"ordert [M] die Gesundheit.} &
  mittel einbauen \\[3pt]
3 & Stadtverkehr &
  \textit{Ein ausgebautes Nahverkehrsnetz} &
  h\"aufiger auf \\
  & & \textit{entlastet [M] den Stadtverkehr} &
  \"offentliche Verkehrs- \\
  & & \textit{deutlich.} &
  mittel umsteigen \\[3pt]
4 & Digitalisierung &
  \textit{Digitale Kompetenzen sind [M] in} &
  regelm\"a\ss{}ig Weiter- \\
  & & \textit{der heutigen Arbeitswelt unverzichtbar.} &
  bildungsangebote nutzen \\[3pt]
5 & Schlaf &
  \textit{Ausreichend Schlaf ist [M] entscheidend} &
  auf feste Schlafzeiten \\
  & & \textit{f\"ur die kognitive Leistungsf\"ahigkeit.} &
  achten \\[3pt]
6 & Plastik &
  \textit{Einwegplastik schadet [M] den} &
  konsequent auf Einweg- \\
  & & \textit{Meeres\"okosystemen erheblich.} &
  plastik verzichten \\[3pt]
7 & Sport &
  \textit{Regelm\"a\ss{}ige k\"orperliche Bewegung} &
  regelm\"a\ss{}ig Sport in \\
  & & \textit{senkt [M] das Risiko f\"ur} &
  den Alltag integrieren \\
  & & \textit{Herz-Kreislauf-Erkrankungen deutlich.} & \\[3pt]
8 & Lokalkauf &
  \textit{Das Kaufen bei lokalen H\"andlern} &
  \"ofter bei lokalen \\
  & & \textit{st\"arkt [M] die regionale Wirtschaft} &
  Gesch\"aften einkaufen \\
  & & \textit{nachhaltig.} & \\
\bottomrule
\end{tabular}
\end{table}

\noindent [M] denotes the position of the pre-filled best-choice marker.

\subsubsection{Vignettes}

\paragraph{Structure.}
Each trial presents a vignette consisting of three parts:
(a) a context paragraph establishing the communicative situation and
fixing the \textsc{pc\textsubscript{prag}} level;
(b) the target sentence with the best-choice marker already embedded
in the Mittelfeld;
(c) the consequent action clause.
The two rating probes follow immediately below the vignette.

\paragraph{Propositional perceived controversy (\textsc{pc\textsubscript{prop}}).}
As in the production experiment, \textsc{pc\textsubscript{prop}} is not
manipulated within the comprehension experiment.
Topic-level propositional acceptance ratings from the norming study
(0--100 scale; e.g., \textit{Digitalisierung}: 85/100,
\textit{Lokalkauf}: 58/100) serve as a continuous, mean-centred
covariate \textsc{pc\textsubscript{prop}}$_c$ in analysis.
Context paragraphs present $p$ without editorial evaluation; variation in
listeners' individual assessments of propositional controversy is
captured by the norming-derived covariate.

\paragraph{Pragmatic perceived controversy (\textsc{pc\textsubscript{prag}}).}
\textsc{pc\textsubscript{prag}} is manipulated between item instances
using the same two context versions as in the production experiment:
the speaker's social circle either \emph{largely agrees} (low) or
\emph{holds divergent views} (high) about $p$.
For example, for the \textit{Digitalisierung} item:

\begin{itemize}
  \item $\textsc{pc}_\text{prag} = \text{low}$: \textit{Deine Kollegin spricht
        mit dir \"uber Berufsausbildung und Weiterbildung.
        In eurem Team teilen die meisten die Ansicht, dass digitale Kompetenzen
        in der heutigen Arbeitswelt unverzichtbar sind.}
  \item $\textsc{pc}_\text{prag} = \text{high}$: \textit{Deine Kollegin spricht
        mit dir \"uber Berufsausbildung und Weiterbildung.
        In eurem Team gehen die Meinungen dazu stark auseinander.}
\end{itemize}

The context is framed from the listener's perspective as a direct addressee,
mirroring the production study where participants adopt the speaker's
perspective.

\paragraph{Marker manipulation.}
The dependent measures of the production experiment (which marker is chosen)
become the independent variable here.
One of the three consensus markers is inserted into the Mittelfeld of the
target sentence.
A complete sample item for the \textit{Digitalisierung} topic
($\textsc{pc}_\text{prag} = \text{high}$, marker = \textit{bekanntlich};
norming $\textsc{pc}_\text{prop}$ rating: 85/100) is given in
(\ref{ex:comp_sample}):

\begin{exe}
  \ex \label{ex:comp_sample}
  \textit{Deine Kollegin spricht mit dir \"uber Berufsausbildung und
  Weiterbildung.
  In eurem Team gehen die Meinungen dazu stark auseinander.}

  \textit{Sie sagt zu dir:} Digitale Kompetenzen sind \textbf{bekanntlich}
  in der heutigen Arbeitswelt unverzichtbar.
  Deshalb solltest du regelm\"a\ss{}ig digitale Weiterbildungsangebote nutzen.

  \noindent\textit{[Rating probes follow.]}
\end{exe}

\subsection{Design}

\subsubsection{Factorial structure}

The comprehension experiment mirrors the production experiment's
$\mathbf{2}$~\textbf{(pc\textsubscript{prag}: low vs.\ high)}
$\mathbf{\times}$~$\mathbf{2}$~\textbf{($g_{\mathrm{implied}}$: low
vs.\ high) design}, yielding \textbf{4 conditions}.
The $g_{\mathrm{implied}}$ dimension is not manipulated via an explicit
goal instruction (as in production) but is instead conveyed entirely
through the choice of marker: for each condition, the marker presented
is the best-choice marker for the corresponding
(\textsc{pc\textsubscript{prag}},~$g$) cell from the production model.
\textsc{pc\textsubscript{prop}} enters as a continuous norming covariate.
With 8 topic items and a Latin square rotation across 4 lists, each
participant sees exactly \textbf{8 critical trials}---one item per
condition, each condition appearing twice---plus 8 filler trials,
for \textbf{16 trials} in total (Table~\ref{tab:design_comp}).

\begin{table}[ht]
\centering
\caption{%
  The $2\times2$ design (4 conditions), mirroring the production experiment.
  The marker presented in each condition is the Bayesian best-choice
  marker from the fitted production model; specific markers are
  determined after production data collection.}
\label{tab:design_comp}
\small
\begin{tabular}{ccll}
\toprule
\textbf{pc\textsubscript{prag}} &
\textbf{$g_{\mathrm{implied}}$} &
\textbf{Marker (TBD from production)} &
\textbf{Condition label} \\
\midrule
low  & low  & best-choice for $(\text{low}, g_{\text{low}})$  &
  (0) community agrees, weak signal \\
low  & high & best-choice for $(\text{low}, g_{\text{high}})$ &
  (1) community agrees, strong signal \\
high & low  & best-choice for $(\text{high}, g_{\text{low}})$ &
  (2) community divided, weak signal \\
high & high & best-choice for $(\text{high}, g_{\text{high}})$ &
  (3) community divided, strong signal \\
\bottomrule
\end{tabular}
\end{table}

\subsubsection{Counterbalancing}
\label{subsub:counterbalancing_comp}

The experiment uses the \textbf{same $4 \times 8$ Latin square} as the
production experiment: each of the 4 lists assigns item $i$
($i = 0,\ldots,7$) to condition $(i + l)\bmod 4$, where
$l \in \{0,1,2,3\}$ is the list number.
This guarantees that:
\begin{itemize}
  \item each participant sees every item exactly once;
  \item each participant sees each of the 4 conditions exactly twice
        (two items per condition), ensuring balanced coverage of the
        $2\times2$ design;
  \item across all 4 lists, each item appears equally often in each
        condition.
\end{itemize}

List assignment follows the same Prolific-and-URL strategy as the production
experiment: each Prolific study embeds \texttt{?list=N} (0--3) in the
experiment URL, capping recruitment at $n = 20$ participants per list
for a total of $N = 80$.

In addition to the 8 critical trials, each participant completed 8
\textbf{filler trials} using the same filler items as the production
experiment (train travel, reading light, cooking, coffee, vacation,
hydration, fresh air, noise), each presented with a randomly assigned
marker.
Fillers are excluded from analysis; they serve to dilute the critical
items and reduce demand characteristics.

\subsubsection{Dependent variables}
\label{subsub:dvs_comp}

Two slider ratings are collected after each vignette (0--100 continuous
scale with no visible numbers, labelled at the endpoints as shown):

\begin{enumerate}
  \item \textbf{Inferred speaker goal strength} (operationalises $\widehat{g}$):
  \begin{quote}
    \textit{Wie stark hast du den Eindruck, dass die Sprecherin m\"ochte,
    dass du ihre Empfehlung befolgst?}\\
    (`How strongly do you get the impression that the speaker wants you to
    follow their recommendation?')\\
    Endpoints: 0 = \textit{\"uberhaupt nicht}
    (`not at all') ---
    100 = \textit{sehr stark} (`very strongly').
  \end{quote}

  \item \textbf{Adoption likelihood of $q$} (operationalises $\widehat{A}$):
  \begin{quote}
    \textit{Wie wahrscheinlich ist es, dass du die empfohlene Handlung
    dann auch tats\"achlich umsetzt?}\\
    (`How likely is it that you will actually carry out the recommended
    action?')\\
    Endpoints: 0 = \textit{sehr unwahrscheinlich}
    (`very unlikely') ---
    100 = \textit{sehr wahrscheinlich} (`very likely').
  \end{quote}
\end{enumerate}

\noindent Perceived controversy is deliberately \emph{not} included as a
rating probe: \textsc{pc\textsubscript{prag}} is experimentally
manipulated via the context paragraph, so a controversy rating would
be redundant and risks obscuring the marker's contribution to inferred
goal strength.

Recorded per trial: \texttt{trial\_id}, \texttt{topic},
\texttt{is\_filler}, \texttt{condition\_index}, \texttt{marker},
\texttt{pc\_prag}, \texttt{g\_implied}, \texttt{rating\_g},
\texttt{rating\_adopt}, \texttt{rt} (ms); and once per participant:
\texttt{list\_num}.
The continuous \texttt{pc\_prop\_rating} covariate (topic-level norming
mean) is joined to the trial data during analysis.

\subsection{Procedure}

Participants complete the experiment online via the magpie platform
\citep{kochetkova2024magpie} in a browser.
After reading instructions in German, participants are assigned to one of
4 counterbalancing lists (via URL \texttt{?list=N} parameter assigned by
Prolific; no random fallback in deployment).
They then proceed through 16 trials---8 critical and 8 fillers---in a
fully randomised order, followed by a post-test questionnaire eliciting
native language, age, gender, education level, and optional comments.

Each trial displays:
(a) a context paragraph establishing the communicative situation and the
\textsc{pc\textsubscript{prag}} level;
(b) the target sentence with the best-choice marker already in place in
the Mittelfeld;
(c) the consequent action clause;
(d) the two rating sliders in fixed order ($\widehat{g}$, $\widehat{A}$),
each labelled in German.
The ``Weiter'' (Next) button is disabled until both sliders have been
moved at least once, ensuring a complete response on every trial.

\subsection{Planned Analysis}

\subsubsection{Empirical plots}

For each dependent variable ($\widehat{g}$, $\widehat{A}$):

\begin{itemize}
  \item \textbf{Condition means}: a dot-and-whisker plot shows mean
        ratings per condition (the $2\times2$ of pc\textsubscript{prag}
        $\times$ $g_{\mathrm{implied}}$) with 95\% bootstrapped CIs.
        The expected pattern: $\widehat{g}$ and $\widehat{A}$ should
        be higher in the high-$g_{\mathrm{implied}}$ conditions
        (where the model assigns a stronger marker), and lower in the
        high-pc\textsubscript{prag} conditions (credibility discounting).
  \item \textbf{Marker identity}: because the best-choice marker may
        differ across items within the same condition, a secondary plot
        shows ratings grouped by actual marker identity, allowing
        inspection of whether marker type (rather than just condition)
        drives the effects.
  \item \textbf{pc\textsubscript{prop} covariate}: scatter plots display
        each DV as a function of the continuous norming rating per
        condition, testing whether higher propositional acceptance
        amplifies the $\widehat{g}$ and $\widehat{A}$ signals for
        strong markers.
\end{itemize}

\subsubsection{Bayesian mixed-effects regression}

Two Bayesian mixed-effects regression models are fitted via
\texttt{brms} \citep{burkner2017}, one per DV ($\widehat{g}$ and
$\widehat{A}$):

\begin{align}
y &\sim \mathrm{Beta}(\mu, \phi)
\quad \text{(ratings rescaled to }(0,1)\text{)} \\
\mathrm{logit}(\mu) &= \alpha
+ \beta_{g} \cdot \mathrm{g\_implied\_c}
+ \beta_{\mathrm{prag}} \cdot \mathrm{pc\_prag\_c}
+ \beta_{\mathrm{prop}} \cdot \mathrm{pc\_prop\_c}
+ \beta_{\mathrm{int}} \cdot \mathrm{g\_implied\_c}
  \times \mathrm{pc\_prag\_c}
+ (1 \mid \mathrm{submission\_id})
+ (1 \mid \mathrm{topic}),
\end{align}

where \texttt{g\_implied\_c} is the $g$ level from the corresponding
production condition ($-0.5$ = low, $+0.5$ = high), effectively
coding the marker's position on the production utility scale.
A secondary specification replaces \texttt{g\_implied\_c} with
actualy marker identity as an unordered factor, to examine
non-monotonic patterns and verify that marker rather than
condition label drives effects.

Model comparison via approximate leave-one-out cross-validation (LOO)
selects among additive, interaction, and marker-as-categorical
formulations.

\subsubsection{RSA listener parameter estimation}
\label{subsub:rsa_fit}

The two empirical rating means (per condition cell: 4 conditions) constrain
the listener parameters of the noisy-threshold RSA model.
Specifically, the model predicts:

\begin{align}
\mathbb{E}[\widehat{g} \mid u, \text{context}]
  &= \int g \cdot P(g \mid u, pc_{\mathrm{prag}})\, dg,
\label{eq:pred_g} \\
P(\widehat{A} \mid u, \text{context})
  &= \iiint \sigma(\eta_0 + \eta_g g - \eta_{pc}\, pc
     - \eta_{\mathrm{int}}\, g \cdot pc)\;
     P(pc_{\mathrm{prop}}, pc_{\mathrm{prag}}, g \mid u)\,
     d{pc_{\mathrm{prop}}}\,d{pc_{\mathrm{prag}}}\,dg,
\label{eq:adoption}
\end{align}

\noindent where the posterior $P(pc_{\mathrm{prop}}, pc_{\mathrm{prag}}, g
\mid u)$ is conditioned on the context's \textsc{pc\textsubscript{prag}}
level as a known prior (since it is explicitly given to the listener).

The threshold parameters $(\mu_1,\mu_2, \sigma_1,\sigma_2,
\lambda)$ together with the marker-specific costs $(c_u)$ are estimated
jointly from the production experiment data via the Stan model.
Given those threshold estimates, the listener parameters
$(\eta_0, \eta_g, \eta_{pc}, \eta_{\mathrm{int}})$ are then estimated by
minimising the weighted squared deviation between the model-predicted
means (Eqs.~\ref{eq:pred_g}--\ref{eq:adoption}) and the empirical rating
means from the comprehension experiment, or via a full joint Bayesian
fit over both experiments simultaneously.

\subsection{Model Predictions}

The noisy-threshold RSA model makes the following predictions for
$\widehat{g}$ and $\widehat{A}$, which serve as pre-registered hypotheses.
In all cases, ``strong-marker condition'' refers to conditions 1 and 3
(high $g_{\mathrm{implied}}$, i.e.\ the marker selected by the production
model under a high-$g$ context), and ``weak-marker condition'' to conditions
0 and 2.

\begin{enumerate}
  \item \textbf{Marker encodes inferred goal $\widehat{g}$}: $\widehat{g}$
        ratings should be higher in strong-marker conditions than in
        weak-marker conditions, across both pc\textsubscript{prag} levels.
        This is the core inversion prediction: if the production model
        correctly captures speaker behaviour, listeners should be able to
        decode $g$ from the marker.

  \item \textbf{Marker encodes adoption $\widehat{A}$}: $\widehat{A}$
        ratings should likewise be higher in strong-marker conditions,
        reflecting that higher inferred $g$ translates into greater
        adoption pressure downstream.

  \item \textbf{pc\textsubscript{prag} modulates $\widehat{A}$ but not
        $\widehat{g}$}: because $g$ is orthogonal to pc\textsubscript{prag}
        in the model's utility function, the effect of the marker on
        $\widehat{g}$ should be largely the same in low- and high-
        pc\textsubscript{prag} contexts.
        By contrast, $\widehat{A}$ should be reduced in high-
        pc\textsubscript{prag} contexts via the interaction term
        $\eta_{\mathrm{int}}\,g \cdot pc$: even a strong-marker signal
        for high $g$ is partly cancelled when controversy is high.

  \item \textbf{Credibility discounting}: for strong-marker conditions
        presented in a high-pc\textsubscript{prag} context (condition~3),
        $\widehat{A}$ should be \emph{lower} than for the same marker in
        a low-pc\textsubscript{prag} context (condition~1).
        This is the model's signature prediction distinguishing Bayesian
        inversion from a simpler marker-strength heuristic: a sceptical
        listener who knows controversy is high will discount the adoption
        implication of even the strongest marker.
        The $g_{\mathrm{implied}} \times \textsc{pc\textsubscript{prag}}$
        interaction on $\widehat{A}$ provides the empirical test.
\end{enumerate}

\end{document}
.
%
% Compile stand-alone:
%   pdflatex comprehension_exp.tex
%
% Embed in main.tex:
%   % comprehension_exp.tex
% Stand-alone plan for the comprehension experiment.
% Can be embedded in main.tex via % comprehension_exp.tex
% Stand-alone plan for the comprehension experiment.
% Can be embedded in main.tex via \input{comprehension_exp}.
%
% Compile stand-alone:
%   pdflatex comprehension_exp.tex
%
% Embed in main.tex:
%   \input{comprehension_exp}   % (inside a \section{Experiment} environment)

\documentclass[12pt]{article}
\usepackage{CSPArticleStyle}
\usepackage{csquotes}
\usepackage[margin=1in]{geometry}
\usepackage{booktabs}
\usepackage{parskip}
\usepackage{amsmath}
\usepackage{times}
\usepackage[T1]{fontenc}
\usepackage[utf8]{inputenc}
\usepackage{xcolor}
\usepackage{tabularray}
\usepackage{gb4e}
\noautomath

\title{Comprehension Experiment: Inferring Speaker Goal and Action Adoption
       from Consensus-Marker Use}
\author{}
\date{}

\begin{document}

%% ============================================================
%% When embedding in main.tex, delete from \documentclass to
%% \begin{document} and delete \end{document} at the bottom.
%% ============================================================

\section{Experiment: Consensus-Marker Comprehension}

\subsection{Overview}

The comprehension experiment is the listener-side counterpart of the
production experiment and tests whether listeners can invert the
speaker's marker choice to recover the latent context that motivated it.
Where the production experiment gave participants a context (fixing
\textsc{pc\textsubscript{prag}} and $g$) and asked which marker they
would use, the comprehension experiment gives participants a context
(fixing \textsc{pc\textsubscript{prag}}) together with a marker already
embedded in the sentence and asks what goal strength the speaker likely
had and how likely the listener is to adopt $q$.

The specific marker presented in each vignette is not chosen arbitrarily
from the three candidate markers. Our preregistered plan was to assign,
for each item$\times$condition cell, the \emph{Bayesian best-choice marker}:
the marker with the highest posterior probability under the cost-augmented
noisy-threshold RSA model fitted to the production data, given the context's
\textsc{pc\textsubscript{prag}} level and a specified $g$ level. This was
intended to mirror the production design closely and to ensure that the
comprehension stimuli reflect model-based marker use rather than arbitrary
marker$\times$item pairings.

Applying this strict modal-assignment rule to the fitted production model,
however, yielded a highly imbalanced stimulus set: most critical cells were
assigned the strongest marker, with very little variation across marker
categories. In addition, many of these modal decisions were supported only by
very small probability margins over the next-best alternative. Because such a
stimulus set would make the comprehension experiment substantially less
informative about marker-based listener inferences, we introduced a documented,
exploratory balancing step for stimulus construction. Starting from the same
posterior predictive probabilities, we selected a near-balanced assignment of
the three markers across the 32 critical cells while maximizing total posterior
support under the fitted production model. The strict modal assignment is
retained as the preregistered baseline, but the balanced assignment is used for
the actual comprehension materials.

\textsc{pc\textsubscript{prag}} is retained as an explicit context
manipulation (low vs.\ high), so that listeners can be expected to
have a well-defined prior on pragmatic controversy when interpreting
the marker. \textsc{pc\textsubscript{prop}} is not manipulated within
the experiment; it enters as a continuous topic-level covariate from
the norming study.

After reading each vignette, participants provide two slider ratings
operationalising the key listener-side quantities:
\begin{enumerate}
  \item \textbf{Inferred speaker goal strength} ($\widehat{g}$):
        how strongly does the speaker seem to want the listener to
        accept and act on $q$?
  \item \textbf{Adoption likelihood} ($\widehat{A}$):
        how likely is the listener to follow the recommended action $q$?
\end{enumerate}

\noindent Perceived controversy is \emph{not} elicited as a rating:
because \textsc{pc\textsubscript{prag}} is already manipulated as a
context factor, asking participants to rate controversy would be
redundant with respect to the manipulation check and circular with
respect to the research question. The theoretically central inference
to test is whether listeners can decode the speaker's \emph{goal} from
marker use, and whether that inference translates into adoption
behaviour.

The two slider measures, together with the \textsc{pc\textsubscript{prag}}
context and the \textsc{pc\textsubscript{prop}} covariate, provide the
data needed to fit the model's listener parameters
$(\eta_0, \eta_g, \eta_{pc}, \eta_{\mathrm{int}})$.

\subsection{Participants}

Participants will be recruited via Prolific and will be required to be
native speakers of German (language-at-home filter).
Data collection is planned after completion of the production experiment
(model fitting required to determine the best-choice markers).
A target of $N = 80$ participants is planned ($n = 20$ per list across
4 counterbalancing lists), each completing 16 trials (8 critical + 8
filler), matching the structure of the production experiment.

\subsection{Materials}

\subsubsection{Markers}

Each critical vignette is presented with one specific consensus marker
already embedded in the Mittelfeld of the sentence frame (see
example~(\ref{ex:frame_comp})):

\begin{exe}
  \ex \label{ex:frame_comp}
  Digitale Kompetenzen sind \textbf{[MARKER]}
  in der heutigen Arbeitswelt unverzichtbar.\\
  `Digital skills are \textbf{[MARKER]} indispensable in today's
  working world.'
\end{exe}

Unlike the production experiment, participants do \emph{not} choose the
marker; they read the sentence with the marker already in place and respond
to the two rating probes.

Critically, the marker assigned to each condition is not drawn arbitrarily
from the three candidates. Under the preregistered plan, for each of the four
production conditions (pc\textsubscript{prag} $\times$ $g$), the marker used in
comprehension was to be the \textbf{Bayesian best-choice marker}: the marker
$u^*$ with the highest posterior probability under the fitted
cost-augmented noisy-threshold RSA production model,
\[
u^* = \arg\max_{u} \; P(u \mid pc_{\mathrm{prag}}, g,
                           pc_{\mathrm{prop}\_c}).
\]
Because $pc_{\mathrm{prop}}$ varies continuously across items,
$u^*$ may differ per item even within the same
(pc\textsubscript{prag}, $g$) cell; the final stimulus list is therefore
constructed item-by-item once the production model has been fitted.
In practice, the fitted production model yielded a nearly degenerate modal map,
with most cells assigned to \textit{bekanntlich}. We therefore retained the
model-based logic of the procedure but replaced strict cellwise modal selection
with an uncertainty-aware balancing step: among the marker assignments with
high posterior support, we selected a near-balanced configuration across the 32
critical cells. This preserves the production--comprehension link while also
ensuring that the comprehension experiment contains sufficient marker variation
to test listener-side inferences.

All three markers remain candidates; the mapping from conditions to
markers is an empirical output of the production analysis, not a design
pre-commitment. In the final comprehension materials, we use the
\textbf{balanced model-based assignment}, rather than the strict modal
assignment, in order to maintain informative variation across marker types.

\subsubsection{Topic Items}

The same eight topic domains used in the norming study and the production
experiment are used here (Table~\ref{tab:items_comp}).
The 4-condition $\times$ 8-item Latin square (Section~\ref{subsub:counterbalancing_comp})
requires no additional items.

\begin{table}[ht]
\centering
\caption{Topic items: sentence frame for $p$ (with marker position [M])
         and consequent $q$.}
\label{tab:items_comp}
\small
\begin{tabular}{clll}
\toprule
\textbf{\#} & \textbf{Topic} & \textbf{Sentence frame} & \textbf{Consequent $q$} \\
\midrule
1 & Klimawandel &
  \textit{Der Klimawandel wird [M]} &
  auf Flugreisen verzichten \\
  & & \textit{durch menschliche Aktivit\"aten verursacht.} & \\[3pt]
2 & Ern\"ahrung &
  \textit{Eine \"uberwiegend pflanzliche Ern\"ahrung} &
  mehr pflanzliche Lebens- \\
  & & \textit{f\"ordert [M] die Gesundheit.} &
  mittel einbauen \\[3pt]
3 & Stadtverkehr &
  \textit{Ein ausgebautes Nahverkehrsnetz} &
  h\"aufiger auf \\
  & & \textit{entlastet [M] den Stadtverkehr} &
  \"offentliche Verkehrs- \\
  & & \textit{deutlich.} &
  mittel umsteigen \\[3pt]
4 & Digitalisierung &
  \textit{Digitale Kompetenzen sind [M] in} &
  regelm\"a\ss{}ig Weiter- \\
  & & \textit{der heutigen Arbeitswelt unverzichtbar.} &
  bildungsangebote nutzen \\[3pt]
5 & Schlaf &
  \textit{Ausreichend Schlaf ist [M] entscheidend} &
  auf feste Schlafzeiten \\
  & & \textit{f\"ur die kognitive Leistungsf\"ahigkeit.} &
  achten \\[3pt]
6 & Plastik &
  \textit{Einwegplastik schadet [M] den} &
  konsequent auf Einweg- \\
  & & \textit{Meeres\"okosystemen erheblich.} &
  plastik verzichten \\[3pt]
7 & Sport &
  \textit{Regelm\"a\ss{}ige k\"orperliche Bewegung} &
  regelm\"a\ss{}ig Sport in \\
  & & \textit{senkt [M] das Risiko f\"ur} &
  den Alltag integrieren \\
  & & \textit{Herz-Kreislauf-Erkrankungen deutlich.} & \\[3pt]
8 & Lokalkauf &
  \textit{Das Kaufen bei lokalen H\"andlern} &
  \"ofter bei lokalen \\
  & & \textit{st\"arkt [M] die regionale Wirtschaft} &
  Gesch\"aften einkaufen \\
  & & \textit{nachhaltig.} & \\
\bottomrule
\end{tabular}
\end{table}

\noindent [M] denotes the position of the pre-filled best-choice marker.

\subsubsection{Vignettes}

\paragraph{Structure.}
Each trial presents a vignette consisting of three parts:
(a) a context paragraph establishing the communicative situation and
fixing the \textsc{pc\textsubscript{prag}} level;
(b) the target sentence with the best-choice marker already embedded
in the Mittelfeld;
(c) the consequent action clause.
The two rating probes follow immediately below the vignette.

\paragraph{Propositional perceived controversy (\textsc{pc\textsubscript{prop}}).}
As in the production experiment, \textsc{pc\textsubscript{prop}} is not
manipulated within the comprehension experiment.
Topic-level propositional acceptance ratings from the norming study
(0--100 scale; e.g., \textit{Digitalisierung}: 85/100,
\textit{Lokalkauf}: 58/100) serve as a continuous, mean-centred
covariate \textsc{pc\textsubscript{prop}}$_c$ in analysis.
Context paragraphs present $p$ without editorial evaluation; variation in
listeners' individual assessments of propositional controversy is
captured by the norming-derived covariate.

\paragraph{Pragmatic perceived controversy (\textsc{pc\textsubscript{prag}}).}
\textsc{pc\textsubscript{prag}} is manipulated between item instances
using the same two context versions as in the production experiment:
the speaker's social circle either \emph{largely agrees} (low) or
\emph{holds divergent views} (high) about $p$.
For example, for the \textit{Digitalisierung} item:

\begin{itemize}
  \item $\textsc{pc}_\text{prag} = \text{low}$: \textit{Deine Kollegin spricht
        mit dir \"uber Berufsausbildung und Weiterbildung.
        In eurem Team teilen die meisten die Ansicht, dass digitale Kompetenzen
        in der heutigen Arbeitswelt unverzichtbar sind.}
  \item $\textsc{pc}_\text{prag} = \text{high}$: \textit{Deine Kollegin spricht
        mit dir \"uber Berufsausbildung und Weiterbildung.
        In eurem Team gehen die Meinungen dazu stark auseinander.}
\end{itemize}

The context is framed from the listener's perspective as a direct addressee,
mirroring the production study where participants adopt the speaker's
perspective.

\paragraph{Marker manipulation.}
The dependent measures of the production experiment (which marker is chosen)
become the independent variable here.
One of the three consensus markers is inserted into the Mittelfeld of the
target sentence.
A complete sample item for the \textit{Digitalisierung} topic
($\textsc{pc}_\text{prag} = \text{high}$, marker = \textit{bekanntlich};
norming $\textsc{pc}_\text{prop}$ rating: 85/100) is given in
(\ref{ex:comp_sample}):

\begin{exe}
  \ex \label{ex:comp_sample}
  \textit{Deine Kollegin spricht mit dir \"uber Berufsausbildung und
  Weiterbildung.
  In eurem Team gehen die Meinungen dazu stark auseinander.}

  \textit{Sie sagt zu dir:} Digitale Kompetenzen sind \textbf{bekanntlich}
  in der heutigen Arbeitswelt unverzichtbar.
  Deshalb solltest du regelm\"a\ss{}ig digitale Weiterbildungsangebote nutzen.

  \noindent\textit{[Rating probes follow.]}
\end{exe}

\subsection{Design}

\subsubsection{Factorial structure}

The comprehension experiment mirrors the production experiment's
$\mathbf{2}$~\textbf{(pc\textsubscript{prag}: low vs.\ high)}
$\mathbf{\times}$~$\mathbf{2}$~\textbf{($g_{\mathrm{implied}}$: low
vs.\ high) design}, yielding \textbf{4 conditions}.
The $g_{\mathrm{implied}}$ dimension is not manipulated via an explicit
goal instruction (as in production) but is instead conveyed entirely
through the choice of marker: for each condition, the marker presented
is the best-choice marker for the corresponding
(\textsc{pc\textsubscript{prag}},~$g$) cell from the production model.
\textsc{pc\textsubscript{prop}} enters as a continuous norming covariate.
With 8 topic items and a Latin square rotation across 4 lists, each
participant sees exactly \textbf{8 critical trials}---one item per
condition, each condition appearing twice---plus 8 filler trials,
for \textbf{16 trials} in total (Table~\ref{tab:design_comp}).

\begin{table}[ht]
\centering
\caption{%
  The $2\times2$ design (4 conditions), mirroring the production experiment.
  The marker presented in each condition is the Bayesian best-choice
  marker from the fitted production model; specific markers are
  determined after production data collection.}
\label{tab:design_comp}
\small
\begin{tabular}{ccll}
\toprule
\textbf{pc\textsubscript{prag}} &
\textbf{$g_{\mathrm{implied}}$} &
\textbf{Marker (TBD from production)} &
\textbf{Condition label} \\
\midrule
low  & low  & best-choice for $(\text{low}, g_{\text{low}})$  &
  (0) community agrees, weak signal \\
low  & high & best-choice for $(\text{low}, g_{\text{high}})$ &
  (1) community agrees, strong signal \\
high & low  & best-choice for $(\text{high}, g_{\text{low}})$ &
  (2) community divided, weak signal \\
high & high & best-choice for $(\text{high}, g_{\text{high}})$ &
  (3) community divided, strong signal \\
\bottomrule
\end{tabular}
\end{table}

\subsubsection{Counterbalancing}
\label{subsub:counterbalancing_comp}

The experiment uses the \textbf{same $4 \times 8$ Latin square} as the
production experiment: each of the 4 lists assigns item $i$
($i = 0,\ldots,7$) to condition $(i + l)\bmod 4$, where
$l \in \{0,1,2,3\}$ is the list number.
This guarantees that:
\begin{itemize}
  \item each participant sees every item exactly once;
  \item each participant sees each of the 4 conditions exactly twice
        (two items per condition), ensuring balanced coverage of the
        $2\times2$ design;
  \item across all 4 lists, each item appears equally often in each
        condition.
\end{itemize}

List assignment follows the same Prolific-and-URL strategy as the production
experiment: each Prolific study embeds \texttt{?list=N} (0--3) in the
experiment URL, capping recruitment at $n = 20$ participants per list
for a total of $N = 80$.

In addition to the 8 critical trials, each participant completed 8
\textbf{filler trials} using the same filler items as the production
experiment (train travel, reading light, cooking, coffee, vacation,
hydration, fresh air, noise), each presented with a randomly assigned
marker.
Fillers are excluded from analysis; they serve to dilute the critical
items and reduce demand characteristics.

\subsubsection{Dependent variables}
\label{subsub:dvs_comp}

Two slider ratings are collected after each vignette (0--100 continuous
scale with no visible numbers, labelled at the endpoints as shown):

\begin{enumerate}
  \item \textbf{Inferred speaker goal strength} (operationalises $\widehat{g}$):
  \begin{quote}
    \textit{Wie stark hast du den Eindruck, dass die Sprecherin m\"ochte,
    dass du ihre Empfehlung befolgst?}\\
    (`How strongly do you get the impression that the speaker wants you to
    follow their recommendation?')\\
    Endpoints: 0 = \textit{\"uberhaupt nicht}
    (`not at all') ---
    100 = \textit{sehr stark} (`very strongly').
  \end{quote}

  \item \textbf{Adoption likelihood of $q$} (operationalises $\widehat{A}$):
  \begin{quote}
    \textit{Wie wahrscheinlich ist es, dass du die empfohlene Handlung
    dann auch tats\"achlich umsetzt?}\\
    (`How likely is it that you will actually carry out the recommended
    action?')\\
    Endpoints: 0 = \textit{sehr unwahrscheinlich}
    (`very unlikely') ---
    100 = \textit{sehr wahrscheinlich} (`very likely').
  \end{quote}
\end{enumerate}

\noindent Perceived controversy is deliberately \emph{not} included as a
rating probe: \textsc{pc\textsubscript{prag}} is experimentally
manipulated via the context paragraph, so a controversy rating would
be redundant and risks obscuring the marker's contribution to inferred
goal strength.

Recorded per trial: \texttt{trial\_id}, \texttt{topic},
\texttt{is\_filler}, \texttt{condition\_index}, \texttt{marker},
\texttt{pc\_prag}, \texttt{g\_implied}, \texttt{rating\_g},
\texttt{rating\_adopt}, \texttt{rt} (ms); and once per participant:
\texttt{list\_num}.
The continuous \texttt{pc\_prop\_rating} covariate (topic-level norming
mean) is joined to the trial data during analysis.

\subsection{Procedure}

Participants complete the experiment online via the magpie platform
\citep{kochetkova2024magpie} in a browser.
After reading instructions in German, participants are assigned to one of
4 counterbalancing lists (via URL \texttt{?list=N} parameter assigned by
Prolific; no random fallback in deployment).
They then proceed through 16 trials---8 critical and 8 fillers---in a
fully randomised order, followed by a post-test questionnaire eliciting
native language, age, gender, education level, and optional comments.

Each trial displays:
(a) a context paragraph establishing the communicative situation and the
\textsc{pc\textsubscript{prag}} level;
(b) the target sentence with the best-choice marker already in place in
the Mittelfeld;
(c) the consequent action clause;
(d) the two rating sliders in fixed order ($\widehat{g}$, $\widehat{A}$),
each labelled in German.
The ``Weiter'' (Next) button is disabled until both sliders have been
moved at least once, ensuring a complete response on every trial.

\subsection{Planned Analysis}

\subsubsection{Empirical plots}

For each dependent variable ($\widehat{g}$, $\widehat{A}$):

\begin{itemize}
  \item \textbf{Condition means}: a dot-and-whisker plot shows mean
        ratings per condition (the $2\times2$ of pc\textsubscript{prag}
        $\times$ $g_{\mathrm{implied}}$) with 95\% bootstrapped CIs.
        The expected pattern: $\widehat{g}$ and $\widehat{A}$ should
        be higher in the high-$g_{\mathrm{implied}}$ conditions
        (where the model assigns a stronger marker), and lower in the
        high-pc\textsubscript{prag} conditions (credibility discounting).
  \item \textbf{Marker identity}: because the best-choice marker may
        differ across items within the same condition, a secondary plot
        shows ratings grouped by actual marker identity, allowing
        inspection of whether marker type (rather than just condition)
        drives the effects.
  \item \textbf{pc\textsubscript{prop} covariate}: scatter plots display
        each DV as a function of the continuous norming rating per
        condition, testing whether higher propositional acceptance
        amplifies the $\widehat{g}$ and $\widehat{A}$ signals for
        strong markers.
\end{itemize}

\subsubsection{Bayesian mixed-effects regression}

Two Bayesian mixed-effects regression models are fitted via
\texttt{brms} \citep{burkner2017}, one per DV ($\widehat{g}$ and
$\widehat{A}$):

\begin{align}
y &\sim \mathrm{Beta}(\mu, \phi)
\quad \text{(ratings rescaled to }(0,1)\text{)} \\
\mathrm{logit}(\mu) &= \alpha
+ \beta_{g} \cdot \mathrm{g\_implied\_c}
+ \beta_{\mathrm{prag}} \cdot \mathrm{pc\_prag\_c}
+ \beta_{\mathrm{prop}} \cdot \mathrm{pc\_prop\_c}
+ \beta_{\mathrm{int}} \cdot \mathrm{g\_implied\_c}
  \times \mathrm{pc\_prag\_c}
+ (1 \mid \mathrm{submission\_id})
+ (1 \mid \mathrm{topic}),
\end{align}

where \texttt{g\_implied\_c} is the $g$ level from the corresponding
production condition ($-0.5$ = low, $+0.5$ = high), effectively
coding the marker's position on the production utility scale.
A secondary specification replaces \texttt{g\_implied\_c} with
actualy marker identity as an unordered factor, to examine
non-monotonic patterns and verify that marker rather than
condition label drives effects.

Model comparison via approximate leave-one-out cross-validation (LOO)
selects among additive, interaction, and marker-as-categorical
formulations.

\subsubsection{RSA listener parameter estimation}
\label{subsub:rsa_fit}

The two empirical rating means (per condition cell: 4 conditions) constrain
the listener parameters of the noisy-threshold RSA model.
Specifically, the model predicts:

\begin{align}
\mathbb{E}[\widehat{g} \mid u, \text{context}]
  &= \int g \cdot P(g \mid u, pc_{\mathrm{prag}})\, dg,
\label{eq:pred_g} \\
P(\widehat{A} \mid u, \text{context})
  &= \iiint \sigma(\eta_0 + \eta_g g - \eta_{pc}\, pc
     - \eta_{\mathrm{int}}\, g \cdot pc)\;
     P(pc_{\mathrm{prop}}, pc_{\mathrm{prag}}, g \mid u)\,
     d{pc_{\mathrm{prop}}}\,d{pc_{\mathrm{prag}}}\,dg,
\label{eq:adoption}
\end{align}

\noindent where the posterior $P(pc_{\mathrm{prop}}, pc_{\mathrm{prag}}, g
\mid u)$ is conditioned on the context's \textsc{pc\textsubscript{prag}}
level as a known prior (since it is explicitly given to the listener).

The threshold parameters $(\mu_1,\mu_2, \sigma_1,\sigma_2,
\lambda)$ together with the marker-specific costs $(c_u)$ are estimated
jointly from the production experiment data via the Stan model.
Given those threshold estimates, the listener parameters
$(\eta_0, \eta_g, \eta_{pc}, \eta_{\mathrm{int}})$ are then estimated by
minimising the weighted squared deviation between the model-predicted
means (Eqs.~\ref{eq:pred_g}--\ref{eq:adoption}) and the empirical rating
means from the comprehension experiment, or via a full joint Bayesian
fit over both experiments simultaneously.

\subsection{Model Predictions}

The noisy-threshold RSA model makes the following predictions for
$\widehat{g}$ and $\widehat{A}$, which serve as pre-registered hypotheses.
In all cases, ``strong-marker condition'' refers to conditions 1 and 3
(high $g_{\mathrm{implied}}$, i.e.\ the marker selected by the production
model under a high-$g$ context), and ``weak-marker condition'' to conditions
0 and 2.

\begin{enumerate}
  \item \textbf{Marker encodes inferred goal $\widehat{g}$}: $\widehat{g}$
        ratings should be higher in strong-marker conditions than in
        weak-marker conditions, across both pc\textsubscript{prag} levels.
        This is the core inversion prediction: if the production model
        correctly captures speaker behaviour, listeners should be able to
        decode $g$ from the marker.

  \item \textbf{Marker encodes adoption $\widehat{A}$}: $\widehat{A}$
        ratings should likewise be higher in strong-marker conditions,
        reflecting that higher inferred $g$ translates into greater
        adoption pressure downstream.

  \item \textbf{pc\textsubscript{prag} modulates $\widehat{A}$ but not
        $\widehat{g}$}: because $g$ is orthogonal to pc\textsubscript{prag}
        in the model's utility function, the effect of the marker on
        $\widehat{g}$ should be largely the same in low- and high-
        pc\textsubscript{prag} contexts.
        By contrast, $\widehat{A}$ should be reduced in high-
        pc\textsubscript{prag} contexts via the interaction term
        $\eta_{\mathrm{int}}\,g \cdot pc$: even a strong-marker signal
        for high $g$ is partly cancelled when controversy is high.

  \item \textbf{Credibility discounting}: for strong-marker conditions
        presented in a high-pc\textsubscript{prag} context (condition~3),
        $\widehat{A}$ should be \emph{lower} than for the same marker in
        a low-pc\textsubscript{prag} context (condition~1).
        This is the model's signature prediction distinguishing Bayesian
        inversion from a simpler marker-strength heuristic: a sceptical
        listener who knows controversy is high will discount the adoption
        implication of even the strongest marker.
        The $g_{\mathrm{implied}} \times \textsc{pc\textsubscript{prag}}$
        interaction on $\widehat{A}$ provides the empirical test.
\end{enumerate}

\end{document}
.
%
% Compile stand-alone:
%   pdflatex comprehension_exp.tex
%
% Embed in main.tex:
%   % comprehension_exp.tex
% Stand-alone plan for the comprehension experiment.
% Can be embedded in main.tex via \input{comprehension_exp}.
%
% Compile stand-alone:
%   pdflatex comprehension_exp.tex
%
% Embed in main.tex:
%   \input{comprehension_exp}   % (inside a \section{Experiment} environment)

\documentclass[12pt]{article}
\usepackage{CSPArticleStyle}
\usepackage{csquotes}
\usepackage[margin=1in]{geometry}
\usepackage{booktabs}
\usepackage{parskip}
\usepackage{amsmath}
\usepackage{times}
\usepackage[T1]{fontenc}
\usepackage[utf8]{inputenc}
\usepackage{xcolor}
\usepackage{tabularray}
\usepackage{gb4e}
\noautomath

\title{Comprehension Experiment: Inferring Speaker Goal and Action Adoption
       from Consensus-Marker Use}
\author{}
\date{}

\begin{document}

%% ============================================================
%% When embedding in main.tex, delete from \documentclass to
%% \begin{document} and delete \end{document} at the bottom.
%% ============================================================

\section{Experiment: Consensus-Marker Comprehension}

\subsection{Overview}

The comprehension experiment is the listener-side counterpart of the
production experiment and tests whether listeners can invert the
speaker's marker choice to recover the latent context that motivated it.
Where the production experiment gave participants a context (fixing
\textsc{pc\textsubscript{prag}} and $g$) and asked which marker they
would use, the comprehension experiment gives participants a context
(fixing \textsc{pc\textsubscript{prag}}) together with a marker already
embedded in the sentence and asks what goal strength the speaker likely
had and how likely the listener is to adopt $q$.

The specific marker presented in each vignette is not chosen arbitrarily
from the three candidate markers. Our preregistered plan was to assign,
for each item$\times$condition cell, the \emph{Bayesian best-choice marker}:
the marker with the highest posterior probability under the cost-augmented
noisy-threshold RSA model fitted to the production data, given the context's
\textsc{pc\textsubscript{prag}} level and a specified $g$ level. This was
intended to mirror the production design closely and to ensure that the
comprehension stimuli reflect model-based marker use rather than arbitrary
marker$\times$item pairings.

Applying this strict modal-assignment rule to the fitted production model,
however, yielded a highly imbalanced stimulus set: most critical cells were
assigned the strongest marker, with very little variation across marker
categories. In addition, many of these modal decisions were supported only by
very small probability margins over the next-best alternative. Because such a
stimulus set would make the comprehension experiment substantially less
informative about marker-based listener inferences, we introduced a documented,
exploratory balancing step for stimulus construction. Starting from the same
posterior predictive probabilities, we selected a near-balanced assignment of
the three markers across the 32 critical cells while maximizing total posterior
support under the fitted production model. The strict modal assignment is
retained as the preregistered baseline, but the balanced assignment is used for
the actual comprehension materials.

\textsc{pc\textsubscript{prag}} is retained as an explicit context
manipulation (low vs.\ high), so that listeners can be expected to
have a well-defined prior on pragmatic controversy when interpreting
the marker. \textsc{pc\textsubscript{prop}} is not manipulated within
the experiment; it enters as a continuous topic-level covariate from
the norming study.

After reading each vignette, participants provide two slider ratings
operationalising the key listener-side quantities:
\begin{enumerate}
  \item \textbf{Inferred speaker goal strength} ($\widehat{g}$):
        how strongly does the speaker seem to want the listener to
        accept and act on $q$?
  \item \textbf{Adoption likelihood} ($\widehat{A}$):
        how likely is the listener to follow the recommended action $q$?
\end{enumerate}

\noindent Perceived controversy is \emph{not} elicited as a rating:
because \textsc{pc\textsubscript{prag}} is already manipulated as a
context factor, asking participants to rate controversy would be
redundant with respect to the manipulation check and circular with
respect to the research question. The theoretically central inference
to test is whether listeners can decode the speaker's \emph{goal} from
marker use, and whether that inference translates into adoption
behaviour.

The two slider measures, together with the \textsc{pc\textsubscript{prag}}
context and the \textsc{pc\textsubscript{prop}} covariate, provide the
data needed to fit the model's listener parameters
$(\eta_0, \eta_g, \eta_{pc}, \eta_{\mathrm{int}})$.

\subsection{Participants}

Participants will be recruited via Prolific and will be required to be
native speakers of German (language-at-home filter).
Data collection is planned after completion of the production experiment
(model fitting required to determine the best-choice markers).
A target of $N = 80$ participants is planned ($n = 20$ per list across
4 counterbalancing lists), each completing 16 trials (8 critical + 8
filler), matching the structure of the production experiment.

\subsection{Materials}

\subsubsection{Markers}

Each critical vignette is presented with one specific consensus marker
already embedded in the Mittelfeld of the sentence frame (see
example~(\ref{ex:frame_comp})):

\begin{exe}
  \ex \label{ex:frame_comp}
  Digitale Kompetenzen sind \textbf{[MARKER]}
  in der heutigen Arbeitswelt unverzichtbar.\\
  `Digital skills are \textbf{[MARKER]} indispensable in today's
  working world.'
\end{exe}

Unlike the production experiment, participants do \emph{not} choose the
marker; they read the sentence with the marker already in place and respond
to the two rating probes.

Critically, the marker assigned to each condition is not drawn arbitrarily
from the three candidates. Under the preregistered plan, for each of the four
production conditions (pc\textsubscript{prag} $\times$ $g$), the marker used in
comprehension was to be the \textbf{Bayesian best-choice marker}: the marker
$u^*$ with the highest posterior probability under the fitted
cost-augmented noisy-threshold RSA production model,
\[
u^* = \arg\max_{u} \; P(u \mid pc_{\mathrm{prag}}, g,
                           pc_{\mathrm{prop}\_c}).
\]
Because $pc_{\mathrm{prop}}$ varies continuously across items,
$u^*$ may differ per item even within the same
(pc\textsubscript{prag}, $g$) cell; the final stimulus list is therefore
constructed item-by-item once the production model has been fitted.
In practice, the fitted production model yielded a nearly degenerate modal map,
with most cells assigned to \textit{bekanntlich}. We therefore retained the
model-based logic of the procedure but replaced strict cellwise modal selection
with an uncertainty-aware balancing step: among the marker assignments with
high posterior support, we selected a near-balanced configuration across the 32
critical cells. This preserves the production--comprehension link while also
ensuring that the comprehension experiment contains sufficient marker variation
to test listener-side inferences.

All three markers remain candidates; the mapping from conditions to
markers is an empirical output of the production analysis, not a design
pre-commitment. In the final comprehension materials, we use the
\textbf{balanced model-based assignment}, rather than the strict modal
assignment, in order to maintain informative variation across marker types.

\subsubsection{Topic Items}

The same eight topic domains used in the norming study and the production
experiment are used here (Table~\ref{tab:items_comp}).
The 4-condition $\times$ 8-item Latin square (Section~\ref{subsub:counterbalancing_comp})
requires no additional items.

\begin{table}[ht]
\centering
\caption{Topic items: sentence frame for $p$ (with marker position [M])
         and consequent $q$.}
\label{tab:items_comp}
\small
\begin{tabular}{clll}
\toprule
\textbf{\#} & \textbf{Topic} & \textbf{Sentence frame} & \textbf{Consequent $q$} \\
\midrule
1 & Klimawandel &
  \textit{Der Klimawandel wird [M]} &
  auf Flugreisen verzichten \\
  & & \textit{durch menschliche Aktivit\"aten verursacht.} & \\[3pt]
2 & Ern\"ahrung &
  \textit{Eine \"uberwiegend pflanzliche Ern\"ahrung} &
  mehr pflanzliche Lebens- \\
  & & \textit{f\"ordert [M] die Gesundheit.} &
  mittel einbauen \\[3pt]
3 & Stadtverkehr &
  \textit{Ein ausgebautes Nahverkehrsnetz} &
  h\"aufiger auf \\
  & & \textit{entlastet [M] den Stadtverkehr} &
  \"offentliche Verkehrs- \\
  & & \textit{deutlich.} &
  mittel umsteigen \\[3pt]
4 & Digitalisierung &
  \textit{Digitale Kompetenzen sind [M] in} &
  regelm\"a\ss{}ig Weiter- \\
  & & \textit{der heutigen Arbeitswelt unverzichtbar.} &
  bildungsangebote nutzen \\[3pt]
5 & Schlaf &
  \textit{Ausreichend Schlaf ist [M] entscheidend} &
  auf feste Schlafzeiten \\
  & & \textit{f\"ur die kognitive Leistungsf\"ahigkeit.} &
  achten \\[3pt]
6 & Plastik &
  \textit{Einwegplastik schadet [M] den} &
  konsequent auf Einweg- \\
  & & \textit{Meeres\"okosystemen erheblich.} &
  plastik verzichten \\[3pt]
7 & Sport &
  \textit{Regelm\"a\ss{}ige k\"orperliche Bewegung} &
  regelm\"a\ss{}ig Sport in \\
  & & \textit{senkt [M] das Risiko f\"ur} &
  den Alltag integrieren \\
  & & \textit{Herz-Kreislauf-Erkrankungen deutlich.} & \\[3pt]
8 & Lokalkauf &
  \textit{Das Kaufen bei lokalen H\"andlern} &
  \"ofter bei lokalen \\
  & & \textit{st\"arkt [M] die regionale Wirtschaft} &
  Gesch\"aften einkaufen \\
  & & \textit{nachhaltig.} & \\
\bottomrule
\end{tabular}
\end{table}

\noindent [M] denotes the position of the pre-filled best-choice marker.

\subsubsection{Vignettes}

\paragraph{Structure.}
Each trial presents a vignette consisting of three parts:
(a) a context paragraph establishing the communicative situation and
fixing the \textsc{pc\textsubscript{prag}} level;
(b) the target sentence with the best-choice marker already embedded
in the Mittelfeld;
(c) the consequent action clause.
The two rating probes follow immediately below the vignette.

\paragraph{Propositional perceived controversy (\textsc{pc\textsubscript{prop}}).}
As in the production experiment, \textsc{pc\textsubscript{prop}} is not
manipulated within the comprehension experiment.
Topic-level propositional acceptance ratings from the norming study
(0--100 scale; e.g., \textit{Digitalisierung}: 85/100,
\textit{Lokalkauf}: 58/100) serve as a continuous, mean-centred
covariate \textsc{pc\textsubscript{prop}}$_c$ in analysis.
Context paragraphs present $p$ without editorial evaluation; variation in
listeners' individual assessments of propositional controversy is
captured by the norming-derived covariate.

\paragraph{Pragmatic perceived controversy (\textsc{pc\textsubscript{prag}}).}
\textsc{pc\textsubscript{prag}} is manipulated between item instances
using the same two context versions as in the production experiment:
the speaker's social circle either \emph{largely agrees} (low) or
\emph{holds divergent views} (high) about $p$.
For example, for the \textit{Digitalisierung} item:

\begin{itemize}
  \item $\textsc{pc}_\text{prag} = \text{low}$: \textit{Deine Kollegin spricht
        mit dir \"uber Berufsausbildung und Weiterbildung.
        In eurem Team teilen die meisten die Ansicht, dass digitale Kompetenzen
        in der heutigen Arbeitswelt unverzichtbar sind.}
  \item $\textsc{pc}_\text{prag} = \text{high}$: \textit{Deine Kollegin spricht
        mit dir \"uber Berufsausbildung und Weiterbildung.
        In eurem Team gehen die Meinungen dazu stark auseinander.}
\end{itemize}

The context is framed from the listener's perspective as a direct addressee,
mirroring the production study where participants adopt the speaker's
perspective.

\paragraph{Marker manipulation.}
The dependent measures of the production experiment (which marker is chosen)
become the independent variable here.
One of the three consensus markers is inserted into the Mittelfeld of the
target sentence.
A complete sample item for the \textit{Digitalisierung} topic
($\textsc{pc}_\text{prag} = \text{high}$, marker = \textit{bekanntlich};
norming $\textsc{pc}_\text{prop}$ rating: 85/100) is given in
(\ref{ex:comp_sample}):

\begin{exe}
  \ex \label{ex:comp_sample}
  \textit{Deine Kollegin spricht mit dir \"uber Berufsausbildung und
  Weiterbildung.
  In eurem Team gehen die Meinungen dazu stark auseinander.}

  \textit{Sie sagt zu dir:} Digitale Kompetenzen sind \textbf{bekanntlich}
  in der heutigen Arbeitswelt unverzichtbar.
  Deshalb solltest du regelm\"a\ss{}ig digitale Weiterbildungsangebote nutzen.

  \noindent\textit{[Rating probes follow.]}
\end{exe}

\subsection{Design}

\subsubsection{Factorial structure}

The comprehension experiment mirrors the production experiment's
$\mathbf{2}$~\textbf{(pc\textsubscript{prag}: low vs.\ high)}
$\mathbf{\times}$~$\mathbf{2}$~\textbf{($g_{\mathrm{implied}}$: low
vs.\ high) design}, yielding \textbf{4 conditions}.
The $g_{\mathrm{implied}}$ dimension is not manipulated via an explicit
goal instruction (as in production) but is instead conveyed entirely
through the choice of marker: for each condition, the marker presented
is the best-choice marker for the corresponding
(\textsc{pc\textsubscript{prag}},~$g$) cell from the production model.
\textsc{pc\textsubscript{prop}} enters as a continuous norming covariate.
With 8 topic items and a Latin square rotation across 4 lists, each
participant sees exactly \textbf{8 critical trials}---one item per
condition, each condition appearing twice---plus 8 filler trials,
for \textbf{16 trials} in total (Table~\ref{tab:design_comp}).

\begin{table}[ht]
\centering
\caption{%
  The $2\times2$ design (4 conditions), mirroring the production experiment.
  The marker presented in each condition is the Bayesian best-choice
  marker from the fitted production model; specific markers are
  determined after production data collection.}
\label{tab:design_comp}
\small
\begin{tabular}{ccll}
\toprule
\textbf{pc\textsubscript{prag}} &
\textbf{$g_{\mathrm{implied}}$} &
\textbf{Marker (TBD from production)} &
\textbf{Condition label} \\
\midrule
low  & low  & best-choice for $(\text{low}, g_{\text{low}})$  &
  (0) community agrees, weak signal \\
low  & high & best-choice for $(\text{low}, g_{\text{high}})$ &
  (1) community agrees, strong signal \\
high & low  & best-choice for $(\text{high}, g_{\text{low}})$ &
  (2) community divided, weak signal \\
high & high & best-choice for $(\text{high}, g_{\text{high}})$ &
  (3) community divided, strong signal \\
\bottomrule
\end{tabular}
\end{table}

\subsubsection{Counterbalancing}
\label{subsub:counterbalancing_comp}

The experiment uses the \textbf{same $4 \times 8$ Latin square} as the
production experiment: each of the 4 lists assigns item $i$
($i = 0,\ldots,7$) to condition $(i + l)\bmod 4$, where
$l \in \{0,1,2,3\}$ is the list number.
This guarantees that:
\begin{itemize}
  \item each participant sees every item exactly once;
  \item each participant sees each of the 4 conditions exactly twice
        (two items per condition), ensuring balanced coverage of the
        $2\times2$ design;
  \item across all 4 lists, each item appears equally often in each
        condition.
\end{itemize}

List assignment follows the same Prolific-and-URL strategy as the production
experiment: each Prolific study embeds \texttt{?list=N} (0--3) in the
experiment URL, capping recruitment at $n = 20$ participants per list
for a total of $N = 80$.

In addition to the 8 critical trials, each participant completed 8
\textbf{filler trials} using the same filler items as the production
experiment (train travel, reading light, cooking, coffee, vacation,
hydration, fresh air, noise), each presented with a randomly assigned
marker.
Fillers are excluded from analysis; they serve to dilute the critical
items and reduce demand characteristics.

\subsubsection{Dependent variables}
\label{subsub:dvs_comp}

Two slider ratings are collected after each vignette (0--100 continuous
scale with no visible numbers, labelled at the endpoints as shown):

\begin{enumerate}
  \item \textbf{Inferred speaker goal strength} (operationalises $\widehat{g}$):
  \begin{quote}
    \textit{Wie stark hast du den Eindruck, dass die Sprecherin m\"ochte,
    dass du ihre Empfehlung befolgst?}\\
    (`How strongly do you get the impression that the speaker wants you to
    follow their recommendation?')\\
    Endpoints: 0 = \textit{\"uberhaupt nicht}
    (`not at all') ---
    100 = \textit{sehr stark} (`very strongly').
  \end{quote}

  \item \textbf{Adoption likelihood of $q$} (operationalises $\widehat{A}$):
  \begin{quote}
    \textit{Wie wahrscheinlich ist es, dass du die empfohlene Handlung
    dann auch tats\"achlich umsetzt?}\\
    (`How likely is it that you will actually carry out the recommended
    action?')\\
    Endpoints: 0 = \textit{sehr unwahrscheinlich}
    (`very unlikely') ---
    100 = \textit{sehr wahrscheinlich} (`very likely').
  \end{quote}
\end{enumerate}

\noindent Perceived controversy is deliberately \emph{not} included as a
rating probe: \textsc{pc\textsubscript{prag}} is experimentally
manipulated via the context paragraph, so a controversy rating would
be redundant and risks obscuring the marker's contribution to inferred
goal strength.

Recorded per trial: \texttt{trial\_id}, \texttt{topic},
\texttt{is\_filler}, \texttt{condition\_index}, \texttt{marker},
\texttt{pc\_prag}, \texttt{g\_implied}, \texttt{rating\_g},
\texttt{rating\_adopt}, \texttt{rt} (ms); and once per participant:
\texttt{list\_num}.
The continuous \texttt{pc\_prop\_rating} covariate (topic-level norming
mean) is joined to the trial data during analysis.

\subsection{Procedure}

Participants complete the experiment online via the magpie platform
\citep{kochetkova2024magpie} in a browser.
After reading instructions in German, participants are assigned to one of
4 counterbalancing lists (via URL \texttt{?list=N} parameter assigned by
Prolific; no random fallback in deployment).
They then proceed through 16 trials---8 critical and 8 fillers---in a
fully randomised order, followed by a post-test questionnaire eliciting
native language, age, gender, education level, and optional comments.

Each trial displays:
(a) a context paragraph establishing the communicative situation and the
\textsc{pc\textsubscript{prag}} level;
(b) the target sentence with the best-choice marker already in place in
the Mittelfeld;
(c) the consequent action clause;
(d) the two rating sliders in fixed order ($\widehat{g}$, $\widehat{A}$),
each labelled in German.
The ``Weiter'' (Next) button is disabled until both sliders have been
moved at least once, ensuring a complete response on every trial.

\subsection{Planned Analysis}

\subsubsection{Empirical plots}

For each dependent variable ($\widehat{g}$, $\widehat{A}$):

\begin{itemize}
  \item \textbf{Condition means}: a dot-and-whisker plot shows mean
        ratings per condition (the $2\times2$ of pc\textsubscript{prag}
        $\times$ $g_{\mathrm{implied}}$) with 95\% bootstrapped CIs.
        The expected pattern: $\widehat{g}$ and $\widehat{A}$ should
        be higher in the high-$g_{\mathrm{implied}}$ conditions
        (where the model assigns a stronger marker), and lower in the
        high-pc\textsubscript{prag} conditions (credibility discounting).
  \item \textbf{Marker identity}: because the best-choice marker may
        differ across items within the same condition, a secondary plot
        shows ratings grouped by actual marker identity, allowing
        inspection of whether marker type (rather than just condition)
        drives the effects.
  \item \textbf{pc\textsubscript{prop} covariate}: scatter plots display
        each DV as a function of the continuous norming rating per
        condition, testing whether higher propositional acceptance
        amplifies the $\widehat{g}$ and $\widehat{A}$ signals for
        strong markers.
\end{itemize}

\subsubsection{Bayesian mixed-effects regression}

Two Bayesian mixed-effects regression models are fitted via
\texttt{brms} \citep{burkner2017}, one per DV ($\widehat{g}$ and
$\widehat{A}$):

\begin{align}
y &\sim \mathrm{Beta}(\mu, \phi)
\quad \text{(ratings rescaled to }(0,1)\text{)} \\
\mathrm{logit}(\mu) &= \alpha
+ \beta_{g} \cdot \mathrm{g\_implied\_c}
+ \beta_{\mathrm{prag}} \cdot \mathrm{pc\_prag\_c}
+ \beta_{\mathrm{prop}} \cdot \mathrm{pc\_prop\_c}
+ \beta_{\mathrm{int}} \cdot \mathrm{g\_implied\_c}
  \times \mathrm{pc\_prag\_c}
+ (1 \mid \mathrm{submission\_id})
+ (1 \mid \mathrm{topic}),
\end{align}

where \texttt{g\_implied\_c} is the $g$ level from the corresponding
production condition ($-0.5$ = low, $+0.5$ = high), effectively
coding the marker's position on the production utility scale.
A secondary specification replaces \texttt{g\_implied\_c} with
actualy marker identity as an unordered factor, to examine
non-monotonic patterns and verify that marker rather than
condition label drives effects.

Model comparison via approximate leave-one-out cross-validation (LOO)
selects among additive, interaction, and marker-as-categorical
formulations.

\subsubsection{RSA listener parameter estimation}
\label{subsub:rsa_fit}

The two empirical rating means (per condition cell: 4 conditions) constrain
the listener parameters of the noisy-threshold RSA model.
Specifically, the model predicts:

\begin{align}
\mathbb{E}[\widehat{g} \mid u, \text{context}]
  &= \int g \cdot P(g \mid u, pc_{\mathrm{prag}})\, dg,
\label{eq:pred_g} \\
P(\widehat{A} \mid u, \text{context})
  &= \iiint \sigma(\eta_0 + \eta_g g - \eta_{pc}\, pc
     - \eta_{\mathrm{int}}\, g \cdot pc)\;
     P(pc_{\mathrm{prop}}, pc_{\mathrm{prag}}, g \mid u)\,
     d{pc_{\mathrm{prop}}}\,d{pc_{\mathrm{prag}}}\,dg,
\label{eq:adoption}
\end{align}

\noindent where the posterior $P(pc_{\mathrm{prop}}, pc_{\mathrm{prag}}, g
\mid u)$ is conditioned on the context's \textsc{pc\textsubscript{prag}}
level as a known prior (since it is explicitly given to the listener).

The threshold parameters $(\mu_1,\mu_2, \sigma_1,\sigma_2,
\lambda)$ together with the marker-specific costs $(c_u)$ are estimated
jointly from the production experiment data via the Stan model.
Given those threshold estimates, the listener parameters
$(\eta_0, \eta_g, \eta_{pc}, \eta_{\mathrm{int}})$ are then estimated by
minimising the weighted squared deviation between the model-predicted
means (Eqs.~\ref{eq:pred_g}--\ref{eq:adoption}) and the empirical rating
means from the comprehension experiment, or via a full joint Bayesian
fit over both experiments simultaneously.

\subsection{Model Predictions}

The noisy-threshold RSA model makes the following predictions for
$\widehat{g}$ and $\widehat{A}$, which serve as pre-registered hypotheses.
In all cases, ``strong-marker condition'' refers to conditions 1 and 3
(high $g_{\mathrm{implied}}$, i.e.\ the marker selected by the production
model under a high-$g$ context), and ``weak-marker condition'' to conditions
0 and 2.

\begin{enumerate}
  \item \textbf{Marker encodes inferred goal $\widehat{g}$}: $\widehat{g}$
        ratings should be higher in strong-marker conditions than in
        weak-marker conditions, across both pc\textsubscript{prag} levels.
        This is the core inversion prediction: if the production model
        correctly captures speaker behaviour, listeners should be able to
        decode $g$ from the marker.

  \item \textbf{Marker encodes adoption $\widehat{A}$}: $\widehat{A}$
        ratings should likewise be higher in strong-marker conditions,
        reflecting that higher inferred $g$ translates into greater
        adoption pressure downstream.

  \item \textbf{pc\textsubscript{prag} modulates $\widehat{A}$ but not
        $\widehat{g}$}: because $g$ is orthogonal to pc\textsubscript{prag}
        in the model's utility function, the effect of the marker on
        $\widehat{g}$ should be largely the same in low- and high-
        pc\textsubscript{prag} contexts.
        By contrast, $\widehat{A}$ should be reduced in high-
        pc\textsubscript{prag} contexts via the interaction term
        $\eta_{\mathrm{int}}\,g \cdot pc$: even a strong-marker signal
        for high $g$ is partly cancelled when controversy is high.

  \item \textbf{Credibility discounting}: for strong-marker conditions
        presented in a high-pc\textsubscript{prag} context (condition~3),
        $\widehat{A}$ should be \emph{lower} than for the same marker in
        a low-pc\textsubscript{prag} context (condition~1).
        This is the model's signature prediction distinguishing Bayesian
        inversion from a simpler marker-strength heuristic: a sceptical
        listener who knows controversy is high will discount the adoption
        implication of even the strongest marker.
        The $g_{\mathrm{implied}} \times \textsc{pc\textsubscript{prag}}$
        interaction on $\widehat{A}$ provides the empirical test.
\end{enumerate}

\end{document}
   % (inside a \section{Experiment} environment)

\documentclass[12pt]{article}
\usepackage{CSPArticleStyle}
\usepackage{csquotes}
\usepackage[margin=1in]{geometry}
\usepackage{booktabs}
\usepackage{parskip}
\usepackage{amsmath}
\usepackage{times}
\usepackage[T1]{fontenc}
\usepackage[utf8]{inputenc}
\usepackage{xcolor}
\usepackage{tabularray}
\usepackage{gb4e}
\noautomath

\title{Comprehension Experiment: Inferring Speaker Goal and Action Adoption
       from Consensus-Marker Use}
\author{}
\date{}

\begin{document}

%% ============================================================
%% When embedding in main.tex, delete from \documentclass to
%% \begin{document} and delete \end{document} at the bottom.
%% ============================================================

\section{Experiment: Consensus-Marker Comprehension}

\subsection{Overview}

The comprehension experiment is the listener-side counterpart of the
production experiment and tests whether listeners can invert the
speaker's marker choice to recover the latent context that motivated it.
Where the production experiment gave participants a context (fixing
\textsc{pc\textsubscript{prag}} and $g$) and asked which marker they
would use, the comprehension experiment gives participants a context
(fixing \textsc{pc\textsubscript{prag}}) together with a marker already
embedded in the sentence and asks what goal strength the speaker likely
had and how likely the listener is to adopt $q$.

The specific marker presented in each vignette is not chosen arbitrarily
from the three candidate markers. Our preregistered plan was to assign,
for each item$\times$condition cell, the \emph{Bayesian best-choice marker}:
the marker with the highest posterior probability under the cost-augmented
noisy-threshold RSA model fitted to the production data, given the context's
\textsc{pc\textsubscript{prag}} level and a specified $g$ level. This was
intended to mirror the production design closely and to ensure that the
comprehension stimuli reflect model-based marker use rather than arbitrary
marker$\times$item pairings.

Applying this strict modal-assignment rule to the fitted production model,
however, yielded a highly imbalanced stimulus set: most critical cells were
assigned the strongest marker, with very little variation across marker
categories. In addition, many of these modal decisions were supported only by
very small probability margins over the next-best alternative. Because such a
stimulus set would make the comprehension experiment substantially less
informative about marker-based listener inferences, we introduced a documented,
exploratory balancing step for stimulus construction. Starting from the same
posterior predictive probabilities, we selected a near-balanced assignment of
the three markers across the 32 critical cells while maximizing total posterior
support under the fitted production model. The strict modal assignment is
retained as the preregistered baseline, but the balanced assignment is used for
the actual comprehension materials.

\textsc{pc\textsubscript{prag}} is retained as an explicit context
manipulation (low vs.\ high), so that listeners can be expected to
have a well-defined prior on pragmatic controversy when interpreting
the marker. \textsc{pc\textsubscript{prop}} is not manipulated within
the experiment; it enters as a continuous topic-level covariate from
the norming study.

After reading each vignette, participants provide two slider ratings
operationalising the key listener-side quantities:
\begin{enumerate}
  \item \textbf{Inferred speaker goal strength} ($\widehat{g}$):
        how strongly does the speaker seem to want the listener to
        accept and act on $q$?
  \item \textbf{Adoption likelihood} ($\widehat{A}$):
        how likely is the listener to follow the recommended action $q$?
\end{enumerate}

\noindent Perceived controversy is \emph{not} elicited as a rating:
because \textsc{pc\textsubscript{prag}} is already manipulated as a
context factor, asking participants to rate controversy would be
redundant with respect to the manipulation check and circular with
respect to the research question. The theoretically central inference
to test is whether listeners can decode the speaker's \emph{goal} from
marker use, and whether that inference translates into adoption
behaviour.

The two slider measures, together with the \textsc{pc\textsubscript{prag}}
context and the \textsc{pc\textsubscript{prop}} covariate, provide the
data needed to fit the model's listener parameters
$(\eta_0, \eta_g, \eta_{pc}, \eta_{\mathrm{int}})$.

\subsection{Participants}

Participants will be recruited via Prolific and will be required to be
native speakers of German (language-at-home filter).
Data collection is planned after completion of the production experiment
(model fitting required to determine the best-choice markers).
A target of $N = 80$ participants is planned ($n = 20$ per list across
4 counterbalancing lists), each completing 16 trials (8 critical + 8
filler), matching the structure of the production experiment.

\subsection{Materials}

\subsubsection{Markers}

Each critical vignette is presented with one specific consensus marker
already embedded in the Mittelfeld of the sentence frame (see
example~(\ref{ex:frame_comp})):

\begin{exe}
  \ex \label{ex:frame_comp}
  Digitale Kompetenzen sind \textbf{[MARKER]}
  in der heutigen Arbeitswelt unverzichtbar.\\
  `Digital skills are \textbf{[MARKER]} indispensable in today's
  working world.'
\end{exe}

Unlike the production experiment, participants do \emph{not} choose the
marker; they read the sentence with the marker already in place and respond
to the two rating probes.

Critically, the marker assigned to each condition is not drawn arbitrarily
from the three candidates. Under the preregistered plan, for each of the four
production conditions (pc\textsubscript{prag} $\times$ $g$), the marker used in
comprehension was to be the \textbf{Bayesian best-choice marker}: the marker
$u^*$ with the highest posterior probability under the fitted
cost-augmented noisy-threshold RSA production model,
\[
u^* = \arg\max_{u} \; P(u \mid pc_{\mathrm{prag}}, g,
                           pc_{\mathrm{prop}\_c}).
\]
Because $pc_{\mathrm{prop}}$ varies continuously across items,
$u^*$ may differ per item even within the same
(pc\textsubscript{prag}, $g$) cell; the final stimulus list is therefore
constructed item-by-item once the production model has been fitted.
In practice, the fitted production model yielded a nearly degenerate modal map,
with most cells assigned to \textit{bekanntlich}. We therefore retained the
model-based logic of the procedure but replaced strict cellwise modal selection
with an uncertainty-aware balancing step: among the marker assignments with
high posterior support, we selected a near-balanced configuration across the 32
critical cells. This preserves the production--comprehension link while also
ensuring that the comprehension experiment contains sufficient marker variation
to test listener-side inferences.

All three markers remain candidates; the mapping from conditions to
markers is an empirical output of the production analysis, not a design
pre-commitment. In the final comprehension materials, we use the
\textbf{balanced model-based assignment}, rather than the strict modal
assignment, in order to maintain informative variation across marker types.

\subsubsection{Topic Items}

The same eight topic domains used in the norming study and the production
experiment are used here (Table~\ref{tab:items_comp}).
The 4-condition $\times$ 8-item Latin square (Section~\ref{subsub:counterbalancing_comp})
requires no additional items.

\begin{table}[ht]
\centering
\caption{Topic items: sentence frame for $p$ (with marker position [M])
         and consequent $q$.}
\label{tab:items_comp}
\small
\begin{tabular}{clll}
\toprule
\textbf{\#} & \textbf{Topic} & \textbf{Sentence frame} & \textbf{Consequent $q$} \\
\midrule
1 & Klimawandel &
  \textit{Der Klimawandel wird [M]} &
  auf Flugreisen verzichten \\
  & & \textit{durch menschliche Aktivit\"aten verursacht.} & \\[3pt]
2 & Ern\"ahrung &
  \textit{Eine \"uberwiegend pflanzliche Ern\"ahrung} &
  mehr pflanzliche Lebens- \\
  & & \textit{f\"ordert [M] die Gesundheit.} &
  mittel einbauen \\[3pt]
3 & Stadtverkehr &
  \textit{Ein ausgebautes Nahverkehrsnetz} &
  h\"aufiger auf \\
  & & \textit{entlastet [M] den Stadtverkehr} &
  \"offentliche Verkehrs- \\
  & & \textit{deutlich.} &
  mittel umsteigen \\[3pt]
4 & Digitalisierung &
  \textit{Digitale Kompetenzen sind [M] in} &
  regelm\"a\ss{}ig Weiter- \\
  & & \textit{der heutigen Arbeitswelt unverzichtbar.} &
  bildungsangebote nutzen \\[3pt]
5 & Schlaf &
  \textit{Ausreichend Schlaf ist [M] entscheidend} &
  auf feste Schlafzeiten \\
  & & \textit{f\"ur die kognitive Leistungsf\"ahigkeit.} &
  achten \\[3pt]
6 & Plastik &
  \textit{Einwegplastik schadet [M] den} &
  konsequent auf Einweg- \\
  & & \textit{Meeres\"okosystemen erheblich.} &
  plastik verzichten \\[3pt]
7 & Sport &
  \textit{Regelm\"a\ss{}ige k\"orperliche Bewegung} &
  regelm\"a\ss{}ig Sport in \\
  & & \textit{senkt [M] das Risiko f\"ur} &
  den Alltag integrieren \\
  & & \textit{Herz-Kreislauf-Erkrankungen deutlich.} & \\[3pt]
8 & Lokalkauf &
  \textit{Das Kaufen bei lokalen H\"andlern} &
  \"ofter bei lokalen \\
  & & \textit{st\"arkt [M] die regionale Wirtschaft} &
  Gesch\"aften einkaufen \\
  & & \textit{nachhaltig.} & \\
\bottomrule
\end{tabular}
\end{table}

\noindent [M] denotes the position of the pre-filled best-choice marker.

\subsubsection{Vignettes}

\paragraph{Structure.}
Each trial presents a vignette consisting of three parts:
(a) a context paragraph establishing the communicative situation and
fixing the \textsc{pc\textsubscript{prag}} level;
(b) the target sentence with the best-choice marker already embedded
in the Mittelfeld;
(c) the consequent action clause.
The two rating probes follow immediately below the vignette.

\paragraph{Propositional perceived controversy (\textsc{pc\textsubscript{prop}}).}
As in the production experiment, \textsc{pc\textsubscript{prop}} is not
manipulated within the comprehension experiment.
Topic-level propositional acceptance ratings from the norming study
(0--100 scale; e.g., \textit{Digitalisierung}: 85/100,
\textit{Lokalkauf}: 58/100) serve as a continuous, mean-centred
covariate \textsc{pc\textsubscript{prop}}$_c$ in analysis.
Context paragraphs present $p$ without editorial evaluation; variation in
listeners' individual assessments of propositional controversy is
captured by the norming-derived covariate.

\paragraph{Pragmatic perceived controversy (\textsc{pc\textsubscript{prag}}).}
\textsc{pc\textsubscript{prag}} is manipulated between item instances
using the same two context versions as in the production experiment:
the speaker's social circle either \emph{largely agrees} (low) or
\emph{holds divergent views} (high) about $p$.
For example, for the \textit{Digitalisierung} item:

\begin{itemize}
  \item $\textsc{pc}_\text{prag} = \text{low}$: \textit{Deine Kollegin spricht
        mit dir \"uber Berufsausbildung und Weiterbildung.
        In eurem Team teilen die meisten die Ansicht, dass digitale Kompetenzen
        in der heutigen Arbeitswelt unverzichtbar sind.}
  \item $\textsc{pc}_\text{prag} = \text{high}$: \textit{Deine Kollegin spricht
        mit dir \"uber Berufsausbildung und Weiterbildung.
        In eurem Team gehen die Meinungen dazu stark auseinander.}
\end{itemize}

The context is framed from the listener's perspective as a direct addressee,
mirroring the production study where participants adopt the speaker's
perspective.

\paragraph{Marker manipulation.}
The dependent measures of the production experiment (which marker is chosen)
become the independent variable here.
One of the three consensus markers is inserted into the Mittelfeld of the
target sentence.
A complete sample item for the \textit{Digitalisierung} topic
($\textsc{pc}_\text{prag} = \text{high}$, marker = \textit{bekanntlich};
norming $\textsc{pc}_\text{prop}$ rating: 85/100) is given in
(\ref{ex:comp_sample}):

\begin{exe}
  \ex \label{ex:comp_sample}
  \textit{Deine Kollegin spricht mit dir \"uber Berufsausbildung und
  Weiterbildung.
  In eurem Team gehen die Meinungen dazu stark auseinander.}

  \textit{Sie sagt zu dir:} Digitale Kompetenzen sind \textbf{bekanntlich}
  in der heutigen Arbeitswelt unverzichtbar.
  Deshalb solltest du regelm\"a\ss{}ig digitale Weiterbildungsangebote nutzen.

  \noindent\textit{[Rating probes follow.]}
\end{exe}

\subsection{Design}

\subsubsection{Factorial structure}

The comprehension experiment mirrors the production experiment's
$\mathbf{2}$~\textbf{(pc\textsubscript{prag}: low vs.\ high)}
$\mathbf{\times}$~$\mathbf{2}$~\textbf{($g_{\mathrm{implied}}$: low
vs.\ high) design}, yielding \textbf{4 conditions}.
The $g_{\mathrm{implied}}$ dimension is not manipulated via an explicit
goal instruction (as in production) but is instead conveyed entirely
through the choice of marker: for each condition, the marker presented
is the best-choice marker for the corresponding
(\textsc{pc\textsubscript{prag}},~$g$) cell from the production model.
\textsc{pc\textsubscript{prop}} enters as a continuous norming covariate.
With 8 topic items and a Latin square rotation across 4 lists, each
participant sees exactly \textbf{8 critical trials}---one item per
condition, each condition appearing twice---plus 8 filler trials,
for \textbf{16 trials} in total (Table~\ref{tab:design_comp}).

\begin{table}[ht]
\centering
\caption{%
  The $2\times2$ design (4 conditions), mirroring the production experiment.
  The marker presented in each condition is the Bayesian best-choice
  marker from the fitted production model; specific markers are
  determined after production data collection.}
\label{tab:design_comp}
\small
\begin{tabular}{ccll}
\toprule
\textbf{pc\textsubscript{prag}} &
\textbf{$g_{\mathrm{implied}}$} &
\textbf{Marker (TBD from production)} &
\textbf{Condition label} \\
\midrule
low  & low  & best-choice for $(\text{low}, g_{\text{low}})$  &
  (0) community agrees, weak signal \\
low  & high & best-choice for $(\text{low}, g_{\text{high}})$ &
  (1) community agrees, strong signal \\
high & low  & best-choice for $(\text{high}, g_{\text{low}})$ &
  (2) community divided, weak signal \\
high & high & best-choice for $(\text{high}, g_{\text{high}})$ &
  (3) community divided, strong signal \\
\bottomrule
\end{tabular}
\end{table}

\subsubsection{Counterbalancing}
\label{subsub:counterbalancing_comp}

The experiment uses the \textbf{same $4 \times 8$ Latin square} as the
production experiment: each of the 4 lists assigns item $i$
($i = 0,\ldots,7$) to condition $(i + l)\bmod 4$, where
$l \in \{0,1,2,3\}$ is the list number.
This guarantees that:
\begin{itemize}
  \item each participant sees every item exactly once;
  \item each participant sees each of the 4 conditions exactly twice
        (two items per condition), ensuring balanced coverage of the
        $2\times2$ design;
  \item across all 4 lists, each item appears equally often in each
        condition.
\end{itemize}

List assignment follows the same Prolific-and-URL strategy as the production
experiment: each Prolific study embeds \texttt{?list=N} (0--3) in the
experiment URL, capping recruitment at $n = 20$ participants per list
for a total of $N = 80$.

In addition to the 8 critical trials, each participant completed 8
\textbf{filler trials} using the same filler items as the production
experiment (train travel, reading light, cooking, coffee, vacation,
hydration, fresh air, noise), each presented with a randomly assigned
marker.
Fillers are excluded from analysis; they serve to dilute the critical
items and reduce demand characteristics.

\subsubsection{Dependent variables}
\label{subsub:dvs_comp}

Two slider ratings are collected after each vignette (0--100 continuous
scale with no visible numbers, labelled at the endpoints as shown):

\begin{enumerate}
  \item \textbf{Inferred speaker goal strength} (operationalises $\widehat{g}$):
  \begin{quote}
    \textit{Wie stark hast du den Eindruck, dass die Sprecherin m\"ochte,
    dass du ihre Empfehlung befolgst?}\\
    (`How strongly do you get the impression that the speaker wants you to
    follow their recommendation?')\\
    Endpoints: 0 = \textit{\"uberhaupt nicht}
    (`not at all') ---
    100 = \textit{sehr stark} (`very strongly').
  \end{quote}

  \item \textbf{Adoption likelihood of $q$} (operationalises $\widehat{A}$):
  \begin{quote}
    \textit{Wie wahrscheinlich ist es, dass du die empfohlene Handlung
    dann auch tats\"achlich umsetzt?}\\
    (`How likely is it that you will actually carry out the recommended
    action?')\\
    Endpoints: 0 = \textit{sehr unwahrscheinlich}
    (`very unlikely') ---
    100 = \textit{sehr wahrscheinlich} (`very likely').
  \end{quote}
\end{enumerate}

\noindent Perceived controversy is deliberately \emph{not} included as a
rating probe: \textsc{pc\textsubscript{prag}} is experimentally
manipulated via the context paragraph, so a controversy rating would
be redundant and risks obscuring the marker's contribution to inferred
goal strength.

Recorded per trial: \texttt{trial\_id}, \texttt{topic},
\texttt{is\_filler}, \texttt{condition\_index}, \texttt{marker},
\texttt{pc\_prag}, \texttt{g\_implied}, \texttt{rating\_g},
\texttt{rating\_adopt}, \texttt{rt} (ms); and once per participant:
\texttt{list\_num}.
The continuous \texttt{pc\_prop\_rating} covariate (topic-level norming
mean) is joined to the trial data during analysis.

\subsection{Procedure}

Participants complete the experiment online via the magpie platform
\citep{kochetkova2024magpie} in a browser.
After reading instructions in German, participants are assigned to one of
4 counterbalancing lists (via URL \texttt{?list=N} parameter assigned by
Prolific; no random fallback in deployment).
They then proceed through 16 trials---8 critical and 8 fillers---in a
fully randomised order, followed by a post-test questionnaire eliciting
native language, age, gender, education level, and optional comments.

Each trial displays:
(a) a context paragraph establishing the communicative situation and the
\textsc{pc\textsubscript{prag}} level;
(b) the target sentence with the best-choice marker already in place in
the Mittelfeld;
(c) the consequent action clause;
(d) the two rating sliders in fixed order ($\widehat{g}$, $\widehat{A}$),
each labelled in German.
The ``Weiter'' (Next) button is disabled until both sliders have been
moved at least once, ensuring a complete response on every trial.

\subsection{Planned Analysis}

\subsubsection{Empirical plots}

For each dependent variable ($\widehat{g}$, $\widehat{A}$):

\begin{itemize}
  \item \textbf{Condition means}: a dot-and-whisker plot shows mean
        ratings per condition (the $2\times2$ of pc\textsubscript{prag}
        $\times$ $g_{\mathrm{implied}}$) with 95\% bootstrapped CIs.
        The expected pattern: $\widehat{g}$ and $\widehat{A}$ should
        be higher in the high-$g_{\mathrm{implied}}$ conditions
        (where the model assigns a stronger marker), and lower in the
        high-pc\textsubscript{prag} conditions (credibility discounting).
  \item \textbf{Marker identity}: because the best-choice marker may
        differ across items within the same condition, a secondary plot
        shows ratings grouped by actual marker identity, allowing
        inspection of whether marker type (rather than just condition)
        drives the effects.
  \item \textbf{pc\textsubscript{prop} covariate}: scatter plots display
        each DV as a function of the continuous norming rating per
        condition, testing whether higher propositional acceptance
        amplifies the $\widehat{g}$ and $\widehat{A}$ signals for
        strong markers.
\end{itemize}

\subsubsection{Bayesian mixed-effects regression}

Two Bayesian mixed-effects regression models are fitted via
\texttt{brms} \citep{burkner2017}, one per DV ($\widehat{g}$ and
$\widehat{A}$):

\begin{align}
y &\sim \mathrm{Beta}(\mu, \phi)
\quad \text{(ratings rescaled to }(0,1)\text{)} \\
\mathrm{logit}(\mu) &= \alpha
+ \beta_{g} \cdot \mathrm{g\_implied\_c}
+ \beta_{\mathrm{prag}} \cdot \mathrm{pc\_prag\_c}
+ \beta_{\mathrm{prop}} \cdot \mathrm{pc\_prop\_c}
+ \beta_{\mathrm{int}} \cdot \mathrm{g\_implied\_c}
  \times \mathrm{pc\_prag\_c}
+ (1 \mid \mathrm{submission\_id})
+ (1 \mid \mathrm{topic}),
\end{align}

where \texttt{g\_implied\_c} is the $g$ level from the corresponding
production condition ($-0.5$ = low, $+0.5$ = high), effectively
coding the marker's position on the production utility scale.
A secondary specification replaces \texttt{g\_implied\_c} with
actualy marker identity as an unordered factor, to examine
non-monotonic patterns and verify that marker rather than
condition label drives effects.

Model comparison via approximate leave-one-out cross-validation (LOO)
selects among additive, interaction, and marker-as-categorical
formulations.

\subsubsection{RSA listener parameter estimation}
\label{subsub:rsa_fit}

The two empirical rating means (per condition cell: 4 conditions) constrain
the listener parameters of the noisy-threshold RSA model.
Specifically, the model predicts:

\begin{align}
\mathbb{E}[\widehat{g} \mid u, \text{context}]
  &= \int g \cdot P(g \mid u, pc_{\mathrm{prag}})\, dg,
\label{eq:pred_g} \\
P(\widehat{A} \mid u, \text{context})
  &= \iiint \sigma(\eta_0 + \eta_g g - \eta_{pc}\, pc
     - \eta_{\mathrm{int}}\, g \cdot pc)\;
     P(pc_{\mathrm{prop}}, pc_{\mathrm{prag}}, g \mid u)\,
     d{pc_{\mathrm{prop}}}\,d{pc_{\mathrm{prag}}}\,dg,
\label{eq:adoption}
\end{align}

\noindent where the posterior $P(pc_{\mathrm{prop}}, pc_{\mathrm{prag}}, g
\mid u)$ is conditioned on the context's \textsc{pc\textsubscript{prag}}
level as a known prior (since it is explicitly given to the listener).

The threshold parameters $(\mu_1,\mu_2, \sigma_1,\sigma_2,
\lambda)$ together with the marker-specific costs $(c_u)$ are estimated
jointly from the production experiment data via the Stan model.
Given those threshold estimates, the listener parameters
$(\eta_0, \eta_g, \eta_{pc}, \eta_{\mathrm{int}})$ are then estimated by
minimising the weighted squared deviation between the model-predicted
means (Eqs.~\ref{eq:pred_g}--\ref{eq:adoption}) and the empirical rating
means from the comprehension experiment, or via a full joint Bayesian
fit over both experiments simultaneously.

\subsection{Model Predictions}

The noisy-threshold RSA model makes the following predictions for
$\widehat{g}$ and $\widehat{A}$, which serve as pre-registered hypotheses.
In all cases, ``strong-marker condition'' refers to conditions 1 and 3
(high $g_{\mathrm{implied}}$, i.e.\ the marker selected by the production
model under a high-$g$ context), and ``weak-marker condition'' to conditions
0 and 2.

\begin{enumerate}
  \item \textbf{Marker encodes inferred goal $\widehat{g}$}: $\widehat{g}$
        ratings should be higher in strong-marker conditions than in
        weak-marker conditions, across both pc\textsubscript{prag} levels.
        This is the core inversion prediction: if the production model
        correctly captures speaker behaviour, listeners should be able to
        decode $g$ from the marker.

  \item \textbf{Marker encodes adoption $\widehat{A}$}: $\widehat{A}$
        ratings should likewise be higher in strong-marker conditions,
        reflecting that higher inferred $g$ translates into greater
        adoption pressure downstream.

  \item \textbf{pc\textsubscript{prag} modulates $\widehat{A}$ but not
        $\widehat{g}$}: because $g$ is orthogonal to pc\textsubscript{prag}
        in the model's utility function, the effect of the marker on
        $\widehat{g}$ should be largely the same in low- and high-
        pc\textsubscript{prag} contexts.
        By contrast, $\widehat{A}$ should be reduced in high-
        pc\textsubscript{prag} contexts via the interaction term
        $\eta_{\mathrm{int}}\,g \cdot pc$: even a strong-marker signal
        for high $g$ is partly cancelled when controversy is high.

  \item \textbf{Credibility discounting}: for strong-marker conditions
        presented in a high-pc\textsubscript{prag} context (condition~3),
        $\widehat{A}$ should be \emph{lower} than for the same marker in
        a low-pc\textsubscript{prag} context (condition~1).
        This is the model's signature prediction distinguishing Bayesian
        inversion from a simpler marker-strength heuristic: a sceptical
        listener who knows controversy is high will discount the adoption
        implication of even the strongest marker.
        The $g_{\mathrm{implied}} \times \textsc{pc\textsubscript{prag}}$
        interaction on $\widehat{A}$ provides the empirical test.
\end{enumerate}

\end{document}
   % (inside a \section{Experiment} environment)

\documentclass[12pt]{article}
\usepackage{CSPArticleStyle}
\usepackage{csquotes}
\usepackage[margin=1in]{geometry}
\usepackage{booktabs}
\usepackage{parskip}
\usepackage{amsmath}
\usepackage{times}
\usepackage[T1]{fontenc}
\usepackage[utf8]{inputenc}
\usepackage{xcolor}
\usepackage{tabularray}
\usepackage{gb4e}
\noautomath

\title{Comprehension Experiment: Inferring Speaker Goal and Action Adoption
       from Consensus-Marker Use}
\author{}
\date{}

\begin{document}

%% ============================================================
%% When embedding in main.tex, delete from \documentclass to
%% \begin{document} and delete \end{document} at the bottom.
%% ============================================================

\section{Experiment: Consensus-Marker Comprehension}

\subsection{Overview}

The comprehension experiment is the listener-side counterpart of the
production experiment and tests whether listeners can invert the
speaker's marker choice to recover the latent context that motivated it.
Where the production experiment gave participants a context (fixing
\textsc{pc\textsubscript{prag}} and $g$) and asked which marker they
would use, the comprehension experiment gives participants a context
(fixing \textsc{pc\textsubscript{prag}}) together with a marker already
embedded in the sentence and asks what goal strength the speaker likely
had and how likely the listener is to adopt $q$.

The specific marker presented in each vignette is not chosen arbitrarily
from the three candidate markers. Our preregistered plan was to assign,
for each item$\times$condition cell, the \emph{Bayesian best-choice marker}:
the marker with the highest posterior probability under the cost-augmented
noisy-threshold RSA model fitted to the production data, given the context's
\textsc{pc\textsubscript{prag}} level and a specified $g$ level. This was
intended to mirror the production design closely and to ensure that the
comprehension stimuli reflect model-based marker use rather than arbitrary
marker$\times$item pairings.

Applying this strict modal-assignment rule to the fitted production model,
however, yielded a highly imbalanced stimulus set: most critical cells were
assigned the strongest marker, with very little variation across marker
categories. In addition, many of these modal decisions were supported only by
very small probability margins over the next-best alternative. Because such a
stimulus set would make the comprehension experiment substantially less
informative about marker-based listener inferences, we introduced a documented,
exploratory balancing step for stimulus construction. Starting from the same
posterior predictive probabilities, we selected a near-balanced assignment of
the three markers across the 32 critical cells while maximizing total posterior
support under the fitted production model. The strict modal assignment is
retained as the preregistered baseline, but the balanced assignment is used for
the actual comprehension materials.

\textsc{pc\textsubscript{prag}} is retained as an explicit context
manipulation (low vs.\ high), so that listeners can be expected to
have a well-defined prior on pragmatic controversy when interpreting
the marker. \textsc{pc\textsubscript{prop}} is not manipulated within
the experiment; it enters as a continuous topic-level covariate from
the norming study.

After reading each vignette, participants provide two slider ratings
operationalising the key listener-side quantities:
\begin{enumerate}
  \item \textbf{Inferred speaker goal strength} ($\widehat{g}$):
        how strongly does the speaker seem to want the listener to
        accept and act on $q$?
  \item \textbf{Adoption likelihood} ($\widehat{A}$):
        how likely is the listener to follow the recommended action $q$?
\end{enumerate}

\noindent Perceived controversy is \emph{not} elicited as a rating:
because \textsc{pc\textsubscript{prag}} is already manipulated as a
context factor, asking participants to rate controversy would be
redundant with respect to the manipulation check and circular with
respect to the research question. The theoretically central inference
to test is whether listeners can decode the speaker's \emph{goal} from
marker use, and whether that inference translates into adoption
behaviour.

The two slider measures, together with the \textsc{pc\textsubscript{prag}}
context and the \textsc{pc\textsubscript{prop}} covariate, provide the
data needed to fit the model's listener parameters
$(\eta_0, \eta_g, \eta_{pc}, \eta_{\mathrm{int}})$.

\subsection{Participants}

Participants will be recruited via Prolific and will be required to be
native speakers of German (language-at-home filter).
Data collection is planned after completion of the production experiment
(model fitting required to determine the best-choice markers).
A target of $N = 80$ participants is planned ($n = 20$ per list across
4 counterbalancing lists), each completing 16 trials (8 critical + 8
filler), matching the structure of the production experiment.

\subsection{Materials}

\subsubsection{Markers}

Each critical vignette is presented with one specific consensus marker
already embedded in the Mittelfeld of the sentence frame (see
example~(\ref{ex:frame_comp})):

\begin{exe}
  \ex \label{ex:frame_comp}
  Digitale Kompetenzen sind \textbf{[MARKER]}
  in der heutigen Arbeitswelt unverzichtbar.\\
  `Digital skills are \textbf{[MARKER]} indispensable in today's
  working world.'
\end{exe}

Unlike the production experiment, participants do \emph{not} choose the
marker; they read the sentence with the marker already in place and respond
to the two rating probes.

Critically, the marker assigned to each condition is not drawn arbitrarily
from the three candidates. Under the preregistered plan, for each of the four
production conditions (pc\textsubscript{prag} $\times$ $g$), the marker used in
comprehension was to be the \textbf{Bayesian best-choice marker}: the marker
$u^*$ with the highest posterior probability under the fitted
cost-augmented noisy-threshold RSA production model,
\[
u^* = \arg\max_{u} \; P(u \mid pc_{\mathrm{prag}}, g,
                           pc_{\mathrm{prop}\_c}).
\]
Because $pc_{\mathrm{prop}}$ varies continuously across items,
$u^*$ may differ per item even within the same
(pc\textsubscript{prag}, $g$) cell; the final stimulus list is therefore
constructed item-by-item once the production model has been fitted.
In practice, the fitted production model yielded a nearly degenerate modal map,
with most cells assigned to \textit{bekanntlich}. We therefore retained the
model-based logic of the procedure but replaced strict cellwise modal selection
with an uncertainty-aware balancing step: among the marker assignments with
high posterior support, we selected a near-balanced configuration across the 32
critical cells. This preserves the production--comprehension link while also
ensuring that the comprehension experiment contains sufficient marker variation
to test listener-side inferences.

All three markers remain candidates; the mapping from conditions to
markers is an empirical output of the production analysis, not a design
pre-commitment. In the final comprehension materials, we use the
\textbf{balanced model-based assignment}, rather than the strict modal
assignment, in order to maintain informative variation across marker types.

\subsubsection{Topic Items}

The same eight topic domains used in the norming study and the production
experiment are used here (Table~\ref{tab:items_comp}).
The 4-condition $\times$ 8-item Latin square (Section~\ref{subsub:counterbalancing_comp})
requires no additional items.

\begin{table}[ht]
\centering
\caption{Topic items: sentence frame for $p$ (with marker position [M])
         and consequent $q$.}
\label{tab:items_comp}
\small
\begin{tabular}{clll}
\toprule
\textbf{\#} & \textbf{Topic} & \textbf{Sentence frame} & \textbf{Consequent $q$} \\
\midrule
1 & Klimawandel &
  \textit{Der Klimawandel wird [M]} &
  auf Flugreisen verzichten \\
  & & \textit{durch menschliche Aktivit\"aten verursacht.} & \\[3pt]
2 & Ern\"ahrung &
  \textit{Eine \"uberwiegend pflanzliche Ern\"ahrung} &
  mehr pflanzliche Lebens- \\
  & & \textit{f\"ordert [M] die Gesundheit.} &
  mittel einbauen \\[3pt]
3 & Stadtverkehr &
  \textit{Ein ausgebautes Nahverkehrsnetz} &
  h\"aufiger auf \\
  & & \textit{entlastet [M] den Stadtverkehr} &
  \"offentliche Verkehrs- \\
  & & \textit{deutlich.} &
  mittel umsteigen \\[3pt]
4 & Digitalisierung &
  \textit{Digitale Kompetenzen sind [M] in} &
  regelm\"a\ss{}ig Weiter- \\
  & & \textit{der heutigen Arbeitswelt unverzichtbar.} &
  bildungsangebote nutzen \\[3pt]
5 & Schlaf &
  \textit{Ausreichend Schlaf ist [M] entscheidend} &
  auf feste Schlafzeiten \\
  & & \textit{f\"ur die kognitive Leistungsf\"ahigkeit.} &
  achten \\[3pt]
6 & Plastik &
  \textit{Einwegplastik schadet [M] den} &
  konsequent auf Einweg- \\
  & & \textit{Meeres\"okosystemen erheblich.} &
  plastik verzichten \\[3pt]
7 & Sport &
  \textit{Regelm\"a\ss{}ige k\"orperliche Bewegung} &
  regelm\"a\ss{}ig Sport in \\
  & & \textit{senkt [M] das Risiko f\"ur} &
  den Alltag integrieren \\
  & & \textit{Herz-Kreislauf-Erkrankungen deutlich.} & \\[3pt]
8 & Lokalkauf &
  \textit{Das Kaufen bei lokalen H\"andlern} &
  \"ofter bei lokalen \\
  & & \textit{st\"arkt [M] die regionale Wirtschaft} &
  Gesch\"aften einkaufen \\
  & & \textit{nachhaltig.} & \\
\bottomrule
\end{tabular}
\end{table}

\noindent [M] denotes the position of the pre-filled best-choice marker.

\subsubsection{Vignettes}

\paragraph{Structure.}
Each trial presents a vignette consisting of three parts:
(a) a context paragraph establishing the communicative situation and
fixing the \textsc{pc\textsubscript{prag}} level;
(b) the target sentence with the best-choice marker already embedded
in the Mittelfeld;
(c) the consequent action clause.
The two rating probes follow immediately below the vignette.

\paragraph{Propositional perceived controversy (\textsc{pc\textsubscript{prop}}).}
As in the production experiment, \textsc{pc\textsubscript{prop}} is not
manipulated within the comprehension experiment.
Topic-level propositional acceptance ratings from the norming study
(0--100 scale; e.g., \textit{Digitalisierung}: 85/100,
\textit{Lokalkauf}: 58/100) serve as a continuous, mean-centred
covariate \textsc{pc\textsubscript{prop}}$_c$ in analysis.
Context paragraphs present $p$ without editorial evaluation; variation in
listeners' individual assessments of propositional controversy is
captured by the norming-derived covariate.

\paragraph{Pragmatic perceived controversy (\textsc{pc\textsubscript{prag}}).}
\textsc{pc\textsubscript{prag}} is manipulated between item instances
using the same two context versions as in the production experiment:
the speaker's social circle either \emph{largely agrees} (low) or
\emph{holds divergent views} (high) about $p$.
For example, for the \textit{Digitalisierung} item:

\begin{itemize}
  \item $\textsc{pc}_\text{prag} = \text{low}$: \textit{Deine Kollegin spricht
        mit dir \"uber Berufsausbildung und Weiterbildung.
        In eurem Team teilen die meisten die Ansicht, dass digitale Kompetenzen
        in der heutigen Arbeitswelt unverzichtbar sind.}
  \item $\textsc{pc}_\text{prag} = \text{high}$: \textit{Deine Kollegin spricht
        mit dir \"uber Berufsausbildung und Weiterbildung.
        In eurem Team gehen die Meinungen dazu stark auseinander.}
\end{itemize}

The context is framed from the listener's perspective as a direct addressee,
mirroring the production study where participants adopt the speaker's
perspective.

\paragraph{Marker manipulation.}
The dependent measures of the production experiment (which marker is chosen)
become the independent variable here.
One of the three consensus markers is inserted into the Mittelfeld of the
target sentence.
A complete sample item for the \textit{Digitalisierung} topic
($\textsc{pc}_\text{prag} = \text{high}$, marker = \textit{bekanntlich};
norming $\textsc{pc}_\text{prop}$ rating: 85/100) is given in
(\ref{ex:comp_sample}):

\begin{exe}
  \ex \label{ex:comp_sample}
  \textit{Deine Kollegin spricht mit dir \"uber Berufsausbildung und
  Weiterbildung.
  In eurem Team gehen die Meinungen dazu stark auseinander.}

  \textit{Sie sagt zu dir:} Digitale Kompetenzen sind \textbf{bekanntlich}
  in der heutigen Arbeitswelt unverzichtbar.
  Deshalb solltest du regelm\"a\ss{}ig digitale Weiterbildungsangebote nutzen.

  \noindent\textit{[Rating probes follow.]}
\end{exe}

\subsection{Design}

\subsubsection{Factorial structure}

The comprehension experiment mirrors the production experiment's
$\mathbf{2}$~\textbf{(pc\textsubscript{prag}: low vs.\ high)}
$\mathbf{\times}$~$\mathbf{2}$~\textbf{($g_{\mathrm{implied}}$: low
vs.\ high) design}, yielding \textbf{4 conditions}.
The $g_{\mathrm{implied}}$ dimension is not manipulated via an explicit
goal instruction (as in production) but is instead conveyed entirely
through the choice of marker: for each condition, the marker presented
is the best-choice marker for the corresponding
(\textsc{pc\textsubscript{prag}},~$g$) cell from the production model.
\textsc{pc\textsubscript{prop}} enters as a continuous norming covariate.
With 8 topic items and a Latin square rotation across 4 lists, each
participant sees exactly \textbf{8 critical trials}---one item per
condition, each condition appearing twice---plus 8 filler trials,
for \textbf{16 trials} in total (Table~\ref{tab:design_comp}).

\begin{table}[ht]
\centering
\caption{%
  The $2\times2$ design (4 conditions), mirroring the production experiment.
  The marker presented in each condition is the Bayesian best-choice
  marker from the fitted production model; specific markers are
  determined after production data collection.}
\label{tab:design_comp}
\small
\begin{tabular}{ccll}
\toprule
\textbf{pc\textsubscript{prag}} &
\textbf{$g_{\mathrm{implied}}$} &
\textbf{Marker (TBD from production)} &
\textbf{Condition label} \\
\midrule
low  & low  & best-choice for $(\text{low}, g_{\text{low}})$  &
  (0) community agrees, weak signal \\
low  & high & best-choice for $(\text{low}, g_{\text{high}})$ &
  (1) community agrees, strong signal \\
high & low  & best-choice for $(\text{high}, g_{\text{low}})$ &
  (2) community divided, weak signal \\
high & high & best-choice for $(\text{high}, g_{\text{high}})$ &
  (3) community divided, strong signal \\
\bottomrule
\end{tabular}
\end{table}

\subsubsection{Counterbalancing}
\label{subsub:counterbalancing_comp}

The experiment uses the \textbf{same $4 \times 8$ Latin square} as the
production experiment: each of the 4 lists assigns item $i$
($i = 0,\ldots,7$) to condition $(i + l)\bmod 4$, where
$l \in \{0,1,2,3\}$ is the list number.
This guarantees that:
\begin{itemize}
  \item each participant sees every item exactly once;
  \item each participant sees each of the 4 conditions exactly twice
        (two items per condition), ensuring balanced coverage of the
        $2\times2$ design;
  \item across all 4 lists, each item appears equally often in each
        condition.
\end{itemize}

List assignment follows the same Prolific-and-URL strategy as the production
experiment: each Prolific study embeds \texttt{?list=N} (0--3) in the
experiment URL, capping recruitment at $n = 20$ participants per list
for a total of $N = 80$.

In addition to the 8 critical trials, each participant completed 8
\textbf{filler trials} using the same filler items as the production
experiment (train travel, reading light, cooking, coffee, vacation,
hydration, fresh air, noise), each presented with a randomly assigned
marker.
Fillers are excluded from analysis; they serve to dilute the critical
items and reduce demand characteristics.

\subsubsection{Dependent variables}
\label{subsub:dvs_comp}

Two slider ratings are collected after each vignette (0--100 continuous
scale with no visible numbers, labelled at the endpoints as shown):

\begin{enumerate}
  \item \textbf{Inferred speaker goal strength} (operationalises $\widehat{g}$):
  \begin{quote}
    \textit{Wie stark hast du den Eindruck, dass die Sprecherin m\"ochte,
    dass du ihre Empfehlung befolgst?}\\
    (`How strongly do you get the impression that the speaker wants you to
    follow their recommendation?')\\
    Endpoints: 0 = \textit{\"uberhaupt nicht}
    (`not at all') ---
    100 = \textit{sehr stark} (`very strongly').
  \end{quote}

  \item \textbf{Adoption likelihood of $q$} (operationalises $\widehat{A}$):
  \begin{quote}
    \textit{Wie wahrscheinlich ist es, dass du die empfohlene Handlung
    dann auch tats\"achlich umsetzt?}\\
    (`How likely is it that you will actually carry out the recommended
    action?')\\
    Endpoints: 0 = \textit{sehr unwahrscheinlich}
    (`very unlikely') ---
    100 = \textit{sehr wahrscheinlich} (`very likely').
  \end{quote}
\end{enumerate}

\noindent Perceived controversy is deliberately \emph{not} included as a
rating probe: \textsc{pc\textsubscript{prag}} is experimentally
manipulated via the context paragraph, so a controversy rating would
be redundant and risks obscuring the marker's contribution to inferred
goal strength.

Recorded per trial: \texttt{trial\_id}, \texttt{topic},
\texttt{is\_filler}, \texttt{condition\_index}, \texttt{marker},
\texttt{pc\_prag}, \texttt{g\_implied}, \texttt{rating\_g},
\texttt{rating\_adopt}, \texttt{rt} (ms); and once per participant:
\texttt{list\_num}.
The continuous \texttt{pc\_prop\_rating} covariate (topic-level norming
mean) is joined to the trial data during analysis.

\subsection{Procedure}

Participants complete the experiment online via the magpie platform
\citep{kochetkova2024magpie} in a browser.
After reading instructions in German, participants are assigned to one of
4 counterbalancing lists (via URL \texttt{?list=N} parameter assigned by
Prolific; no random fallback in deployment).
They then proceed through 16 trials---8 critical and 8 fillers---in a
fully randomised order, followed by a post-test questionnaire eliciting
native language, age, gender, education level, and optional comments.

Each trial displays:
(a) a context paragraph establishing the communicative situation and the
\textsc{pc\textsubscript{prag}} level;
(b) the target sentence with the best-choice marker already in place in
the Mittelfeld;
(c) the consequent action clause;
(d) the two rating sliders in fixed order ($\widehat{g}$, $\widehat{A}$),
each labelled in German.
The ``Weiter'' (Next) button is disabled until both sliders have been
moved at least once, ensuring a complete response on every trial.

\subsection{Planned Analysis}

\subsubsection{Empirical plots}

For each dependent variable ($\widehat{g}$, $\widehat{A}$):

\begin{itemize}
  \item \textbf{Condition means}: a dot-and-whisker plot shows mean
        ratings per condition (the $2\times2$ of pc\textsubscript{prag}
        $\times$ $g_{\mathrm{implied}}$) with 95\% bootstrapped CIs.
        The expected pattern: $\widehat{g}$ and $\widehat{A}$ should
        be higher in the high-$g_{\mathrm{implied}}$ conditions
        (where the model assigns a stronger marker), and lower in the
        high-pc\textsubscript{prag} conditions (credibility discounting).
  \item \textbf{Marker identity}: because the best-choice marker may
        differ across items within the same condition, a secondary plot
        shows ratings grouped by actual marker identity, allowing
        inspection of whether marker type (rather than just condition)
        drives the effects.
  \item \textbf{pc\textsubscript{prop} covariate}: scatter plots display
        each DV as a function of the continuous norming rating per
        condition, testing whether higher propositional acceptance
        amplifies the $\widehat{g}$ and $\widehat{A}$ signals for
        strong markers.
\end{itemize}

\subsubsection{Bayesian mixed-effects regression}

Two Bayesian mixed-effects regression models are fitted via
\texttt{brms} \citep{burkner2017}, one per DV ($\widehat{g}$ and
$\widehat{A}$):

\begin{align}
y &\sim \mathrm{Beta}(\mu, \phi)
\quad \text{(ratings rescaled to }(0,1)\text{)} \\
\mathrm{logit}(\mu) &= \alpha
+ \beta_{g} \cdot \mathrm{g\_implied\_c}
+ \beta_{\mathrm{prag}} \cdot \mathrm{pc\_prag\_c}
+ \beta_{\mathrm{prop}} \cdot \mathrm{pc\_prop\_c}
+ \beta_{\mathrm{int}} \cdot \mathrm{g\_implied\_c}
  \times \mathrm{pc\_prag\_c}
+ (1 \mid \mathrm{submission\_id})
+ (1 \mid \mathrm{topic}),
\end{align}

where \texttt{g\_implied\_c} is the $g$ level from the corresponding
production condition ($-0.5$ = low, $+0.5$ = high), effectively
coding the marker's position on the production utility scale.
A secondary specification replaces \texttt{g\_implied\_c} with
actualy marker identity as an unordered factor, to examine
non-monotonic patterns and verify that marker rather than
condition label drives effects.

Model comparison via approximate leave-one-out cross-validation (LOO)
selects among additive, interaction, and marker-as-categorical
formulations.

\subsubsection{RSA listener parameter estimation}
\label{subsub:rsa_fit}

The two empirical rating means (per condition cell: 4 conditions) constrain
the listener parameters of the noisy-threshold RSA model.
Specifically, the model predicts:

\begin{align}
\mathbb{E}[\widehat{g} \mid u, \text{context}]
  &= \int g \cdot P(g \mid u, pc_{\mathrm{prag}})\, dg,
\label{eq:pred_g} \\
P(\widehat{A} \mid u, \text{context})
  &= \iiint \sigma(\eta_0 + \eta_g g - \eta_{pc}\, pc
     - \eta_{\mathrm{int}}\, g \cdot pc)\;
     P(pc_{\mathrm{prop}}, pc_{\mathrm{prag}}, g \mid u)\,
     d{pc_{\mathrm{prop}}}\,d{pc_{\mathrm{prag}}}\,dg,
\label{eq:adoption}
\end{align}

\noindent where the posterior $P(pc_{\mathrm{prop}}, pc_{\mathrm{prag}}, g
\mid u)$ is conditioned on the context's \textsc{pc\textsubscript{prag}}
level as a known prior (since it is explicitly given to the listener).

The threshold parameters $(\mu_1,\mu_2, \sigma_1,\sigma_2,
\lambda)$ together with the marker-specific costs $(c_u)$ are estimated
jointly from the production experiment data via the Stan model.
Given those threshold estimates, the listener parameters
$(\eta_0, \eta_g, \eta_{pc}, \eta_{\mathrm{int}})$ are then estimated by
minimising the weighted squared deviation between the model-predicted
means (Eqs.~\ref{eq:pred_g}--\ref{eq:adoption}) and the empirical rating
means from the comprehension experiment, or via a full joint Bayesian
fit over both experiments simultaneously.

\subsection{Model Predictions}

The noisy-threshold RSA model makes the following predictions for
$\widehat{g}$ and $\widehat{A}$, which serve as pre-registered hypotheses.
In all cases, ``strong-marker condition'' refers to conditions 1 and 3
(high $g_{\mathrm{implied}}$, i.e.\ the marker selected by the production
model under a high-$g$ context), and ``weak-marker condition'' to conditions
0 and 2.

\begin{enumerate}
  \item \textbf{Marker encodes inferred goal $\widehat{g}$}: $\widehat{g}$
        ratings should be higher in strong-marker conditions than in
        weak-marker conditions, across both pc\textsubscript{prag} levels.
        This is the core inversion prediction: if the production model
        correctly captures speaker behaviour, listeners should be able to
        decode $g$ from the marker.

  \item \textbf{Marker encodes adoption $\widehat{A}$}: $\widehat{A}$
        ratings should likewise be higher in strong-marker conditions,
        reflecting that higher inferred $g$ translates into greater
        adoption pressure downstream.

  \item \textbf{pc\textsubscript{prag} modulates $\widehat{A}$ but not
        $\widehat{g}$}: because $g$ is orthogonal to pc\textsubscript{prag}
        in the model's utility function, the effect of the marker on
        $\widehat{g}$ should be largely the same in low- and high-
        pc\textsubscript{prag} contexts.
        By contrast, $\widehat{A}$ should be reduced in high-
        pc\textsubscript{prag} contexts via the interaction term
        $\eta_{\mathrm{int}}\,g \cdot pc$: even a strong-marker signal
        for high $g$ is partly cancelled when controversy is high.

  \item \textbf{Credibility discounting}: for strong-marker conditions
        presented in a high-pc\textsubscript{prag} context (condition~3),
        $\widehat{A}$ should be \emph{lower} than for the same marker in
        a low-pc\textsubscript{prag} context (condition~1).
        This is the model's signature prediction distinguishing Bayesian
        inversion from a simpler marker-strength heuristic: a sceptical
        listener who knows controversy is high will discount the adoption
        implication of even the strongest marker.
        The $g_{\mathrm{implied}} \times \textsc{pc\textsubscript{prag}}$
        interaction on $\widehat{A}$ provides the empirical test.
\end{enumerate}

\end{document}
   % (inside a \section{Experiment} environment)

\documentclass[12pt]{article}
\usepackage{CSPArticleStyle}
\usepackage{csquotes}
\usepackage[margin=1in]{geometry}
\usepackage{booktabs}
\usepackage{parskip}
\usepackage{amsmath}
\usepackage{times}
\usepackage[T1]{fontenc}
\usepackage[utf8]{inputenc}
\usepackage{xcolor}
\usepackage{tabularray}
\usepackage{gb4e}
\noautomath

\title{Comprehension Experiment: Inferring Speaker Goal and Action Adoption
       from Consensus-Marker Use}
\author{}
\date{}

\begin{document}

%% ============================================================
%% When embedding in main.tex, delete from \documentclass to
%% \begin{document} and delete \end{document} at the bottom.
%% ============================================================

\section{Experiment: Consensus-Marker Comprehension}

\subsection{Overview}

The comprehension experiment is the listener-side counterpart of the
production experiment and tests whether listeners can invert the
speaker's marker choice to recover the latent context that motivated it.
Where the production experiment gave participants a context (fixing
\textsc{pc\textsubscript{prag}} and $g$) and asked which marker they
would use, the comprehension experiment gives participants a context
(fixing \textsc{pc\textsubscript{prag}}) together with a marker already
embedded in the sentence and asks what goal strength the speaker likely
had and how likely the listener is to adopt $q$.

The specific marker presented in each vignette is not chosen arbitrarily
from the three candidate markers. Our preregistered plan was to assign,
for each item$\times$condition cell, the \emph{Bayesian best-choice marker}:
the marker with the highest posterior probability under the cost-augmented
noisy-threshold RSA model fitted to the production data, given the context's
\textsc{pc\textsubscript{prag}} level and a specified $g$ level. This was
intended to mirror the production design closely and to ensure that the
comprehension stimuli reflect model-based marker use rather than arbitrary
marker$\times$item pairings.

Applying this strict modal-assignment rule to the fitted production model,
however, yielded a highly imbalanced stimulus set: most critical cells were
assigned the strongest marker, with very little variation across marker
categories. In addition, many of these modal decisions were supported only by
very small probability margins over the next-best alternative. Because such a
stimulus set would make the comprehension experiment substantially less
informative about marker-based listener inferences, we introduced a documented,
exploratory balancing step for stimulus construction. Starting from the same
posterior predictive probabilities, we selected a near-balanced assignment of
the three markers across the 32 critical cells while maximizing total posterior
support under the fitted production model. The strict modal assignment is
retained as the preregistered baseline, but the balanced assignment is used for
the actual comprehension materials.

\textsc{pc\textsubscript{prag}} is retained as an explicit context
manipulation (low vs.\ high), so that listeners can be expected to
have a well-defined prior on pragmatic controversy when interpreting
the marker. \textsc{pc\textsubscript{prop}} is not manipulated within
the experiment; it enters as a continuous topic-level covariate from
the norming study.

After reading each vignette, participants provide two slider ratings
operationalising the key listener-side quantities:
\begin{enumerate}
  \item \textbf{Inferred speaker goal strength} ($\widehat{g}$):
        how strongly does the speaker seem to want the listener to
        accept and act on $q$?
  \item \textbf{Adoption likelihood} ($\widehat{A}$):
        how likely is the listener to follow the recommended action $q$?
\end{enumerate}

\noindent Perceived controversy is \emph{not} elicited as a rating:
because \textsc{pc\textsubscript{prag}} is already manipulated as a
context factor, asking participants to rate controversy would be
redundant with respect to the manipulation check and circular with
respect to the research question. The theoretically central inference
to test is whether listeners can decode the speaker's \emph{goal} from
marker use, and whether that inference translates into adoption
behaviour.

The two slider measures, together with the \textsc{pc\textsubscript{prag}}
context and the \textsc{pc\textsubscript{prop}} covariate, provide the
data needed to fit the model's listener parameters
$(\eta_0, \eta_g, \eta_{pc}, \eta_{\mathrm{int}})$.

\subsection{Participants}

Participants will be recruited via Prolific and will be required to be
native speakers of German (language-at-home filter).
Data collection is planned after completion of the production experiment
(model fitting required to determine the best-choice markers).
A target of $N = 80$ participants is planned ($n = 20$ per list across
4 counterbalancing lists), each completing 16 trials (8 critical + 8
filler), matching the structure of the production experiment.

\subsection{Materials}

\subsubsection{Markers}

Each critical vignette is presented with one specific consensus marker
already embedded in the Mittelfeld of the sentence frame (see
example~(\ref{ex:frame_comp})):

\begin{exe}
  \ex \label{ex:frame_comp}
  Digitale Kompetenzen sind \textbf{[MARKER]}
  in der heutigen Arbeitswelt unverzichtbar.\\
  `Digital skills are \textbf{[MARKER]} indispensable in today's
  working world.'
\end{exe}

Unlike the production experiment, participants do \emph{not} choose the
marker; they read the sentence with the marker already in place and respond
to the two rating probes.

Critically, the marker assigned to each condition is not drawn arbitrarily
from the three candidates. Under the preregistered plan, for each of the four
production conditions (pc\textsubscript{prag} $\times$ $g$), the marker used in
comprehension was to be the \textbf{Bayesian best-choice marker}: the marker
$u^*$ with the highest posterior probability under the fitted
cost-augmented noisy-threshold RSA production model,
\[
u^* = \arg\max_{u} \; P(u \mid pc_{\mathrm{prag}}, g,
                           pc_{\mathrm{prop}\_c}).
\]
Because $pc_{\mathrm{prop}}$ varies continuously across items,
$u^*$ may differ per item even within the same
(pc\textsubscript{prag}, $g$) cell; the final stimulus list is therefore
constructed item-by-item once the production model has been fitted.
In practice, the fitted production model yielded a nearly degenerate modal map,
with most cells assigned to \textit{bekanntlich}. We therefore retained the
model-based logic of the procedure but replaced strict cellwise modal selection
with an uncertainty-aware balancing step: among the marker assignments with
high posterior support, we selected a near-balanced configuration across the 32
critical cells. This preserves the production--comprehension link while also
ensuring that the comprehension experiment contains sufficient marker variation
to test listener-side inferences.

All three markers remain candidates; the mapping from conditions to
markers is an empirical output of the production analysis, not a design
pre-commitment. In the final comprehension materials, we use the
\textbf{balanced model-based assignment}, rather than the strict modal
assignment, in order to maintain informative variation across marker types.

\subsubsection{Topic Items}

The same eight topic domains used in the norming study and the production
experiment are used here (Table~\ref{tab:items_comp}).
The 4-condition $\times$ 8-item Latin square (Section~\ref{subsub:counterbalancing_comp})
requires no additional items.

\begin{table}[ht]
\centering
\caption{Topic items: sentence frame for $p$ (with marker position [M])
         and consequent $q$.}
\label{tab:items_comp}
\small
\begin{tabular}{clll}
\toprule
\textbf{\#} & \textbf{Topic} & \textbf{Sentence frame} & \textbf{Consequent $q$} \\
\midrule
1 & Klimawandel &
  \textit{Der Klimawandel wird [M]} &
  auf Flugreisen verzichten \\
  & & \textit{durch menschliche Aktivit\"aten verursacht.} & \\[3pt]
2 & Ern\"ahrung &
  \textit{Eine \"uberwiegend pflanzliche Ern\"ahrung} &
  mehr pflanzliche Lebens- \\
  & & \textit{f\"ordert [M] die Gesundheit.} &
  mittel einbauen \\[3pt]
3 & Stadtverkehr &
  \textit{Ein ausgebautes Nahverkehrsnetz} &
  h\"aufiger auf \\
  & & \textit{entlastet [M] den Stadtverkehr} &
  \"offentliche Verkehrs- \\
  & & \textit{deutlich.} &
  mittel umsteigen \\[3pt]
4 & Digitalisierung &
  \textit{Digitale Kompetenzen sind [M] in} &
  regelm\"a\ss{}ig Weiter- \\
  & & \textit{der heutigen Arbeitswelt unverzichtbar.} &
  bildungsangebote nutzen \\[3pt]
5 & Schlaf &
  \textit{Ausreichend Schlaf ist [M] entscheidend} &
  auf feste Schlafzeiten \\
  & & \textit{f\"ur die kognitive Leistungsf\"ahigkeit.} &
  achten \\[3pt]
6 & Plastik &
  \textit{Einwegplastik schadet [M] den} &
  konsequent auf Einweg- \\
  & & \textit{Meeres\"okosystemen erheblich.} &
  plastik verzichten \\[3pt]
7 & Sport &
  \textit{Regelm\"a\ss{}ige k\"orperliche Bewegung} &
  regelm\"a\ss{}ig Sport in \\
  & & \textit{senkt [M] das Risiko f\"ur} &
  den Alltag integrieren \\
  & & \textit{Herz-Kreislauf-Erkrankungen deutlich.} & \\[3pt]
8 & Lokalkauf &
  \textit{Das Kaufen bei lokalen H\"andlern} &
  \"ofter bei lokalen \\
  & & \textit{st\"arkt [M] die regionale Wirtschaft} &
  Gesch\"aften einkaufen \\
  & & \textit{nachhaltig.} & \\
\bottomrule
\end{tabular}
\end{table}

\noindent [M] denotes the position of the pre-filled best-choice marker.

\subsubsection{Vignettes}

\paragraph{Structure.}
Each trial presents a vignette consisting of three parts:
(a) a context paragraph establishing the communicative situation and
fixing the \textsc{pc\textsubscript{prag}} level;
(b) the target sentence with the best-choice marker already embedded
in the Mittelfeld;
(c) the consequent action clause.
The two rating probes follow immediately below the vignette.

\paragraph{Propositional perceived controversy (\textsc{pc\textsubscript{prop}}).}
As in the production experiment, \textsc{pc\textsubscript{prop}} is not
manipulated within the comprehension experiment.
Topic-level propositional acceptance ratings from the norming study
(0--100 scale; e.g., \textit{Digitalisierung}: 85/100,
\textit{Lokalkauf}: 58/100) serve as a continuous, mean-centred
covariate \textsc{pc\textsubscript{prop}}$_c$ in analysis.
Context paragraphs present $p$ without editorial evaluation; variation in
listeners' individual assessments of propositional controversy is
captured by the norming-derived covariate.

\paragraph{Pragmatic perceived controversy (\textsc{pc\textsubscript{prag}}).}
\textsc{pc\textsubscript{prag}} is manipulated between item instances
using the same two context versions as in the production experiment:
the speaker's social circle either \emph{largely agrees} (low) or
\emph{holds divergent views} (high) about $p$.
For example, for the \textit{Digitalisierung} item:

\begin{itemize}
  \item $\textsc{pc}_\text{prag} = \text{low}$: \textit{Deine Kollegin spricht
        mit dir \"uber Berufsausbildung und Weiterbildung.
        In eurem Team teilen die meisten die Ansicht, dass digitale Kompetenzen
        in der heutigen Arbeitswelt unverzichtbar sind.}
  \item $\textsc{pc}_\text{prag} = \text{high}$: \textit{Deine Kollegin spricht
        mit dir \"uber Berufsausbildung und Weiterbildung.
        In eurem Team gehen die Meinungen dazu stark auseinander.}
\end{itemize}

The context is framed from the listener's perspective as a direct addressee,
mirroring the production study where participants adopt the speaker's
perspective.

\paragraph{Marker manipulation.}
The dependent measures of the production experiment (which marker is chosen)
become the independent variable here.
One of the three consensus markers is inserted into the Mittelfeld of the
target sentence.
A complete sample item for the \textit{Digitalisierung} topic
($\textsc{pc}_\text{prag} = \text{high}$, marker = \textit{bekanntlich};
norming $\textsc{pc}_\text{prop}$ rating: 85/100) is given in
(\ref{ex:comp_sample}):

\begin{exe}
  \ex \label{ex:comp_sample}
  \textit{Deine Kollegin spricht mit dir \"uber Berufsausbildung und
  Weiterbildung.
  In eurem Team gehen die Meinungen dazu stark auseinander.}

  \textit{Sie sagt zu dir:} Digitale Kompetenzen sind \textbf{bekanntlich}
  in der heutigen Arbeitswelt unverzichtbar.
  Deshalb solltest du regelm\"a\ss{}ig digitale Weiterbildungsangebote nutzen.

  \noindent\textit{[Rating probes follow.]}
\end{exe}

\subsection{Design}

\subsubsection{Factorial structure}

The comprehension experiment mirrors the production experiment's
$\mathbf{2}$~\textbf{(pc\textsubscript{prag}: low vs.\ high)}
$\mathbf{\times}$~$\mathbf{2}$~\textbf{($g_{\mathrm{implied}}$: low
vs.\ high) design}, yielding \textbf{4 conditions}.
The $g_{\mathrm{implied}}$ dimension is not manipulated via an explicit
goal instruction (as in production) but is instead conveyed entirely
through the choice of marker: for each condition, the marker presented
is the best-choice marker for the corresponding
(\textsc{pc\textsubscript{prag}},~$g$) cell from the production model.
\textsc{pc\textsubscript{prop}} enters as a continuous norming covariate.
With 8 topic items and a Latin square rotation across 4 lists, each
participant sees exactly \textbf{8 critical trials}---one item per
condition, each condition appearing twice---plus 8 filler trials,
for \textbf{16 trials} in total (Table~\ref{tab:design_comp}).

\begin{table}[ht]
\centering
\caption{%
  The $2\times2$ design (4 conditions), mirroring the production experiment.
  The marker presented in each condition is the Bayesian best-choice
  marker from the fitted production model; specific markers are
  determined after production data collection.}
\label{tab:design_comp}
\small
\begin{tabular}{ccll}
\toprule
\textbf{pc\textsubscript{prag}} &
\textbf{$g_{\mathrm{implied}}$} &
\textbf{Marker (TBD from production)} &
\textbf{Condition label} \\
\midrule
low  & low  & best-choice for $(\text{low}, g_{\text{low}})$  &
  (0) community agrees, weak signal \\
low  & high & best-choice for $(\text{low}, g_{\text{high}})$ &
  (1) community agrees, strong signal \\
high & low  & best-choice for $(\text{high}, g_{\text{low}})$ &
  (2) community divided, weak signal \\
high & high & best-choice for $(\text{high}, g_{\text{high}})$ &
  (3) community divided, strong signal \\
\bottomrule
\end{tabular}
\end{table}

\subsubsection{Counterbalancing}
\label{subsub:counterbalancing_comp}

The experiment uses the \textbf{same $4 \times 8$ Latin square} as the
production experiment: each of the 4 lists assigns item $i$
($i = 0,\ldots,7$) to condition $(i + l)\bmod 4$, where
$l \in \{0,1,2,3\}$ is the list number.
This guarantees that:
\begin{itemize}
  \item each participant sees every item exactly once;
  \item each participant sees each of the 4 conditions exactly twice
        (two items per condition), ensuring balanced coverage of the
        $2\times2$ design;
  \item across all 4 lists, each item appears equally often in each
        condition.
\end{itemize}

List assignment follows the same Prolific-and-URL strategy as the production
experiment: each Prolific study embeds \texttt{?list=N} (0--3) in the
experiment URL, capping recruitment at $n = 20$ participants per list
for a total of $N = 80$.

In addition to the 8 critical trials, each participant completed 8
\textbf{filler trials} using the same filler items as the production
experiment (train travel, reading light, cooking, coffee, vacation,
hydration, fresh air, noise), each presented with a randomly assigned
marker.
Fillers are excluded from analysis; they serve to dilute the critical
items and reduce demand characteristics.

\subsubsection{Dependent variables}
\label{subsub:dvs_comp}

Two slider ratings are collected after each vignette (0--100 continuous
scale with no visible numbers, labelled at the endpoints as shown):

\begin{enumerate}
  \item \textbf{Inferred speaker goal strength} (operationalises $\widehat{g}$):
  \begin{quote}
    \textit{Wie stark hast du den Eindruck, dass die Sprecherin m\"ochte,
    dass du ihre Empfehlung befolgst?}\\
    (`How strongly do you get the impression that the speaker wants you to
    follow their recommendation?')\\
    Endpoints: 0 = \textit{\"uberhaupt nicht}
    (`not at all') ---
    100 = \textit{sehr stark} (`very strongly').
  \end{quote}

  \item \textbf{Adoption likelihood of $q$} (operationalises $\widehat{A}$):
  \begin{quote}
    \textit{Wie wahrscheinlich ist es, dass du die empfohlene Handlung
    dann auch tats\"achlich umsetzt?}\\
    (`How likely is it that you will actually carry out the recommended
    action?')\\
    Endpoints: 0 = \textit{sehr unwahrscheinlich}
    (`very unlikely') ---
    100 = \textit{sehr wahrscheinlich} (`very likely').
  \end{quote}
\end{enumerate}

\noindent Perceived controversy is deliberately \emph{not} included as a
rating probe: \textsc{pc\textsubscript{prag}} is experimentally
manipulated via the context paragraph, so a controversy rating would
be redundant and risks obscuring the marker's contribution to inferred
goal strength.

Recorded per trial: \texttt{trial\_id}, \texttt{topic},
\texttt{is\_filler}, \texttt{condition\_index}, \texttt{marker},
\texttt{pc\_prag}, \texttt{g\_implied}, \texttt{rating\_g},
\texttt{rating\_adopt}, \texttt{rt} (ms); and once per participant:
\texttt{list\_num}.
The continuous \texttt{pc\_prop\_rating} covariate (topic-level norming
mean) is joined to the trial data during analysis.

\subsection{Procedure}

Participants complete the experiment online via the magpie platform
\citep{kochetkova2024magpie} in a browser.
After reading instructions in German, participants are assigned to one of
4 counterbalancing lists (via URL \texttt{?list=N} parameter assigned by
Prolific; no random fallback in deployment).
They then proceed through 16 trials---8 critical and 8 fillers---in a
fully randomised order, followed by a post-test questionnaire eliciting
native language, age, gender, education level, and optional comments.

Each trial displays:
(a) a context paragraph establishing the communicative situation and the
\textsc{pc\textsubscript{prag}} level;
(b) the target sentence with the best-choice marker already in place in
the Mittelfeld;
(c) the consequent action clause;
(d) the two rating sliders in fixed order ($\widehat{g}$, $\widehat{A}$),
each labelled in German.
The ``Weiter'' (Next) button is disabled until both sliders have been
moved at least once, ensuring a complete response on every trial.

\subsection{Planned Analysis}

\subsubsection{Empirical plots}

For each dependent variable ($\widehat{g}$, $\widehat{A}$):

\begin{itemize}
  \item \textbf{Condition means}: a dot-and-whisker plot shows mean
        ratings per condition (the $2\times2$ of pc\textsubscript{prag}
        $\times$ $g_{\mathrm{implied}}$) with 95\% bootstrapped CIs.
        The expected pattern: $\widehat{g}$ and $\widehat{A}$ should
        be higher in the high-$g_{\mathrm{implied}}$ conditions
        (where the model assigns a stronger marker), and lower in the
        high-pc\textsubscript{prag} conditions (credibility discounting).
  \item \textbf{Marker identity}: because the best-choice marker may
        differ across items within the same condition, a secondary plot
        shows ratings grouped by actual marker identity, allowing
        inspection of whether marker type (rather than just condition)
        drives the effects.
  \item \textbf{pc\textsubscript{prop} covariate}: scatter plots display
        each DV as a function of the continuous norming rating per
        condition, testing whether higher propositional acceptance
        amplifies the $\widehat{g}$ and $\widehat{A}$ signals for
        strong markers.
\end{itemize}

\subsubsection{Bayesian mixed-effects regression}

Two Bayesian mixed-effects regression models are fitted via
\texttt{brms} \citep{burkner2017}, one per DV ($\widehat{g}$ and
$\widehat{A}$):

\begin{align}
y &\sim \mathrm{Beta}(\mu, \phi)
\quad \text{(ratings rescaled to }(0,1)\text{)} \\
\mathrm{logit}(\mu) &= \alpha
+ \beta_{g} \cdot \mathrm{g\_implied\_c}
+ \beta_{\mathrm{prag}} \cdot \mathrm{pc\_prag\_c}
+ \beta_{\mathrm{prop}} \cdot \mathrm{pc\_prop\_c}
+ \beta_{\mathrm{int}} \cdot \mathrm{g\_implied\_c}
  \times \mathrm{pc\_prag\_c}
+ (1 \mid \mathrm{submission\_id})
+ (1 \mid \mathrm{topic}),
\end{align}

where \texttt{g\_implied\_c} is the $g$ level from the corresponding
production condition ($-0.5$ = low, $+0.5$ = high), effectively
coding the marker's position on the production utility scale.
A secondary specification replaces \texttt{g\_implied\_c} with
actualy marker identity as an unordered factor, to examine
non-monotonic patterns and verify that marker rather than
condition label drives effects.

Model comparison via approximate leave-one-out cross-validation (LOO)
selects among additive, interaction, and marker-as-categorical
formulations.

\subsubsection{RSA listener parameter estimation}
\label{subsub:rsa_fit}

The two empirical rating means (per condition cell: 4 conditions) constrain
the listener parameters of the noisy-threshold RSA model.
Specifically, the model predicts:

\begin{align}
\mathbb{E}[\widehat{g} \mid u, \text{context}]
  &= \int g \cdot P(g \mid u, pc_{\mathrm{prag}})\, dg,
\label{eq:pred_g} \\
P(\widehat{A} \mid u, \text{context})
  &= \iiint \sigma(\eta_0 + \eta_g g - \eta_{pc}\, pc
     - \eta_{\mathrm{int}}\, g \cdot pc)\;
     P(pc_{\mathrm{prop}}, pc_{\mathrm{prag}}, g \mid u)\,
     d{pc_{\mathrm{prop}}}\,d{pc_{\mathrm{prag}}}\,dg,
\label{eq:adoption}
\end{align}

\noindent where the posterior $P(pc_{\mathrm{prop}}, pc_{\mathrm{prag}}, g
\mid u)$ is conditioned on the context's \textsc{pc\textsubscript{prag}}
level as a known prior (since it is explicitly given to the listener).

The threshold parameters $(\mu_1,\mu_2, \sigma_1,\sigma_2,
\lambda)$ together with the marker-specific costs $(c_u)$ are estimated
jointly from the production experiment data via the Stan model.
Given those threshold estimates, the listener parameters
$(\eta_0, \eta_g, \eta_{pc}, \eta_{\mathrm{int}})$ are then estimated by
minimising the weighted squared deviation between the model-predicted
means (Eqs.~\ref{eq:pred_g}--\ref{eq:adoption}) and the empirical rating
means from the comprehension experiment, or via a full joint Bayesian
fit over both experiments simultaneously.

\subsection{Model Predictions}

The noisy-threshold RSA model makes the following predictions for
$\widehat{g}$ and $\widehat{A}$, which serve as pre-registered hypotheses.
In all cases, ``strong-marker condition'' refers to conditions 1 and 3
(high $g_{\mathrm{implied}}$, i.e.\ the marker selected by the production
model under a high-$g$ context), and ``weak-marker condition'' to conditions
0 and 2.

\begin{enumerate}
  \item \textbf{Marker encodes inferred goal $\widehat{g}$}: $\widehat{g}$
        ratings should be higher in strong-marker conditions than in
        weak-marker conditions, across both pc\textsubscript{prag} levels.
        This is the core inversion prediction: if the production model
        correctly captures speaker behaviour, listeners should be able to
        decode $g$ from the marker.

  \item \textbf{Marker encodes adoption $\widehat{A}$}: $\widehat{A}$
        ratings should likewise be higher in strong-marker conditions,
        reflecting that higher inferred $g$ translates into greater
        adoption pressure downstream.

  \item \textbf{pc\textsubscript{prag} modulates $\widehat{A}$ but not
        $\widehat{g}$}: because $g$ is orthogonal to pc\textsubscript{prag}
        in the model's utility function, the effect of the marker on
        $\widehat{g}$ should be largely the same in low- and high-
        pc\textsubscript{prag} contexts.
        By contrast, $\widehat{A}$ should be reduced in high-
        pc\textsubscript{prag} contexts via the interaction term
        $\eta_{\mathrm{int}}\,g \cdot pc$: even a strong-marker signal
        for high $g$ is partly cancelled when controversy is high.

  \item \textbf{Credibility discounting}: for strong-marker conditions
        presented in a high-pc\textsubscript{prag} context (condition~3),
        $\widehat{A}$ should be \emph{lower} than for the same marker in
        a low-pc\textsubscript{prag} context (condition~1).
        This is the model's signature prediction distinguishing Bayesian
        inversion from a simpler marker-strength heuristic: a sceptical
        listener who knows controversy is high will discount the adoption
        implication of even the strongest marker.
        The $g_{\mathrm{implied}} \times \textsc{pc\textsubscript{prag}}$
        interaction on $\widehat{A}$ provides the empirical test.
\end{enumerate}

\end{document}
